% Lesson 12 — Writing: Compositions about taking part in recreational activities — Anglais 1ère A — Cours signé OnBuch+
% Programme officiel MINESEC. Compiler avec : tectonic ENA12-writing-compositions-about-taking-part-in-recreational.tex
\newcommand{\DOCMATIERE}{Anglais}
\newcommand{\DOCNIVEAU}{1ère A}
\newcommand{\DOCMODULE}{Chapter 1 : Family and Social Life}
\newcommand{\DOCLECON}{ENA12}
\newcommand{\DOCTITRE}{Lesson 12 — Writing: Compositions about taking part in recreational activities}
% OnBuch+ — préambule « cours signés » ENRICHI (compilé avec tectonic / XeLaTeX)
% Système visuel « Pop » d'OnBuch+ : fond crème (jamais blanc), bordures encre
% épaisses, ombres « dures » décalées, titres Archivo Black en capitales.
% Version MATHÉMATIQUES : graphes (pgfplots/tikz), boîtes pédagogiques
% (définition, propriété, méthode, attention, le savais-tu, exercice résolu,
% exercice, TP, bilan), page de garde et signature OnBuch+ ancrée en bas.
%
% Macros attendues dans main.tex AVANT \input{preamble} :
%   \DOCMATIERE  \DOCNIVEAU  \DOCMODULE  \DOCLECON (numéro)  \DOCTITRE
\documentclass[11pt]{article}

\usepackage{fontspec}
\usepackage[english,french]{babel} % français = langue principale ; l'anglais via \eng{..} / textebox
\usepackage[a4paper,margin=2cm,top=3.4cm,bottom=2.8cm,headheight=1.4cm,footskip=1.3cm]{geometry}
\usepackage{amsmath,amssymb,amsfonts}
\usepackage{mathtools} % \xmapsto, \coloneqq... souvent utilisés par les modèles
\usepackage{cancel} % simplifications barrées (bilans de forces)
% \abs{z} (module, valeur absolue) et \norm{u} (norme), utilisés par les modèles
\providecommand{\abs}{}\let\abs\relax\DeclarePairedDelimiter{\abs}{\lvert}{\rvert}
\providecommand{\norm}{}\let\norm\relax\DeclarePairedDelimiter{\norm}{\lVert}{\rVert}
\usepackage[table]{xcolor} % [table] : fournit aussi \rowcolors (alternance de lignes)
\usepackage{pagecolor}
\usepackage{tikz}
\usetikzlibrary{arrows.meta,positioning,calc,shapes.geometric,shapes.misc,shapes.arrows,decorations.pathmorphing,decorations.pathreplacing,decorations.markings,patterns,shadows,mindmap,trees,fit,backgrounds,matrix,intersections,angles,quotes}
\usepackage{pgfplots}
\usepackage[version=4]{mhchem} % équations nucléaires \ce{^{235}_{92}U}
\usepackage[european,siunitx]{circuitikz} % schémas électriques
\pgfplotsset{compat=1.18}
\usepgfplotslibrary{fillbetween,groupplots} % aires entre courbes (calcul intégral)
\usepackage{enumitem}
\usepackage{fancyhdr}
% [most] charge toutes les bibliothèques tcolorbox (listings, minted, raster,
% documentation...) et gonfle inutilement la mémoire TeX sur un document long
% (~300 boîtes) : on ne charge que ce qui est réellement utilisé.
\usepackage[skins,breakable]{tcolorbox}
\usepackage{graphicx}
\usepackage{colortbl}
\usepackage{booktabs}
\usepackage{tabularx}
\usepackage{diagbox} % case d'en-tête diagonale des tableaux à double entrée
% Colonnes extensibles centrée / gauche / droite (C, L, R), souvent utilisées par les modèles
\newcolumntype{C}{>{\centering\arraybackslash}X}
\newcolumntype{L}{>{\raggedright\arraybackslash}X}
\newcolumntype{R}{>{\raggedleft\arraybackslash}X}
\usepackage{array}
\usepackage{multicol}
\usepackage{multirow}
\usepackage{siunitx}
% parse-numbers=false : accepte aussi \SI{2,5 \times 10^{-3}}{...}
\sisetup{locale=FR,output-decimal-marker={,},group-minimum-digits=4,parse-numbers=false}
\DeclareSIUnit\ohm{\ensuremath{\symup{\Omega}}} % U+2126 absent de la police
\DeclareSIUnit\division{div} % réglages d'oscilloscope (ms/div, V/div)
\DeclareSIUnit\tr{tr} % tours (vitesse de rotation, ex. \SI{1500}{\tr\per\minute})
\DeclareSIUnit\molar{M} % concentration molaire (ex. \SI{3}{\molar}, \SI{1}{\micro\molar})
\DeclareSIUnit\an{an} % année (le modèle confond parfois avec \year, un compteur TeX) — ex. \SI{5}{\kilo\meter\per\an}
\DeclareSIUnit\FCFA{FCFA} % franc CFA (ex. \SI{500}{\FCFA\per\kilo\gram})
\DeclareSIUnit\degreeBrix{\textdegree Brix} % teneur en sucre d'un sirop/jus (réfractométrie)
\DeclareSIUnit\month{mois} % le modèle utilise parfois \month, un compteur TeX, comme unité de durée
\usepackage{qrcode}
% \degree : commande fréquemment utilisée par les modèles pour un angle ou une
% température en dehors de \SI{}{} (non fournie par les paquets chargés).
\providecommand{\degree}{^{\circ}}
% \qed (fin de démonstration), utilisé par les modèles sans amsthm
\providecommand{\qed}{\hfill$\square$}
% \barre : double barre des valeurs interdites dans un tableau de variations (tkz-tab)
\providecommand{\barre}{\ensuremath{\|}}
% environnement proof (amsthm non chargé) utilisé par les modèles
\ifdefined\proof\else\newenvironment{proof}[1][Démonstration]{\par\noindent\textit{#1.}\ }{\qed\par}\fi
\usepackage{lastpage}
\usepackage{hyperref}
\hypersetup{hidelinks}

% ============================================================
% Palette « Pop » OnBuch+ — mêmes hex que la classe P (plus_ui.dart)
% ============================================================
\definecolor{poporange}{HTML}{F59321}
\definecolor{popdark}{HTML}{E07A0C}
\definecolor{popcream}{HTML}{FFF6E8}
\definecolor{popink}{HTML}{1C1714}
\definecolor{poppurple}{HTML}{7A5AE0}
\definecolor{popgreen}{HTML}{2E9E6A}
\definecolor{poppink}{HTML}{E0457B}
\definecolor{popblue}{HTML}{2F7DE1}
\definecolor{popgold}{HTML}{C98A12}
\definecolor{popmuted}{HTML}{8A7F76}
\definecolor{popline}{HTML}{EADFD2}
\definecolor{popgray}{HTML}{8A7F76}
\colorlet{popgrey}{popgray}
% Alias tolérants (le modèle nomme parfois les couleurs différemment)
\colorlet{popwhite}{white}
\colorlet{popblack}{popink}
\colorlet{popred}{poppink}
\colorlet{popyellow}{popgold}
% Teintes claires (fonds de tableaux/encadrés) — le modèle nomme parfois
% les couleurs avec un suffixe L (ex. popblueL) sans les avoir définies.
\colorlet{poporangeL}{poporange!20}
\colorlet{popdarkL}{popdark!20}
\colorlet{poppurpleL}{poppurple!20}
\colorlet{popgreenL}{popgreen!20}
\colorlet{poppinkL}{poppink!20}
\colorlet{popblueL}{popblue!20}
\colorlet{popgoldL}{popgold!20}
\colorlet{popredL}{popred!20}
\colorlet{popgrayL}{popgray!20}
% Teintes claires pour fonds de tableaux / zones
\colorlet{popblueL}{popblue!10!white}
\colorlet{poporangeL}{poporange!14!white}
\colorlet{popgreenL}{popgreen!12!white}
\colorlet{poppurpleL}{poppurple!12!white}
\colorlet{poppinkL}{poppink!12!white}
\colorlet{popgoldL}{popgold!14!white}

\pagecolor{popcream}

% ============================================================
% Typographie
% ============================================================
\defaultfontfeatures{Ligatures=TeX}
\setmainfont{PlusJakartaSans-Medium.ttf}[Path=../../../fonts/,BoldFont=PlusJakartaSans-Bold.ttf]
% Maths en sans-serif (Fira Math) avec chiffres et lettres droites de Plus
% Jakarta, pour que les formules se fondent dans le texte.
\usepackage[math-style=ISO,bold-style=ISO]{unicode-math}
\setmathfont{FiraMath-Regular.otf}[Path=../../../fonts/]
\setmathfont{PlusJakartaSans-Medium.ttf}[Path=../../../fonts/,range={up/num,up/{Latin,latin},\mathup}]
\usepackage{icomma} % virgule décimale sans espace en mode math
% Caractères mathématiques Unicode absents des polices de texte (Plus Jakarta,
% Archivo Black) : rendus par leur équivalent mathématique, en texte comme en
% formule — sinon ils s'affichent en carré vide (ex. « Arithmétique dans ℤ »).
\usepackage{newunicodechar}
\newunicodechar{ℝ}{\ensuremath{\mathbb{R}}}
\newunicodechar{ℤ}{\ensuremath{\mathbb{Z}}}
\newunicodechar{ℂ}{\ensuremath{\mathbb{C}}}
\newunicodechar{ℕ}{\ensuremath{\mathbb{N}}}
\newunicodechar{ℚ}{\ensuremath{\mathbb{Q}}}
\newunicodechar{𝔻}{\ensuremath{\mathbb{D}}}
\newunicodechar{α}{\ensuremath{\alpha}}
\newunicodechar{β}{\ensuremath{\beta}}
\newunicodechar{γ}{\ensuremath{\gamma}}
\newunicodechar{δ}{\ensuremath{\delta}}
\newunicodechar{ε}{\ensuremath{\varepsilon}}
\newunicodechar{θ}{\ensuremath{\theta}}
\newunicodechar{λ}{\ensuremath{\lambda}}
\newunicodechar{μ}{\ensuremath{\mu}}
\newunicodechar{σ}{\ensuremath{\sigma}}
\newunicodechar{τ}{\ensuremath{\tau}}
\newunicodechar{φ}{\ensuremath{\varphi}}
\newunicodechar{ψ}{\ensuremath{\psi}}
\newunicodechar{ω}{\ensuremath{\omega}}
\newunicodechar{Σ}{\ensuremath{\Sigma}}
% majuscules produites par \MakeUppercase dans les titres (α-aminés -> Α)
\newunicodechar{Α}{\ensuremath{\alpha}}
\newunicodechar{Β}{\ensuremath{\beta}}
\newunicodechar{Γ}{\ensuremath{\Gamma}}
\newunicodechar{Δ}{\ensuremath{\Delta}}
\newunicodechar{Ε}{\ensuremath{\varepsilon}}
\newunicodechar{Θ}{\ensuremath{\Theta}}
\newunicodechar{Λ}{\ensuremath{\Lambda}}
\newunicodechar{Μ}{\ensuremath{\mu}}
\newunicodechar{Τ}{\ensuremath{\tau}}
\newunicodechar{Φ}{\ensuremath{\Phi}}
\newunicodechar{Ψ}{\ensuremath{\Psi}}
\newunicodechar{Ω}{\ensuremath{\Omega}}
\newunicodechar{↦}{\ensuremath{\mapsto}}
\newunicodechar{∈}{\ensuremath{\in}}
\newunicodechar{∉}{\ensuremath{\notin}}
\newunicodechar{∧}{\ensuremath{\wedge}}
\newunicodechar{⇔}{\ensuremath{\Leftrightarrow}}
\newunicodechar{⟺}{\ensuremath{\Longleftrightarrow}}
\newunicodechar{⇒}{\ensuremath{\Rightarrow}}
\newunicodechar{⟹}{\ensuremath{\Longrightarrow}}
\newunicodechar{∘}{\ensuremath{\circ}}
\newunicodechar{∪}{\ensuremath{\cup}}
\newunicodechar{∩}{\ensuremath{\cap}}
\newunicodechar{≡}{\ensuremath{\equiv}}
\newunicodechar{⊂}{\ensuremath{\subset}}
\newunicodechar{⊥}{\ensuremath{\perp}}
\newunicodechar{∃}{\ensuremath{\exists}}
\newunicodechar{∀}{\ensuremath{\forall}}
\newunicodechar{∛}{\ensuremath{\sqrt[3]{\,}}}
\newunicodechar{∜}{\ensuremath{\sqrt[4]{\,}}}
\newunicodechar{⃗}{\ensuremath{{}^{\to}}}
\newunicodechar{ⁿ}{\ifmmode^{n}\else\textsuperscript{\ensuremath{n}}\fi}
\newunicodechar{ˣ}{\ifmmode^{x}\else\textsuperscript{\ensuremath{x}}\fi}
\newunicodechar{ᵇ}{\ifmmode^{b}\else\textsuperscript{\ensuremath{b}}\fi}
\newunicodechar{ʳ}{\ifmmode^{r}\else\textsuperscript{\ensuremath{r}}\fi}
\newunicodechar{ᵘ}{\ifmmode^{u}\else\textsuperscript{\ensuremath{u}}\fi}
\newunicodechar{ʸ}{\ifmmode^{y}\else\textsuperscript{\ensuremath{y}}\fi}
\newunicodechar{ᵃ}{\ifmmode^{a}\else\textsuperscript{\ensuremath{a}}\fi}
\newunicodechar{ᵉ}{\ifmmode^{e}\else\textsuperscript{\ensuremath{e}}\fi}
\newunicodechar{ᵏ}{\ifmmode^{k}\else\textsuperscript{\ensuremath{k}}\fi}
\newunicodechar{ᵖ}{\ifmmode^{p}\else\textsuperscript{\ensuremath{p}}\fi}
\newunicodechar{ᴺ}{\ifmmode^{N}\else\textsuperscript{\ensuremath{N}}\fi}
\newunicodechar{ᶜ}{\ifmmode^{c}\else\textsuperscript{\ensuremath{c}}\fi}
\newunicodechar{ʰ}{\ifmmode^{h}\else\textsuperscript{\ensuremath{h}}\fi}
\newunicodechar{ᵠ}{\ifmmode^{q}\else\textsuperscript{\ensuremath{q}}\fi}
\newunicodechar{ₙ}{\ifmmode_{n}\else\textsubscript{\ensuremath{n}}\fi}
\newunicodechar{ₐ}{\ifmmode_{a}\else\textsubscript{\ensuremath{a}}\fi}
\newunicodechar{ᵢ}{\ifmmode_{i}\else\textsubscript{\ensuremath{i}}\fi}
\newunicodechar{ᵣ}{\ifmmode_{r}\else\textsubscript{\ensuremath{r}}\fi}
\newunicodechar{ₖ}{\ifmmode_{k}\else\textsubscript{\ensuremath{k}}\fi}
\newunicodechar{ᵦ}{\ifmmode_{\beta}\else\textsubscript{\ensuremath{\beta}}\fi}
\newunicodechar{ₓ}{\ifmmode_{x}\else\textsubscript{\ensuremath{x}}\fi}
\newfontfamily\popdisplay{ArchivoBlack-Regular.ttf}[Path=../../../fonts/]
\newfontfamily\poplogo{SpaceGrotesk-Bold.ttf}[Path=../../../fonts/]
\newfontfamily\popsemi{PlusJakartaSans-SemiBold.ttf}[Path=../../../fonts/]
\newfontfamily\popxbold{PlusJakartaSans-ExtraBold.ttf}[Path=../../../fonts/]
\newfontfamily\popxboldls{PlusJakartaSans-ExtraBold.ttf}[Path=../../../fonts/,LetterSpace=3.0]

\setlist[enumerate]{leftmargin=1.7em,itemsep=5pt,topsep=4pt}
\setlist[itemize]{leftmargin=1.7em,itemsep=4pt,topsep=4pt,label={\color{poporange}\textbullet}}
\setlength{\parindent}{0pt}
\setlength{\parskip}{5pt}
\renewcommand{\arraystretch}{1.25}

% pgfplots : style Pop par défaut
\pgfplotsset{
  every axis/.append style={axis line style={popink,line width=1pt},tick style={popink},
    grid style={popline,line width=0.5pt},label style={font=\small},tick label style={font=\footnotesize}},
  popcurve/.style={poporange,line width=1.8pt,no markers,smooth},
}
% Même style disponible dans un \draw TikZ simple (hors pgfplots)
\tikzset{popcurve/.style={poporange,line width=1.8pt}}
\tikzset{popgrid/.style={popline,very thin}}
\colorlet{popcurve}{poporange} % le nom du style est parfois utilisé comme couleur

% ============================================================
% Pill / squiggle
% ============================================================
\newcommand{\poppill}[3]{%
  \tikz[baseline=-0.55ex]\node[draw=popink,line width=1.2pt,rounded corners=2.6pt,%
    fill=#2,inner xsep=6.5pt,inner ysep=3pt]{%
    \popxboldls\fontsize{7}{7}\selectfont\color{#3}\MakeUppercase{#1}};%
}
\newcommand{\popsquiggle}[1]{%
  \tikz[baseline=-0.4ex]{
    \draw[#1,line width=2.6pt,line cap=round]
      plot [smooth] coordinates {(0,0.09) (0.34,0.01) (0.68,0.21) (1.02,0.01) (1.36,0.10)};
  }%
}
% Mot-clé surligné (marqueur orange sous le texte)
\newcommand{\cle}[1]{{\popxbold\color{popdark}#1}}

% ============================================================
% En-tête / pied
% ============================================================
\pagestyle{fancy}
\fancyhf{}
\renewcommand{\headrulewidth}{0pt}
\renewcommand{\footrulewidth}{0pt}
\lhead{\raisebox{-0.15ex}{\poplogo\fontsize{13}{13}\selectfont\color{popink}OnBuch\color{poporange}+}}
\rhead{\poppill{\DOCMATIERE\ \textbullet\ \DOCNIVEAU\ \textbullet\ Leçon \DOCLECON}{white}{popink}}
\lfoot{\popxbold\fontsize{7.6}{7.6}\selectfont\color{popmuted}\MakeUppercase{Cours signé OnBuch+}\ \textbullet\ {\popsemi\DOCTITRE}}
\rfoot{\tikz[baseline=-0.6ex]\node[rounded corners=8.5pt,fill=popink,minimum height=17pt,minimum width=17pt,inner xsep=5pt,inner sep=0]{\popxbold\fontsize{8}{8}\selectfont\color{popcream}\thepage\,/\,\pageref*{LastPage}};}

% ============================================================
% Titres
% ============================================================
\newcommand{\chaptitre}[2]{%
  \begin{tcolorbox}[enhanced,colback=poporange,colframe=popink,boxrule=2.6pt,arc=11pt,
      drop shadow=popink,left=15pt,right=15pt,top=12pt,bottom=11pt,before skip=3pt,after skip=18pt]
    \poppill{#2}{popink}{popcream}\par\vspace{9pt}
    {\raggedright\hyphenpenalty=10000\exhyphenpenalty=10000\popdisplay\fontsize{22}{24}\selectfont\color{popcream}\MakeUppercase{#1}\par}
  \end{tcolorbox}
}
% Section numérotée : pastille orange avec le numéro + titre Archivo
\newcounter{popsec}
\newcommand{\coursec}[1]{%
  \stepcounter{popsec}\setcounter{popsub}{0}%
  \par\vspace{16pt}\noindent
  \begin{minipage}{\linewidth}
  \tikz[baseline=-0.7ex]\node[draw=popink,line width=1.6pt,fill=poporange,rounded corners=4pt,minimum size=22pt,inner sep=0]{\popdisplay\fontsize{12}{12}\selectfont\color{popcream}\thepopsec};%
  \hspace{8pt}{\popdisplay\fontsize{15}{17}\selectfont\color{popink}\MakeUppercase{#1}}\par
  \vspace{4pt}\noindent{\color{poporange}\rule{36pt}{3.6pt}}
  \end{minipage}\par\nopagebreak\vspace{8pt}
}
\newcounter{popsub}
\newcommand{\courssub}[1]{%
  \stepcounter{popsub}%
  \par\vspace{9pt}\noindent{\popxbold\fontsize{11.5}{13}\selectfont\color{popink}{\color{poporange}\thepopsec.\thepopsub}\hspace{6pt}#1}\par\nopagebreak\vspace{4pt}
}

% ============================================================
% Boîtes pédagogiques — même recette : fond blanc, bordure épaisse colorée,
% ombre dure encre, badge plein. Titre optionnel [#1].
% ============================================================
\tcbset{popbase/.style={breakable,enhanced,colback=white,boxrule=2.3pt,arc=9pt,
  shadow={3pt}{-3pt}{0pt}{popink},left=14pt,right=14pt,top=11pt,bottom=11pt,before skip=11pt,after skip=15pt,
  fontupper=\popsemi\fontsize{10.6}{14.5}\selectfont\color{popink},
  fontlower=\popsemi\fontsize{10.6}{14.5}\selectfont\color{popink}}}
% Coche dessinée (le glyphe ✓ n'existe pas dans les polices du document)
\newcommand{\popcheck}[1]{\tikz[baseline=-0.5ex]\draw[#1,line width=1.6pt,line cap=round,line join=round](-0.13,0.0)--(-0.04,-0.09)--(0.13,0.1);}
\providecommand{\coche}{\popcheck{popgreen}}
\providecommand{\resultat}[1]{\par\smallskip\noindent\fcolorbox{popgreen}{popgreenL}{\parbox{\dimexpr\linewidth-2\fboxsep-2\fboxrule\relax}{\textbf{Résultat.} #1}}\par\smallskip}
\newcommand{\popboxhead}[3]{\poppill{#1}{#2}{white}\ifx\relax#3\relax\else\hspace{6pt}{\popxbold\fontsize{10.6}{12}\selectfont\color{popink}#3}\fi\par\vspace{7pt}}

\newtcolorbox{aretenir}[1][]{popbase,colframe=popgreen,before upper={\popboxhead{À retenir}{popgreen}{#1}}}
% Texte en anglais (document de lecture, dialogue, extrait) : cadre
% d'étude. Un titre optionnel (source, auteur) s'affiche dans l'en-tête.
\newtcolorbox{textebox}[1][]{popbase,colframe=popblue,before upper={\popboxhead{Text}{popblue}{#1}\begin{otherlanguage}{english}},after upper={\end{otherlanguage}}}
% Marques ✓ / ✗ (forme correcte / fautive) souvent utilisées par les modèles
\providecommand{\cross}{\textcolor{poppink}{\ensuremath{\times}}}
\providecommand{\xmark}{\cross}\providecommand{\crossmark}{\cross}\providecommand{\cmark}{\ensuremath{\checkmark}}
% Ligne de réponse / justification à compléter (inventée par les modèles)
\providecommand{\justification}[1]{\par\noindent\textit{Justification~:} \dotfill\par}
% \textipa (package tipa non chargé) : la police ne couvre pas l'API -> simple passage
\providecommand{\textipa}[1]{#1}
% Soulignement (ulem non chargé) : \uline, \ul
\providecommand{\uline}[1]{\underline{#1}}\providecommand{\ul}[1]{\underline{#1}}
% \glossaire{\item ...} inventée par les modèles : liste de glossaire
\providecommand{\glossaire}[1]{\begin{itemize}#1\end{itemize}}
% \alert (beamer) inventée par les modèles : texte mis en évidence
\providecommand{\alert}[1]{\textbf{\textcolor{poppink}{#1}}}
% variantes ASCII d'ordinaux écrites en macro
\providecommand{\iere}{\textsuperscript{re}}\providecommand{\ieme}{\textsuperscript{e}}\providecommand{\ere}{\textsuperscript{re}}\providecommand{\eme}{\textsuperscript{e}}\providecommand{\er}{\textsuperscript{er}}\providecommand{\ie}{\textsuperscript{e}}
% Ordinaux français écrits en macro par les modèles : 1\ière, 3\ième
\newcommand{\ière}{\textsuperscript{re}}\newcommand{\ième}{\textsuperscript{e}}
% \exemple (inventée par les modèles) : étiquette « Exemple : »
\providecommand{\exemple}{\textbf{Exemple~:}\ }
% \square utilisable aussi hors mode math (cases à cocher)
\AtBeginDocument{\let\squareMath\square\renewcommand{\square}{\ensuremath{\squareMath}}}
% Mot ou phrase en anglais à l'intérieur d'un paragraphe français
\newcommand{\eng}[1]{\ifmmode\text{\textit{#1}}\else\begingroup\selectlanguage{english}\itshape #1\endgroup\fi}
\let\esp\eng % alias toléré
% flèches d'intonation inventées par les modèles
\providecommand{\textsearrow}{}\renewcommand{\textsearrow}{\ensuremath{\searrow}}\providecommand{\textnearrow}{}\renewcommand{\textnearrow}{\ensuremath{\nearrow}}
% \fr{..} : mot français (inventé par les modèles)
\providecommand{\fr}[1]{\textit{#1}}
\newtcolorbox{exemplebox}[1][]{popbase,colframe=poporange,before upper={\popboxhead{Exemple}{poporange}{#1}}}
\newtcolorbox{definition}[1][]{popbase,colframe=popblue,before upper={\popboxhead{Définition}{popblue}{#1}}}
\newtcolorbox{propriete}[1][]{popbase,colframe=poppurple,before upper={\popboxhead{Propriété}{poppurple}{#1}}}
\newtcolorbox{methode}[1][]{popbase,colframe=popgold,before upper={\popboxhead{Méthode}{popgold}{#1}}}
\newtcolorbox{attention}[1][]{popbase,colframe=poppink,before upper={\popboxhead{Attention}{poppink}{#1}}}
\newtcolorbox{experience}[1][]{popbase,colframe=popblue,colback=popblueL,before upper={\popboxhead{Expérience}{popblue}{#1}}}
% « Le savais-tu ? » : carte violette pleine (contexte, culture, Cameroun)
\newtcolorbox{savaistu}[1][]{popbase,colback=poppurple,colframe=popink,
  fontupper=\popsemi\fontsize{10.4}{14.2}\selectfont\color{white},
  before upper={\poppill{Le savais-tu\,?}{white}{poppurple}\ifx\relax#1\relax\else\hspace{6pt}{\popxbold\color{white}#1}\fi\par\vspace{7pt}}}
% Exercice résolu : énoncé en haut, \tcblower puis la solution sur fond crème
\newcounter{popexo}\newcounter{popexr}
\newtcolorbox{exoresolu}[1][]{popbase,colframe=poporange,colbacklower=popcream,
  segmentation style={popink,line width=1.2pt,dashed},
  before upper={\stepcounter{popexr}\popboxhead{Exercice résolu \thepopexr}{poporange}{#1}},
  before lower={\poppill{Solution}{popink}{popcream}\par\vspace{6pt}}}
% Exercice (non résolu en place) — la correction est dans la partie Corrigés
\newtcolorbox{exercice}[1][]{popbase,colframe=popink,
  before upper={\stepcounter{popexo}\popboxhead{Exercice \thepopexo}{popink}{#1}}}
% Corrigé
\newtcolorbox{corrige}[1][]{popbase,colframe=popgreen,colback=popgreenL,
  before upper={\popboxhead{Corrigé}{popgreen}{#1}}}
% Objectifs / prérequis / bilan
\newtcolorbox{objectifs}{popbase,colframe=popink,colback=poporangeL,
  before upper={\popboxhead{Objectifs de la leçon}{popink}{}}}
\newtcolorbox{prerequis}{popbase,colframe=popink,colback=popblueL,
  before upper={\popboxhead{Prérequis}{popblue}{}}}
\newtcolorbox{bilan}{popbase,colframe=popink,colback=white,boxrule=2.8pt,
  before upper={\popboxhead{Fiche bilan}{poporange}{L'essentiel à connaître pour l'examen}}}
% Figure encadrée avec légende
\newtcolorbox{popfigure}[1][]{enhanced,breakable=false,colback=white,colframe=popline,boxrule=1.4pt,arc=8pt,
  left=10pt,right=10pt,top=10pt,bottom=8pt,before skip=10pt,after skip=12pt,halign=center,
  lower separated=false,halign lower=center,
  fontlower=\popsemi\fontsize{9}{11.5}\selectfont\color{popmuted},
  #1}
\newcommand{\legende}[1]{\par\vspace{6pt}{\popsemi\fontsize{9}{11.5}\selectfont\color{popmuted}#1}}

% ============================================================
% Signature OnBuch+
% ============================================================
\newcommand{\poplogoinline}[1]{{\poplogo\fontsize{#1}{#1}\selectfont\color{popink}OnBuch\color{poporange}+}}
\newcommand{\signatureonbuch}{%
  \begin{tcolorbox}[enhanced,colback=popink,colframe=popink,boxrule=2.2pt,arc=10pt,
      drop shadow=poporange,left=14pt,right=14pt,top=10pt,bottom=10pt,before skip=0pt,after skip=0pt]
    \tikz[baseline=-0.7ex]\node[circle,fill=popgreen,draw=popcream,line width=1.3pt,minimum size=22pt,inner sep=0]{\popcheck{white}};%
    \hspace{9pt}\begin{minipage}[c]{0.62\linewidth}
      {\popxbold\fontsize{12}{13}\selectfont\color{popcream}Signé\ \textbullet\ {\poplogo OnBuch\color{poporange}+}}\par\vspace{2pt}
      {\raggedright\popsemi\fontsize{8}{10.5}\selectfont\color{popline}\MakeUppercase{Équipe pédagogique OnBuch+}\\\MakeUppercase{Conforme au programme officiel MINESEC}\par}
    \end{minipage}\hfill
    {\poplogo\fontsize{22}{22}\selectfont\color{popcream}OnBuch\color{poporange}+}
  \end{tcolorbox}%
}

% ============================================================
% Page de garde
% #1 titre · #2 matière · #3 niveau · #4 module · #5 n° leçon · #6 durée ·
% #7 description / objectifs (texte ou itemize)
% ============================================================
\newcommand{\pagegarde}[7]{%
  \thispagestyle{empty}
  \newgeometry{a4paper,margin=1.7cm,top=1.6cm,bottom=1.4cm}%
  \begin{tikzpicture}[remember picture,overlay]
    \fill[poporange!18] ([xshift=-3.2cm,yshift=-3.6cm]current page.north east) circle (4.6cm);
    \fill[poppurple!14] ([xshift=2.4cm,yshift=7.6cm]current page.south west) circle (3.8cm);
  \end{tikzpicture}%
  {\poplogo\fontsize{30}{30}\selectfont\color{popink}OnBuch\color{poporange}+}\hfill
  \raisebox{0.6ex}{\poppill{Cours signé \textbullet\ Premium}{popink}{popcream}}\par
  \vspace{1pt}\popsquiggle{poppurple}\par
  \vspace{8pt}
  \begin{tcolorbox}[enhanced,colback=poporange,colframe=popink,boxrule=2.8pt,arc=13pt,
      drop shadow=popink,left=18pt,right=18pt,top=12pt,bottom=13pt,after skip=10pt]
    \poppill{#2 \textbullet\ #3}{popink}{popcream}\hspace{5pt}\poppill{#4}{white}{popink}\par\vspace{8pt}
    {\popdisplay\fontsize{14}{14}\selectfont\color{popink}LEÇON #5}\par\vspace{4pt}
    {\raggedright\hyphenpenalty=10000\exhyphenpenalty=10000\popdisplay\fontsize{25}{27}\selectfont\color{popcream}\MakeUppercase{#1}\par}
    \vspace{8pt}
    {\popxbold\fontsize{9.5}{9.5}\selectfont\color{popink}\MakeUppercase{Durée conseillée : #6}}
  \end{tcolorbox}
  \begin{tcolorbox}[enhanced,colback=white,colframe=popink,boxrule=2.2pt,arc=10pt,
      drop shadow=popink,left=16pt,right=16pt,top=10pt,bottom=10pt,after skip=8pt]
    \poppill{Dans cette leçon}{poppurple}{white}\par\vspace{5pt}
    \popsemi\fontsize{9.3}{12.6}\selectfont\color{popink}#7
  \end{tcolorbox}
  \noindent\begin{minipage}[t]{0.487\linewidth}
    \begin{tcolorbox}[enhanced,colback=popgreen,colframe=popink,boxrule=2.2pt,arc=10pt,
        drop shadow=popink,left=11pt,right=11pt,top=10pt,bottom=10pt,height=3.1cm,valign=center]
      \tcbox[colback=white,colframe=popink,boxrule=1.3pt,arc=5pt,boxsep=0pt,nobeforeafter,
        left=3pt,right=3pt,top=3pt,bottom=3pt,valign=center]{\qrcode[height=1.9cm]{https://play.google.com/store/apps/details?id=com.luvvix.onbuch&hl=fr}}%
      \hspace{8pt}\begin{minipage}[c]{0.5\linewidth}
        {\popdisplay\fontsize{11}{12.5}\selectfont\color{white}\MakeUppercase{Télécharge OnBuch}}\par\vspace{4pt}
        {\raggedright\popsemi\fontsize{8.4}{11}\selectfont\color{white}Gratuit \textbullet\ Google Play\\Tuteur IA, annales, concours, résultats\par}
      \end{minipage}
    \end{tcolorbox}
  \end{minipage}\hfill
  \begin{minipage}[t]{0.487\linewidth}
    \begin{tcolorbox}[enhanced,colback=poppurple,colframe=popink,boxrule=2.2pt,arc=10pt,
        drop shadow=popink,left=11pt,right=11pt,top=10pt,bottom=10pt,height=3.1cm,valign=center]
      \tikz[baseline=-0.7ex]\node[circle,fill=white,draw=popink,line width=1.3pt,minimum size=1.6cm,inner sep=0]{\popdisplay\fontsize{24}{24}\selectfont\color{poppurple}+};%
      \hspace{8pt}\begin{minipage}[c]{0.6\linewidth}
        {\popdisplay\fontsize{11}{12.5}\selectfont\color{white}\MakeUppercase{Passe à OnBuch+}}\par\vspace{4pt}
        {\raggedright\popsemi\fontsize{8.4}{11}\selectfont\color{white}Tous les cours signés, Tuteur IA illimité, fiches PDF — onglet OnBuch+ de l'app\par}
      \end{minipage}
    \end{tcolorbox}
  \end{minipage}
  \vspace{8pt}
  \popcontenu
  \vfill
  \signatureonbuch
  \restoregeometry
  \newpage
}

% Rangée « Ce cours contient » (page de garde)
\newcommand{\poptile}[3]{% #1 couleur #2 gros texte #3 légende
  \begin{tcolorbox}[enhanced,colback=#1,colframe=popink,boxrule=2pt,arc=9pt,shadow={2.5pt}{-2.5pt}{0pt}{popink},
      left=6pt,right=6pt,top=9pt,bottom=9pt,halign=center,valign=center,height=1.75cm,width=\linewidth,nobeforeafter]
    {\popdisplay\fontsize{12.5}{13}\selectfont\color{white}\MakeUppercase{#2}}\par\vspace{3pt}
    {\centering\popsemi\fontsize{7.8}{9.5}\selectfont\color{white}#3\par}
  \end{tcolorbox}}
\newcommand{\popcontenu}{%
  \par\noindent{\popxboldls\fontsize{7.5}{7.5}\selectfont\color{popmuted}CE COURS SIGNÉ CONTIENT}\par\vspace{7pt}
  \noindent\begin{minipage}[t]{0.235\linewidth}\poptile{popblue}{Cours}{complet, illustré, pas à pas}\end{minipage}\hfill
  \begin{minipage}[t]{0.235\linewidth}\poptile{poporange}{Méthodes}{et exercices résolus}\end{minipage}\hfill
  \begin{minipage}[t]{0.235\linewidth}\poptile{poppink}{Exercices}{type examen + corrigés}\end{minipage}\hfill
  \begin{minipage}[t]{0.235\linewidth}\poptile{popgreen}{Fiche bilan}{l'essentiel à retenir}\end{minipage}\par}

% Sommaire
\newtcolorbox{sommaire}{enhanced,colback=white,colframe=popink,boxrule=2.2pt,arc=10pt,
  drop shadow=popink,left=18pt,right=18pt,top=13pt,bottom=13pt,before skip=6pt,after skip=20pt,
  before upper={\poppill{Sommaire}{poppurple}{white}\par\vspace{9pt}\popsemi\fontsize{10.6}{14.5}\selectfont\color{popink}}}

% Signature de fin de cours — poussée tout en bas de la dernière page
\newcommand{\findecours}{%
  \par\vfill\nopagebreak
  \begin{center}\popsquiggle{poporange}\end{center}\vspace{4pt}
  \signatureonbuch
}
\AtBeginDocument{\def\llbracket{\mathopen{[\mkern-3.5mu[}}\def\rrbracket{\mathclose{]\mkern-3.5mu]}}} % intervalles d'entiers (glyphe ⟦ absent de la police)
% Glyphes absents de Fira Math : redéfinis à partir de glyphes disponibles
\AtBeginDocument{%
  \def\setminus{\mathbin{\backslash}}\let\smallsetminus\setminus
  \def\vdots{\mathord{\vbox{\baselineskip3.2pt\lineskiplimit0pt\kern4pt\hbox{.}\hbox{.}\hbox{.}}}}%
  \def\ddots{\mathinner{\mkern1mu\raise6.5pt\hbox{.}\mkern2mu\raise3.6pt\hbox{.}\mkern2mu\raise0.7pt\hbox{.}\mkern1mu}}%
  \def\triangle{\mathord{\scalebox{0.9}{$\symup{\Delta}$}}}%
  \def\nabla{\mathord{\raisebox{\depth}{\scalebox{1}[-1]{$\symup{\Delta}$}}}}%
  \def\checkmark{\popcheck{popdark}}%
  \def\star{\mathbin{\ast}}\def\bigstar{\mathord{\ast}}%
  \def\diamond{\mathbin{\scalebox{0.6}{$\Diamond$}}}%
  \def\bigcup{\mathop{\mathchoice{\raisebox{-0.3em}{\scalebox{1.7}{$\cup$}}}{\scalebox{1.25}{$\cup$}}{\cup}{\cup}}\displaylimits}%
  \def\bigcap{\mathop{\mathchoice{\raisebox{-0.3em}{\scalebox{1.7}{$\cap$}}}{\scalebox{1.25}{$\cap$}}{\cap}{\cap}}\displaylimits}%
  \def\lesseqqgtr{\mathrel{\lessgtr}}\def\bigoplus{\mathop{\oplus}\displaylimits}%
}
\providecommand{\ssi}{si et seulement si} % abréviation fréquente
\AtBeginDocument{\providecommand{\wideparen}{\overparen}} % arc de cercle

%% FIGFIX-BEGIN v5 — correctifs globaux des figures (docs/onbuch-generation/tools/figfix)
\usepackage{adjustbox}\usepackage{etoolbox}\tcbuselibrary{raster}
% toute figure TikZ/pgfplots plus large que la ligne est réduite à la largeur disponible
\BeforeBeginEnvironment{tikzpicture}{\begin{adjustbox}{max width=\linewidth}}
\AfterEndEnvironment{tikzpicture}{\end{adjustbox}}
% légendes pgfplots lisibles par-dessus les courbes
\pgfplotsset{every axis/.append style={legend style={fill=white,fill opacity=0.92,text opacity=1,draw=black!15,inner sep=3pt}}}
% étiquettes d'arêtes (syntaxe edge["..."]) avec fond blanc
\tikzset{every edge quotes/.append style={fill=white,inner sep=1.2pt}}
%% FIGFIX-END
\begin{document}

\pagegarde{Lesson 12 — Writing: Compositions about taking part in recreational activities}{Anglais}{1ère A}{Chapter 1 : Family and Social Life}{ENA12}{2 h}{Cette leçon d'expression écrite apprend aux élèves de 1ère A à planifier, rédiger et corriger une composition sur le thème « taking part in recreational activities ». Elle couvre la structure introduction-développement-conclusion, les connecteurs logiques, les temps verbaux et le vocabulaire utiles, avec un modèle commenté et une grille d'évaluation.
\vspace{4pt}
\begin{multicols}{2}\small
\begin{itemize}
\item À la fin de cette leçon, je sais identifier la structure d'une composition (introduction, développement, conclusion)
\item À la fin de cette leçon, je sais utiliser le vocabulaire des loisirs et des activités récréatives
\item À la fin de cette leçon, je sais employer les connecteurs logiques pour organiser mes idées (first, moreover, however, finally...)
\item À la fin de cette leçon, je sais choisir les temps verbaux adaptés (present simple, present continuous, past simple)
\item À la fin de cette leçon, je sais rédiger un plan avant d'écrire ma composition
\item À la fin de cette leçon, je sais exprimer mes goûts et préférences avec des structures variées (I enjoy, I am fond of, I am keen on...)
\end{itemize}\end{multicols}}

\begin{sommaire}
\begin{multicols}{2}
\begin{enumerate}
\item Découverte : lire un modèle de composition
\item La structure d'une composition
\item Connecteurs logiques et temps verbaux
\item Vocabulaire des activités récréatives
\item Planifier et rédiger sa composition
\item Relire, corriger et évaluer
\item Activité d'intégration
\item Méthodes et exercices résolus
\item Exercices
\item Corrigés des exercices
\item Fiche bilan
\end{enumerate}
\end{multicols}
\end{sommaire}

\begin{objectifs}
\begin{itemize}
\item À la fin de cette leçon, je sais identifier la structure d'une composition (introduction, développement, conclusion)
\item À la fin de cette leçon, je sais utiliser le vocabulaire des loisirs et des activités récréatives
\item À la fin de cette leçon, je sais employer les connecteurs logiques pour organiser mes idées (first, moreover, however, finally...)
\item À la fin de cette leçon, je sais choisir les temps verbaux adaptés (present simple, present continuous, past simple)
\item À la fin de cette leçon, je sais rédiger un plan avant d'écrire ma composition
\item À la fin de cette leçon, je sais exprimer mes goûts et préférences avec des structures variées (I enjoy, I am fond of, I am keen on...)
\item À la fin de cette leçon, je sais relire et corriger ma production (orthographe, grammaire, ponctuation)
\item À la fin de cette leçon, je sais éviter les erreurs fréquentes des francophones en expression écrite
\item À la fin de cette leçon, je sais évaluer une composition à l'aide d'une grille critériée
\end{itemize}
\end{objectifs}

\begin{prerequis}
\begin{itemize}
\item Le présent simple et le présent continu (Seconde)
\item Le vocabulaire de base des loisirs et des sports (Seconde)
\item L'expression des goûts : like, love, hate + -ing (Seconde)
\item Les connecteurs simples : and, but, because (début d'année)
\item La rédaction d'un paragraphe simple (Seconde)
\end{itemize}
\end{prerequis}

\newpage

\coursec{Découverte : lire un modèle de composition}

\textbf{Situation d'accroche.} \eng{Last weekend, your cousin in Bamenda spent Saturday playing football with friends, singing in the church choir on Sunday, and ended the weekend at a video game centre. When your English teacher asks you to write a composition about how you spend your free time, will you simply list activities, or will you organise your ideas like a real writer? In Britain and the USA, students are trained to plan, draft and revise their essays — and at the Cameroonian BAC, a well-structured composition can earn you top marks.}

« Le week-end dernier, ton cousin de Bamenda a passé son samedi à jouer au football avec ses amis, a chanté dans la chorale de l'église le dimanche, et a terminé le week-end dans une salle de jeux vidéo. Quand ton professeur d'anglais te demande de rédiger une composition sur la façon dont tu passes ton temps libre, vas-tu simplement énumérer des activités, ou organiser tes idées comme un vrai écrivain ? En Grande-Bretagne et aux États-Unis, les élèves apprennent à planifier, rédiger et relire leurs essais — et au BAC camerounais, une composition bien structurée peut te rapporter les meilleures notes. »

\textbf{Problématique.} Comment passer d'une simple liste d'activités (\eng{I play football. I read. I watch TV.}) à une véritable \cle{composition} claire, organisée et agréable à lire ? Pour le découvrir, nous allons d'abord \textbf{lire et analyser un modèle} rédigé par un élève camerounais, afin de dégager les secrets d'une bonne composition.

\courssub{Lecture du texte modèle}

\begin{methode}[Comment aborder la lecture d'un modèle]
Avant de lire, adopte la bonne démarche :
\begin{enumerate}
\item \textbf{Première lecture silencieuse} : lis le texte entièrement pour saisir le sens global, sans t'arrêter sur les mots inconnus.
\item \textbf{Deuxième lecture à voix haute} : lis phrase par phrase, en respectant les pauses (virgules, points). Cela entraîne ta prononciation et ta fluidité.
\item \textbf{Repérage} : identifie les paragraphes et demande-toi : \emph{quelle est la fonction de chaque paragraphe ?}
\item \textbf{Soulignement} : souligne les mots de liaison (\eng{first of all, moreover, however, in conclusion}) et entoure le vocabulaire utile que tu pourrais réutiliser.
\end{enumerate}
\end{methode}

Lis maintenant le texte modèle ci-dessous, d'abord silencieusement, puis à voix haute.

\begin{textebox}[My Favourite Recreational Activity — a model composition]
In my free time, I enjoy playing football with my friends in our neighbourhood in Yaoundé. Every Saturday afternoon, we meet on the small field behind the market, and we play until the sun goes down.

First of all, football keeps me fit and healthy. After a long week of classes and homework, running on the field helps me forget my stress. Moreover, football teaches me teamwork and discipline: I have learnt that we cannot win a match if everyone plays alone. We must pass the ball, encourage each other and respect the referee's decisions.

However, I also like quieter activities. On Sunday mornings, I sometimes read novels in English, because reading improves my vocabulary and my imagination. In the evening, I enjoy listening to music with my elder sister. She loves makossa, but I prefer gospel songs.

In conclusion, recreational activities are very important in my life. They help me relax after a long week at school, they bring me closer to my friends and family, and they make me a healthier and happier person. That is why I believe every student should take part in at least one recreational activity.
\end{textebox}

\textbf{Glossaire.}

\begin{center}
\begin{tabularx}{\linewidth}{|l|l|X|}
\hline
\rowcolor{popblueL} \textbf{English} & \textbf{Nature} & \textbf{Français} \\
\hline
free time & noun phrase & temps libre, loisirs \\
\hline
neighbourhood & noun & quartier \\
\hline
field & noun & terrain (de sport) \\
\hline
fit & adjective & en forme, en bonne santé \\
\hline
teamwork & noun & travail d'équipe \\
\hline
referee & noun & arbitre \\
\hline
elder sister & noun phrase & grande sœur \\
\hline
to bring closer & verb phrase & rapprocher \\
\hline
to take part in & verb phrase & participer à, prendre part à \\
\hline
\end{tabularx}
\end{center}

\begin{attention}[Prononciation et faux amis]
\begin{itemize}
\item \eng{neighbourhood} se prononce « né-beur-houd » : le \emph{-gh-} est muet.
\item \eng{actually} ne veut pas dire « actuellement » mais « en réalité » ; « actuellement » se dit \eng{currently} ou \eng{at present}.
\item \eng{to assist} ne signifie pas « assister à » (au sens d'être présent) mais « aider » ; « assister à un match » se dit \eng{to attend a match} ou \eng{to watch a match}.
\item Attention à \eng{sensible} (raisonnable) et \eng{sensitive} (sensible) : deux faux amis classiques.
\end{itemize}
\end{attention}

\courssub{Qu'est-ce qu'une composition ?}

\begin{definition}[Composition]
A \cle{composition} is \eng{a short piece of writing on a single subject, organised in paragraphs: an introduction, a body and a conclusion.} « Une composition est un court texte écrit sur un seul sujet, organisé en paragraphes : une introduction, un développement et une conclusion. »
\end{definition}

En français comme en anglais, une composition n'est \textbf{pas} une liste de phrases juxtaposées. C'est un texte \textbf{uni} : un seul sujet, une progression logique, des paragraphes qui ont chacun une fonction. La grande différence avec la rédaction française : en anglais, on privilégie la \textbf{clarté} et la \textbf{simplicité} — des phrases courtes et correctes valent mieux que de longues phrases compliquées pleines d'erreurs.

\courssub{Repérage des trois grandes parties}

Relis le texte modèle et observe : il contient \textbf{quatre paragraphes} qui remplissent \textbf{trois fonctions}.

\begin{popfigure}
\begin{tikzpicture}[node distance=0.55cm]
\node[draw=popblue, fill=popblueL, rounded corners, minimum width=8.5cm, minimum height=1.1cm, align=center] (intro) {\textbf{Introduction}\\ \small Paragraph 1 — presents the topic and the favourite activity};
\node[draw=popgreen, fill=popgreenL, rounded corners, minimum width=8.5cm, minimum height=1.6cm, align=center, below=of intro] (body) {\textbf{Body (2--3 paragraphs)}\\ \small Paragraphs 2--3 — reasons, examples, contrasting activities};
\node[draw=poporange, fill=poporangeL, rounded corners, minimum width=8.5cm, minimum height=1.1cm, align=center, below=of body] (concl) {\textbf{Conclusion}\\ \small Paragraph 4 — summary and personal opinion};
\draw[-{Stealth}, thick, popdark] (intro) -- (body);
\draw[-{Stealth}, thick, popdark] (body) -- (concl);
\end{tikzpicture}
\legende{La structure en trois parties d'une composition : introduction, développement, conclusion.}
\end{popfigure}

\begin{aretenir}[La fonction de chaque partie]
\begin{itemize}
\item \textbf{Introduction} : elle présente le sujet et capte l'attention du lecteur. Dans le modèle : \eng{In my free time, I enjoy playing football...} — l'auteur annonce immédiatement son activité préférée et le lieu (Yaoundé).
\item \textbf{Body (développement)} : deux ou trois paragraphes qui développent les idées — ici, les \emph{raisons} d'aimer le football (paragraphe 2) puis les \emph{autres activités} plus calmes (paragraphe 3).
\item \textbf{Conclusion} : elle résume et donne une opinion personnelle finale : \eng{That is why I believe every student should take part in at least one recreational activity.}
\end{itemize}
\end{aretenir}

\courssub{Questions de compréhension}

\begin{exercice}[Comprehension questions]
Réponds en anglais, par des phrases complètes.

\textbf{A. Global comprehension}
\begin{enumerate}
\item What is the subject of the composition?
\item Who is speaking, and what is his or her point of view?
\item What is the writer's favourite recreational activity?
\end{enumerate}

\textbf{B. Detailed comprehension}
\begin{enumerate}
\item Where and when does the writer play football?
\item Give two reasons why the writer loves football.
\item Which quieter activities does the writer mention?
\item What type of music does the writer's sister love?
\item In the conclusion, what advice does the writer give to every student?
\end{enumerate}

\textbf{C. Focus on link words}
\begin{enumerate}
\item Underline all the link words (connectors) in the model text.
\item Which connector introduces an addition? Which one introduces a contrast?
\end{enumerate}
\end{exercice}

\begin{corrige}[Comprehension questions]
\textbf{A.} 1. The subject is recreational activities and how the writer spends his free time. 2. A student in Yaoundé is speaking; he gives his personal point of view (first person \eng{I}). 3. His favourite activity is playing football with his friends.

\textbf{B.} 1. He plays on the small field behind the market, every Saturday afternoon. 2. Football keeps him fit and healthy, and it teaches him teamwork and discipline. 3. Reading novels in English and listening to music. 4. She loves makossa. 5. He says every student should take part in at least one recreational activity.

\textbf{C.} 1. Connectors: \eng{first of all, moreover, however, in conclusion, that is why} (plus \eng{because, but, and}). 2. Addition: \eng{moreover}. Contrast: \eng{however}.
\end{corrige}

\courssub{Les connecteurs du texte modèle : observation}

Observons comment les connecteurs fonctionnent dans le modèle.

\begin{center}
\begin{tabularx}{\linewidth}{|l|X|X|}
\hline
\rowcolor{popblueL} \textbf{Connecteur} & \textbf{Fonction} & \textbf{Exemple du texte} \\
\hline
\eng{first of all} & première idée & \eng{First of all, football keeps me fit.} \\
\hline
\eng{moreover} & ajout & \eng{Moreover, football teaches me teamwork.} \\
\hline
\eng{however} & contraste & \eng{However, I also like quieter activities.} \\
\hline
\eng{because} & cause & \eng{...because reading improves my vocabulary.} \\
\hline
\eng{in conclusion} & conclusion & \eng{In conclusion, recreational activities are important.} \\
\hline
\eng{that is why} & conséquence & \eng{That is why I believe every student should...} \\
\hline
\end{tabularx}
\end{center}

\begin{attention}[Erreurs fréquentes des francophones]
\begin{itemize}
\item Ne traduis pas « d'abord » par \eng{at first} dans une énumération d'arguments : utilise \eng{first of all} ou \eng{firstly}. \eng{At first} signifie « au début » (opposé à « ensuite » dans le temps).
\item En anglais, on met une \textbf{virgule} après le connecteur en début de phrase : \eng{Moreover, it teaches me discipline.} — et non \eng{Moreover it teaches me discipline}.
\item N'écris jamais \eng{but} et \eng{however} dans la même phrase pour la même opposition : un seul connecteur de contraste suffit.
\item \eng{Because} et \eng{so} ne s'emploient pas ensemble : dis \eng{I was tired, so I slept.} ou \eng{Because I was tired, I slept.} — jamais \eng{Because I was tired, so I slept.}
\end{itemize}
\end{attention}

\courssub{Exercice résolu : analyser un paragraphe}

\begin{exoresolu}[Identify the function of each paragraph]
\textbf{Énoncé.} Voici les quatre paragraphes du texte modèle, mélangés. Remets-les dans l'ordre logique et indique la fonction de chacun (introduction, body, conclusion).

\begin{enumerate}
\item[A.] \eng{In conclusion, recreational activities are very important in my life. They help me relax after a long week at school.}
\item[B.] \eng{First of all, football keeps me fit and healthy. Moreover, football teaches me teamwork and discipline.}
\item[C.] \eng{In my free time, I enjoy playing football with my friends in our neighbourhood in Yaoundé.}
\item[D.] \eng{However, I also like quieter activities. On Sunday mornings, I sometimes read novels in English.}
\end{enumerate}

\tcblower
\textbf{Solution détaillée.}
\begin{enumerate}
\item Le paragraphe \textbf{C} ouvre le texte : il présente le sujet (\eng{in my free time}) et l'activité préférée. C'est l'\textbf{introduction}.
\item Le paragraphe \textbf{B} commence par \eng{first of all} : il donne les premières raisons. C'est le \textbf{premier paragraphe du body}.
\item Le paragraphe \textbf{D} commence par \eng{however} : il introduit un contraste (activités calmes). C'est le \textbf{deuxième paragraphe du body}.
\item Le paragraphe \textbf{A} commence par \eng{in conclusion} : c'est la \textbf{conclusion}.
\end{enumerate}
Ordre correct : \textbf{C — B — D — A}. Les connecteurs en début de paragraphe sont les indices qui permettent de reconstituer l'ordre.
\end{exoresolu}

\courssub{Les caractéristiques d'une bonne composition}

À partir de notre analyse, nous pouvons dégager les qualités qui font une bonne composition.

\begin{aretenir}[What makes a good composition?]
\begin{itemize}
\item \textbf{Clarity (clarté)} : des phrases courtes et correctes ; une idée par phrase.
\item \textbf{Organisation} : trois parties visibles — introduction, body, conclusion — avec un paragraphe par idée.
\item \textbf{Linking (liaison)} : des connecteurs variés qui guident le lecteur (\eng{first of all, moreover, however, in conclusion}).
\item \textbf{Rich vocabulary (richesse du vocabulaire)} : des mots précis sur le thème (\eng{teamwork, discipline, to relax, to take part in}) au lieu de répéter \eng{good, nice, like}.
\item \textbf{Correct grammar} : des temps verbaux cohérents (ici, le présent simple pour les habitudes : \eng{I enjoy, we meet, I read}).
\item \textbf{Personal touch} : une opinion personnelle et des détails concrets (Yaoundé, le marché, makossa) qui rendent le texte vivant.
\end{itemize}
\end{aretenir}

\begin{savaistu}[La composition au BAC]
Au baccalauréat camerounais, l'épreuve d'anglais comporte une partie \eng{essay writing}. Ta composition n'est pas notée seulement sur tes idées : le correcteur évalue quatre critères — le \textbf{contenu} (content), l'\textbf{organisation} (organisation), la \textbf{grammaire} (grammar) et le \textbf{vocabulaire} (vocabulary). Un élève aux idées simples mais au texte bien structuré, sans fautes et avec des connecteurs variés obtient souvent une meilleure note qu'un élève aux idées brillantes mais désorganisées. En Grande-Bretagne et aux États-Unis, les élèves apprennent dès le collège la méthode \eng{plan — draft — revise} (planifier — rédiger — relire) : c'est exactement la méthode que tu vas pratiquer dans cette leçon.
\end{savaistu}

\begin{experience}[Class activity — reading aloud and spotting connectors]
Par groupes de quatre :
\begin{enumerate}
\item Chaque élève lit un paragraphe du texte modèle à voix haute, avec une bonne prononciation et les bonnes pauses.
\item Les autres membres du groupe lèvent la main à chaque connecteur entendu et le notent.
\item Comparez vos listes : le groupe qui a repéré le plus de connecteurs (\eng{first of all, moreover, however, because, but, in conclusion, that is why...}) gagne.
\item Pour finir, chaque groupe dit à voix haute la fonction de chaque paragraphe : \eng{Paragraph one is the introduction...}
\end{enumerate}
\end{experience}

\textbf{Bilan de la section.} Tu as lu un modèle de composition, repéré ses trois parties (introduction, body, conclusion), identifié le sujet et le point de vue de l'auteur, souligné les connecteurs et dégagé les qualités d'une bonne composition : clarté, organisation, liaison et richesse du vocabulaire. Dans la prochaine section, nous étudierons en détail la structure d'une composition et le rôle précis de chaque paragraphe.

\coursec{La structure d'une composition}

Dans la section précédente, nous avons lu et commenté un modèle de composition intitulé \eng{“Why I love playing football after school”}. Nous avons remarqué que ce texte n'était pas écrit « d'un seul jet » : il était clairement divisé en plusieurs paragraphes, chacun jouant un rôle précis. Dans cette section, nous allons démonter cette machine bien huilée pour comprendre \cle{la structure d'une composition} : l'\cle{introduction}, le \cle{développement} (\eng{body}) et la \cle{conclusion}. Maîtriser cette architecture, c'est déjà gagner la moitié des points à l'examen.

\courssub{Vue d'ensemble : les trois parties obligatoires}

\begin{propriete}[La règle d'or de la composition]
Une composition (\eng{a composition} ou \eng{an essay}) se compose toujours de \textbf{trois parties} :
\begin{enumerate}
\item une \cle{introduction} (2 à 3 phrases) ;
\item un \cle{développement} (\eng{body}) de 2 à 3 paragraphes ;
\item une \cle{conclusion} (2 à 3 phrases).
\end{enumerate}
Règle fondamentale : \textbf{chaque paragraphe développe UNE seule idée principale}. On ne mélange jamais deux idées différentes dans le même paragraphe, et on ne répartit jamais une même idée sur deux paragraphes.
\end{propriete}

\begin{popfigure}
\begin{tikzpicture}[
node distance=0.55cm,
box/.style={rectangle, rounded corners=3pt, draw=popblue, fill=popblueL, text width=8.6cm, align=center, inner sep=5pt, font=\small},
lab/.style={font=\small\bfseries, text=popdark}
]
\node[box, align=center] (intro){\textbf{Introduction}\\ hook (accroche) + topic (sujet annoncé)};
\node[box, below=of intro, align=center] (b1){\textbf{Body paragraph 1}\\ idea 1 (topic sentence) + arguments + example};
\node[box, below=of b1, align=center] (b2){\textbf{Body paragraph 2}\\ idea 2 (topic sentence) + arguments + example};
\node[box, below=of b2, align=center] (conc){\textbf{Conclusion}\\ summary (résumé) + personal opinion (opinion)};
\draw[-{Stealth}, thick, poporange] (intro) -- (b1);
\draw[-{Stealth}, thick, poporange] (b1) -- (b2);
\draw[-{Stealth}, thick, poporange] (b2) -- (conc);
\end{tikzpicture}
\legende{Organigramme de la structure d'une composition : de l'accroche à l'opinion finale.}
\end{popfigure}

\begin{attention}[L'erreur qui coûte cher à l'examen]
Ne jamais écrire toute la composition \textbf{en un seul bloc}. Les examinateurs sanctionnent systématiquement l'absence de paragraphes : un texte sans alinéas donne l'impression d'un désordre d'idées, même si le contenu est bon. À chaque nouveau paragraphe, on va à la ligne et on marque un \cle{alinéa} (un léger retrait en début de ligne). Pensez aussi à l'\cle{équilibre} : un paragraphe de 2 lignes suivi d'un paragraphe de 15 lignes est maladroit. Visez des paragraphes de longueur comparable (4 à 6 phrases chacun).
\end{attention}

\courssub{L'introduction : présenter le sujet et capter l'attention}

Observons d'abord trois introductions possibles pour une composition sur les activités récréatives :

\begin{exemplebox}[Trois façons de commencer]
\textbf{1. Avec une question :} \eng{What do you do when the school bell rings at 4 p.m.? Many students in Yaoundé rush to the football field, and I am one of them.}\\
« Que fais-tu quand la cloche sonne à 16 h ? Beaucoup d'élèves de Yaoundé se précipitent sur le terrain de football, et j'en fais partie. »

\textbf{2. Avec un fait général :} \eng{All over the world, young people need moments of relaxation after long hours of study. In my school, recreational activities play a central role in our daily life.}\\
« Partout dans le monde, les jeunes ont besoin de moments de détente après de longues heures d'étude. Dans mon école, les activités récréatives jouent un rôle central dans notre vie quotidienne. »

\textbf{3. Avec une courte anecdote :} \eng{Last term, I scored the winning goal during the inter-class final. That day, I understood how much sport matters to me.}\\
« Le trimestre dernier, j'ai marqué le but de la victoire pendant la finale inter-classes. Ce jour-là, j'ai compris combien le sport compte pour moi. »
\end{exemplebox}

\begin{definition}[Hook et topic]
L'introduction contient deux éléments :
\begin{itemize}
\item le \cle{hook} (« l'accroche ») : la première phrase, qui attire l'attention du lecteur (question, fait général, anecdote) ;
\item le \cle{topic} (« le sujet ») : la dernière phrase de l'introduction, qui annonce clairement ce dont la composition va parler.
\end{itemize}
\end{definition}

\begin{methode}[Écrire une bonne introduction en trois étapes]
\begin{enumerate}
\item \textbf{Phrase générale} : commencez par une remarque large sur les loisirs ou la vie des jeunes (\eng{Recreational activities are important for every student.}).
\item \textbf{Resserrement} : rapprochez-vous de votre situation (votre école, votre quartier, votre famille).
\item \textbf{Annonce du sujet} : terminez par une phrase qui présente le thème précis de votre composition (\eng{In this composition, I will explain why playing football is my favourite recreational activity.}).
\end{enumerate}
Comptez 2 à 3 phrases au total : une introduction trop longue déséquilibre toute la composition.
\end{methode}

\courssub{Le développement : une idée par paragraphe}

Le développement (\eng{the body}) est le cœur de la composition. Il contient en général \textbf{2 ou 3 paragraphes}. Chaque paragraphe suit une architecture interne rigoureuse :

\begin{propriete}[La structure interne d'un paragraphe de développement]
Un bon paragraphe contient :
\begin{enumerate}
\item une \cle{topic sentence} (« phrase d'introduction du paragraphe ») : la première phrase, qui annonce l'idée principale du paragraphe ;
\item des \cle{arguments} (\eng{supporting sentences}) : 2 ou 3 phrases qui expliquent ou justifient cette idée ;
\item un \cle{exemple} concret (\eng{for example}, \eng{for instance}) qui illustre l'idée.
\end{enumerate}
Formule à retenir : \textbf{topic sentence + arguments + exemple}.
\end{propriete}

\begin{exemplebox}[Un paragraphe de développement bien construit]
\eng{First of all, football keeps me fit.} (« Tout d'abord, le football me garde en forme. ») \emph{— topic sentence.}\\
\eng{When I play every evening, I run, jump and breathe fresh air. As a result, I sleep well and I feel strong in class the next morning.} \emph{— arguments.}\\
\eng{For example, since I joined the school team, I no longer feel tired during the long mathematics lessons.} \emph{— exemple.}

Autre exemple de topic sentence : \eng{Moreover, recreational activities help me make new friends.} (« De plus, les activités récréatives m'aident à me faire de nouveaux amis. »)
\end{exemplebox}

\begin{attention}[Piège fréquent des francophones]
Ne traduisez pas mot à mot le français « le football me garde en forme » par \eng{football guards me in form} : c'est un calque incorrect. Le verbe anglais est \eng{to keep} : \eng{football keeps me fit}. De même, on dit \eng{to make friends} (« se faire des amis ») et non \eng{to win friends}. Attention aussi à l'ordre des mots : \eng{I play football every evening} et non \eng{I play every evening football} si la phrase continue ensuite.
\end{attention}

\courssub{La conclusion : résumer et donner son opinion}

La conclusion est la dernière image que le lecteur garde de votre travail. Elle remplit deux fonctions :

\begin{propriete}[Les deux fonctions de la conclusion]
\begin{enumerate}
\item \textbf{Résumer} (\eng{to sum up}) les idées principales du développement, sans répéter les mêmes phrases mot pour mot ;
\item \textbf{Donner son opinion personnelle} (\eng{in my opinion}, \eng{I believe that}, \eng{to my mind}) ou ouvrir sur un conseil, un souhait.
\end{enumerate}
La conclusion ne doit \textbf{jamais} introduire une idée nouvelle qui n'a pas été traitée dans le développement.
\end{propriete}

\begin{exemplebox}[Une conclusion modèle]
\eng{In conclusion, recreational activities help us relax.} (« En conclusion, les activités récréatives nous aident à nous détendre. »)\\
\eng{They keep our bodies healthy and our minds fresh. In my opinion, every school in Cameroon should give students more time for sport and games.}\\
« Elles gardent notre corps en bonne santé et notre esprit alerte. À mon avis, chaque école du Cameroun devrait accorder aux élèves plus de temps pour le sport et les jeux. »
\end{exemplebox}

\courssub{Le titre de la composition}

\begin{aretenir}[Règles pour le titre]
Le titre (\eng{the title}) doit être :
\begin{itemize}
\item \textbf{court} : de 3 à 8 mots environ ;
\item \textbf{en rapport direct avec le sujet} : il annonce le thème, pas autre chose ;
\item écrit \textbf{sans point final} et sans guillemets ;
\item avec une majuscule aux mots importants : \eng{My Favourite Recreational Activity}, \eng{Why Students Need Sports and Games}.
\end{itemize}
Si le sujet d'examen impose un titre, recopiez-le exactement tel quel.
\end{aretenir}

\courssub{Tableau récapitulatif : les trois parties et leurs fonctions}

\begin{center}
\begin{tabularx}{\linewidth}{|l|X|X|}
\hline
\rowcolor{popblueL} \textbf{Partie} & \textbf{English} & \textbf{Français / fonction} \\
\hline
Introduction & hook + topic sentence & accroche + annonce du sujet (2-3 phrases) \\
\hline
Body paragraph 1 & topic sentence + arguments + example & idée 1 développée et illustrée \\
\hline
Body paragraph 2 & topic sentence + arguments + example & idée 2 développée et illustrée \\
\hline
Conclusion & summary + personal opinion & résumé + avis personnel (2-3 phrases) \\
\hline
Title & a short, relevant title & titre court, en rapport avec le sujet \\
\hline
\end{tabularx}
\end{center}

\begin{center}
\begin{tabularx}{\linewidth}{|X|X|X|}
\hline
\rowcolor{popblueL} Français & English & Remarque \\
\hline
Tout d'abord & First of all & ouvre le 1\textsuperscript{er} paragraphe \\
\hline
De plus / Ensuite & Moreover / Then & ouvre le 2\textsuperscript{e} paragraphe \\
\hline
Par exemple & For example & introduit l'exemple \\
\hline
En conclusion & In conclusion & ouvre la conclusion \\
\hline
À mon avis & In my opinion & exprime l'opinion personnelle \\
\hline
\end{tabularx}
\end{center}

\courssub{Exercice résolu et entraînement}

\begin{exoresolu}[Remettre une composition dans l'ordre]
Énoncé — Les cinq phrases suivantes forment une composition complète, mais elles sont mélangées. Remettez-les dans l'ordre logique (introduction, développement, conclusion) et justifiez.

a) \eng{In conclusion, dancing is more than a hobby for me; it is a source of joy.}\\
b) \eng{Moreover, dancing helps me forget my worries after difficult school days.}\\
c) \eng{Many students enjoy sports, but my favourite recreational activity is traditional dancing.}\\
d) \eng{For instance, last Saturday, I danced with my friends at a school party and I felt completely free.}\\
e) \eng{First of all, dancing keeps my body flexible and strong.}
\tcblower
Solution détaillée :
\begin{enumerate}
\item \textbf{Phrase c} : c'est l'introduction. Elle oppose une idée générale (\eng{many students enjoy sports}) au sujet précis (\eng{traditional dancing}) : c'est le topic.
\item \textbf{Phrase e} : le connecteur \eng{first of all} signale le premier paragraphe du développement (idée 1 : le corps).
\item \textbf{Phrase b} : le connecteur \eng{moreover} signale le deuxième paragraphe (idée 2 : le moral).
\item \textbf{Phrase d} : \eng{for instance} introduit l'exemple qui illustre la phrase b ; elle suit donc immédiatement b.
\item \textbf{Phrase a} : \eng{in conclusion} marque la conclusion, avec un résumé et une opinion implicite.
\end{enumerate}
Ordre correct : \textbf{c — e — b — d — a}.
\end{exoresolu}

\begin{exercice}[À vous de structurer]
Sujet : \eng{Write a composition about your favourite recreational activity.}
\begin{enumerate}
\item Écrivez un titre court et adapté.
\item Rédigez uniquement le plan : la phrase d'accroche (\eng{hook}), les deux \eng{topic sentences} du développement et la phrase de conclusion.
\item Indiquez pour chaque paragraphe l'exemple concret que vous utiliserez.
\end{enumerate}
\end{exercice}

\begin{corrige}[Exercice : éléments de réponse]
Exemple de plan attendu :
\begin{itemize}
\item Titre : \eng{Why I Love Reading Storybooks}.
\item Hook : \eng{Some students relax by playing games, but I relax with a good book.}
\item Topic sentence 1 : \eng{First of all, reading opens my mind to new worlds and new ideas.}
\item Topic sentence 2 : \eng{Moreover, reading improves my English vocabulary every day.}
\item Exemples : un livre précis lu récemment ; un mot nouveau appris grâce à la lecture.
\item Conclusion : \eng{In conclusion, reading is my best companion. In my opinion, every student should read at least one storybook every month.}
\end{itemize}
Vérifiez que chaque paragraphe ne contient qu'une seule idée et que les paragraphes sont de longueur équilibrée.
\end{corrige}

\begin{savaistu}[Pourquoi les paragraphes ?]
Le mot \eng{paragraph} vient du grec \emph{paragraphos}, une petite marque que les scribes grecs traçaient dans la marge pour signaler un changement d'idée. Deux mille cinq cents ans plus tard, la règle n'a pas changé : un paragraphe, une idée. Dans les examens anglophones comme au BEPC et au Probatoire, la clarté de la structure (\eng{organisation}) représente une part importante de la note d'expression écrite — souvent autant que la grammaire elle-même.
\end{savaistu}

\coursec{Connecteurs logiques et temps verbaux}

Dans la section précédente, nous avons appris qu'une composition bien construite comporte trois parties : une \cle{introduction}, un \cle{développement} et une \cle{conclusion}. Mais une suite de phrases juxtaposées ne suffit pas : pour que le lecteur suive votre pensée, il faut \emph{relier} les idées entre elles. C'est le rôle des \cle{connecteurs logiques} (en anglais \eng{linking words} ou \eng{connectors}). De plus, pour parler de vos activités récréatives, vous devrez choisir le bon \cle{temps verbal} : habitudes, activités en cours, souvenirs passés. Cette section vous donne les deux outils indispensables : les connecteurs et les temps.

\courssub{Pourquoi les connecteurs sont-ils indispensables ?}

Lisez les deux paragraphes suivants. Lequel est le plus agréable à lire ?

\begin{textebox}[Two versions of the same paragraph]
\textbf{Version A (without connectors)}\\
I play football. I like basketball. I train on Saturdays. I am sometimes tired. I never miss a match. Sport is good for my health. I will continue.

\textbf{Version B (with connectors)}\\
I play football, \emph{but} I also like basketball. \emph{First of all}, I train every Saturday. \emph{However}, I am sometimes tired. \emph{Moreover}, I never miss a match \emph{because} sport is good for my health. \emph{In conclusion}, I will certainly continue.
\end{textebox}

La version B est clairement supérieure : les connecteurs montrent les \emph{relations logiques} entre les idées (opposition, ordre, addition, cause, conclusion). Les examinateurs du Probatoire et du Baccalauréat accordent des points précis à la \cle{cohérence} du texte : une composition sans connecteurs perd facilement plusieurs points.

\begin{definition}[Connecteur logique]
Un \cle{connecteur logique} (\eng{linking word}) est un mot ou une expression qui relie deux phrases ou deux idées en précisant leur relation : addition, ordre chronologique, opposition, cause, conséquence ou conclusion.
\end{definition}

\begin{popfigure}
\begin{tikzpicture}[mindmap, grow cyclic,
  every node/.style={concept, font=\small},
  root concept/.append style={concept color=popblue, font=\bfseries\small},
  level 1 concept/.append style={level distance=3.4cm, sibling angle=90, font=\small},
  level 2 concept/.append style={level distance=2.6cm, font=\scriptsize}]
\node[root concept]{Connectors}
  child[concept color=popgreen] { node {Addition}
    child { node[concept color=popgreenL] {and, also} }
    child { node[concept color=popgreenL] {moreover} }
    child { node[concept color=popgreenL] {in addition} }
  }
  child[concept color=poporange] { node {Order}
    child { node[concept color=poporangeL] {first of all} }
    child { node[concept color=poporangeL] {then, next} }
    child { node[concept color=poporangeL] {finally} }
  }
  child[concept color=poppink] { node {Opposition}
    child { node[concept color=poppinkL] {but} }
    child { node[concept color=poppinkL] {however} }
    child { node[concept color=poppinkL] {whereas} }
  }
  child[concept color=poppurple] { node {Conclusion}
    child { node[concept color=poppurpleL] {in conclusion} }
    child { node[concept color=poppurpleL] {to sum up} }
    child { node[concept color=poppurpleL] {all in all} }
  };
\end{tikzpicture}
\legende{Carte mentale des principales familles de connecteurs logiques.}
\end{popfigure}

\courssub{Les connecteurs d'addition}

Les connecteurs d'addition servent à \emph{ajouter} une idée à une idée précédente, comme « et », « aussi », « de plus » en français.

\begin{definition}[Connecteurs d'addition]
\eng{and} (et), \eng{also} (aussi), \eng{moreover} (de plus), \eng{furthermore} (en outre), \eng{in addition} (en plus). Ils renforcent une idée en ajoutant un argument ou un détail.
\end{definition}

Observez la position de chaque connecteur dans la phrase :

\begin{exemplebox}[Exemples traduits]
\begin{itemize}
\item \eng{I play football and I swim on Sundays.} « Je joue au football et je nage le dimanche. »
\item \eng{I also enjoy dancing.} « J'aime aussi danser. » (\eng{also} se place avant le verbe principal.)
\item \eng{Sport keeps me fit. Moreover, it teaches me teamwork.} « Le sport me garde en forme. De plus, il m'apprend le travail en équipe. »
\item \eng{Furthermore, it is completely free.} « En outre, c'est entièrement gratuit. »
\item \eng{In addition, I meet many friends there.} « En plus, j'y rencontre beaucoup d'amis. »
\end{itemize}
\end{exemplebox}

\begin{attention}[La place de « also »]
En français, « aussi » peut commencer une phrase (« Aussi, j'aime danser »). En anglais, \eng{also} se place \emph{avant le verbe principal} et \emph{après l'auxiliaire} ou le verbe \eng{be} : \eng{I also play tennis.} / \eng{She is also a singer.} Ne dites jamais \eng{Also I play tennis} en tête de phrase dans une composition soignée ; préférez \eng{Moreover} ou \eng{In addition} en début de phrase.
\end{attention}

\courssub{Les connecteurs d'ordre}

Ils organisent le récit ou l'énumération dans le temps, comme « d'abord », « ensuite », « enfin ».

\begin{definition}[Connecteurs d'ordre]
\eng{first of all} (tout d'abord), \eng{then} (ensuite), \eng{next} (puis), \eng{after that} (après cela), \eng{finally} (finalement). Ils sont parfaits pour décrire le déroulement d'une journée d'activités ou d'une compétition.
\end{definition}

\begin{exemplebox}[Une journée sportive racontée]
\eng{First of all, we warmed up for ten minutes. Then, the coach divided us into two teams. Next, we played the first half. After that, we rested and drank water. Finally, our team won the match 2--1.}\\
« Tout d'abord, nous nous sommes échauffés pendant dix minutes. Ensuite, l'entraîneur nous a répartis en deux équipes. Puis, nous avons joué la première mi-temps. Après cela, nous nous sommes reposés et avons bu de l'eau. Finalement, notre équipe a gagné le match 2 à 1. »
\end{exemplebox}

\courssub{Les connecteurs d'opposition}

Ils expriment un contraste, comme « mais », « cependant », « alors que ».

\begin{definition}[Connecteurs d'opposition]
\eng{but} (mais), \eng{however} (cependant), \eng{whereas} (alors que), \eng{on the other hand} (d'un autre côté).
\end{definition}

\begin{exemplebox}[Exemples traduits]
\begin{itemize}
\item \eng{I like football, but I prefer basketball.} « J'aime le football, mais je préfère le basket-ball. »
\item \eng{Football is exciting. However, I also like reading.} « Le football est passionnant. Cependant, j'aime aussi lire. »
\item \eng{My brother likes video games, whereas I prefer outdoor games.} « Mon frère aime les jeux vidéo, alors que je préfère les jeux de plein air. »
\item \eng{Swimming is tiring. On the other hand, it is very relaxing.} « La natation est fatigante. D'un autre côté, elle est très relaxante. »
\end{itemize}
\end{exemplebox}

\begin{attention}[« However » n'est pas « mais »]
Erreur classique du francophone : utiliser \eng{however} comme \eng{but} au milieu d'une phrase, séparé par une simple virgule. C'est incorrect. \eng{However} est un adverbe : il commence généralement une \emph{nouvelle phrase} (ou suit un point-virgule) et est suivi d'une virgule.\\
Correct : \eng{I like football. However, I prefer basketball.}\\
Correct : \eng{I like football, but I prefer basketball.}\\
Incorrect : \eng{I like football, however I prefer basketball.}
\end{attention}

\courssub{Les connecteurs de cause et de conséquence}

Ils expriment la raison (\eng{because}, \eng{since}) ou le résultat (\eng{so}, \eng{therefore}, \eng{as a result}).

\begin{exemplebox}[Exemples traduits]
\begin{itemize}
\item \eng{I joined the choir because I love singing.} « Je me suis inscrit à la chorale parce que j'adore chanter. »
\item \eng{Since it was raining, we stayed indoors and played ludo.} « Comme il pleuvait, nous sommes restés à l'intérieur et avons joué au ludo. »
\item \eng{The stadium was closed, so we trained in the school yard.} « Le stade était fermé, donc nous nous sommes entraînés dans la cour de l'école. »
\item \eng{I practised every day; therefore, I won the race.} « Je m'entraînais chaque jour ; par conséquent, j'ai gagné la course. »
\item \eng{As a result, our team qualified for the final.} « En conséquence, notre équipe s'est qualifiée pour la finale. »
\end{itemize}
\end{exemplebox}

\begin{attention}[Jamais « because ... so » dans la même phrase]
En français on peut dire « parce qu'il pleuvait, donc nous sommes restés ». En anglais, on choisit \emph{un seul} connecteur : soit \eng{Because it was raining, we stayed at home}, soit \eng{It was raining, so we stayed at home}. Jamais les deux ensemble. Remarquez aussi l'ordre des mots : \eng{at home} (sans article), et non \eng{at the home}.
\end{attention}

\courssub{Les connecteurs de conclusion}

Ils annoncent la fin de la composition : \eng{in conclusion} (en conclusion), \eng{to sum up} (pour résumer), \eng{all in all} (tout compte fait).

\begin{exemplebox}[Exemples traduits]
\begin{itemize}
\item \eng{In conclusion, recreational activities are essential for a balanced life.} « En conclusion, les activités récréatives sont essentielles pour une vie équilibrée. »
\item \eng{To sum up, sport keeps me healthy and happy.} « Pour résumer, le sport me garde en bonne santé et heureux. »
\item \eng{All in all, I would recommend dancing to everyone.} « Tout compte fait, je recommanderais la danse à tout le monde. »
\end{itemize}
\end{exemplebox}

\begin{center}
\begin{tabularx}{\linewidth}{|l|X|X|}
\hline
\rowcolor{popblueL} \textbf{Relation} & \textbf{English} & \textbf{Français} \\
\hline
Addition & and, also, moreover, furthermore, in addition & et, aussi, de plus, en outre, en plus \\
\hline
Ordre & first of all, then, next, after that, finally & tout d'abord, ensuite, puis, après cela, enfin \\
\hline
Opposition & but, however, whereas, on the other hand & mais, cependant, alors que, d'un autre côté \\
\hline
Cause & because, since & parce que, comme/puisque \\
\hline
Conséquence & so, therefore, as a result & donc, par conséquent, en conséquence \\
\hline
Conclusion & in conclusion, to sum up, all in all & en conclusion, pour résumer, tout compte fait \\
\hline
\end{tabularx}
\end{center}

\courssub{Les temps verbaux pour parler des activités récréatives}

Trois temps suffisent pour rédiger une excellente composition sur ce thème : le \cle{présent simple}, le \cle{présent continu} et le \cle{past simple}.

\begin{propriete}[Le présent simple : habitudes et goûts]
\textbf{Forme} : sujet + verbe de base (avec \emph{-s} à la 3\up{e} personne du singulier : \eng{he plays}, \eng{she sings}).\\
\textbf{Emploi} : habitudes, routines, goûts généraux, vérités permanentes. Marqueurs typiques : \eng{every Saturday}, \eng{always}, \eng{usually}, \eng{often}, \eng{never}.\\
\eng{She sings in the choir every Sunday.} « Elle chante à la chorale chaque dimanche. »\\
\eng{I play football every Saturday.} « Je joue au football chaque samedi. »
\end{propriete}

\begin{propriete}[Le présent continu : activité en cours]
\textbf{Forme} : sujet + \eng{be} (am/is/are) + verbe en \emph{-ing}.\\
\textbf{Emploi} : action en cours au moment où l'on parle, ou activité temporaire de la période présente. Marqueurs : \eng{now}, \eng{at the moment}, \eng{these days}, \eng{this term}.\\
\eng{I am learning to swim these days.} « J'apprends à nager en ce moment. »\\
\eng{We are preparing a drama competition this term.} « Nous préparons un concours de théâtre ce trimestre. »
\end{propriete}

\begin{propriete}[Le past simple : expérience passée terminée]
\textbf{Forme} : verbe régulier + \emph{-ed} (\eng{played}, \eng{visited}) ou forme irrégulière (\eng{won}, \eng{went}, \eng{took}).\\
\textbf{Emploi} : action passée, terminée, souvent datée. Marqueurs : \eng{last week}, \eng{yesterday}, \eng{two months ago}, \eng{in 2023}.\\
\eng{We visited the stadium last week.} « Nous avons visité le stade la semaine dernière. »\\
\eng{Last month, we won the match.} « Le mois dernier, nous avons gagné le match. »
\end{propriete}

\begin{center}
\begin{tabular}{|l|l|l|}
\hline
\rowcolor{popblueL} \textbf{Temps} & \textbf{Emploi} & \textbf{Marqueurs} \\
\hline
Present simple & habitudes, goûts & every day, always, usually \\
\hline
Present continuous & activité en cours & now, these days, at the moment \\
\hline
Past simple & expérience terminée & last week, yesterday, ago \\
\hline
\end{tabular}
\end{center}

\begin{attention}[Le « s » de la troisième personne]
N'oubliez jamais le \emph{-s} au présent simple avec \eng{he}, \eng{she}, \eng{it} : \eng{She play\textbf{s} tennis}, et non \eng{She play tennis}. C'est la faute la plus fréquente des élèves francophones, car le français ne marque pas cette différence à l'oral. Attention aux verbes terminés par \emph{-ch}, \emph{-sh}, \emph{-s}, \emph{-o} : ils prennent \emph{-es} : \eng{he watches}, \eng{she goes}, \eng{it finishes}.
\end{attention}

\begin{methode}[Varier ses connecteurs]
\begin{enumerate}
\item Avant de rédiger, préparez sur votre brouillon une liste de \textbf{8 à 10 connecteurs variés} (au moins un par famille : addition, ordre, opposition, cause, conclusion).
\item Interdisez-vous de répéter \eng{and} et \eng{but} plus d'une fois : remplacez-les par \eng{moreover}, \eng{in addition}, \eng{however}, \eng{whereas}.
\item Placez un connecteur d'ordre au début de chaque étape d'un récit (\eng{first of all... then... finally}).
\item À la relecture, soulignez chaque connecteur : s'il y a moins de six connecteurs dans votre composition, ajoutez-en.
\end{enumerate}
\end{methode}

\begin{exoresolu}[Choisir le bon connecteur et le bon temps]
Énoncé : complétez avec le connecteur ou la forme verbale correcte.\\
1. I love music. \_\_\_\_\_\_\_, I joined the school choir. (\eng{However} / \eng{Therefore})\\
2. She \_\_\_\_\_\_\_ (play) basketball every Friday.\\
3. Last holidays, we \_\_\_\_\_\_\_ (go) to the beach in Kribi.\\
4. I like chess, \_\_\_\_\_\_\_ my sister prefers dancing. (\eng{because} / \eng{whereas})\\
5. Look! The children \_\_\_\_\_\_\_ (dance) in the playground.
\tcblower
Solution détaillée :\\
1. \eng{Therefore} : la deuxième phrase est une conséquence de la première (j'aime la musique, donc je rejoins la chorale). \eng{However} exprimerait une opposition, ce qui n'a pas de sens ici.\\
2. \eng{plays} : présent simple avec marqueur d'habitude \eng{every Friday}, et \emph{-s} à la 3\up{e} personne.\\
3. \eng{went} : past simple du verbe irrégulier \eng{go}, marqueur \eng{last holidays}.\\
4. \eng{whereas} : opposition entre deux goûts différents ; \eng{because} exprimerait une cause, illogique ici.\\
5. \eng{are dancing} : présent continu, signalé par \eng{Look!} qui indique une action en cours.
\end{exoresolu}

\begin{exercice}[À vous de jouer]
Reliez les phrases avec le connecteur indiqué, puis mettez les verbes au bon temps.\\
1. (however) I am tired. / I will finish the game.\\
2. (because) I joined the drama club. / I love acting.\\
3. We \_\_\_\_\_\_\_ (win) the singing competition last term.\\
4. My friend \_\_\_\_\_\_\_ (watch) football every weekend.\\
5. (first of all / then / finally) racontez en trois phrases votre dernière sortie récréative.
\end{exercice}

\begin{corrige}[Exercice : connecteurs et temps]
1. \eng{I am tired. However, I will finish the game.} (nouvelle phrase + virgule après \eng{However}.)\\
2. \eng{I joined the drama club because I love acting.}\\
3. \eng{won} : past simple irrégulier de \eng{win}, marqueur \eng{last term}.\\
4. \eng{watches} : présent simple, 3\up{e} personne ; attention, \eng{watch} prend \emph{-es}.\\
5. Exemple : \eng{First of all, we took a bendskin (a motorbike taxi) to the stadium. Then, we watched the match and cheered our team. Finally, we ate grilled fish and went home happy.}
\end{corrige}

\begin{aretenir}[L'essentiel de la section]
\begin{itemize}
\item Les connecteurs donnent de la \cle{cohérence} à votre composition : addition (\eng{moreover}), ordre (\eng{first of all}), opposition (\eng{however}), cause (\eng{because}), conséquence (\eng{therefore}), conclusion (\eng{in conclusion}).
\item \eng{However} commence une nouvelle phrase suivi d'une virgule ; il ne remplace pas \eng{but} au milieu d'une phrase.
\item Présent simple = habitudes (\emph{-s} à la 3\up{e} personne) ; présent continu = action en cours (\eng{be + -ing}) ; past simple = expérience passée terminée.
\item Variez vos connecteurs : préparez une liste de 8 à 10 mots de liaison avant de rédiger.
\end{itemize}
\end{aretenir}

\coursec{Vocabulaire des activités récréatives}

Dans la section précédente, nous avons vu comment relier nos idées grâce aux connecteurs logiques (\eng{first, moreover, finally}) et comment choisir les bons temps verbaux. Mais de beaux connecteurs ne suffisent pas : pour écrire une composition riche sur les loisirs, il faut disposer d'un \cle{vocabulaire} précis et varié. C'est l'objectif de cette section : constituer votre « boîte à mots » sur les activités récréatives, apprendre à exprimer vos goûts et vos opinions, et éviter les pièges lexicaux qui guettent les francophones.

\courssub{Qu'est-ce qu'une activité récréative ?}

\begin{definition}[Recreational activity]
A \cle{recreational activity} is an activity done for pleasure and relaxation during free time. Une activité récréative est une activité pratiquée pour le plaisir et la détente pendant les loisirs. On parle aussi de \eng{leisure activity} ou \eng{hobby} (passe-temps).
\end{definition}

Observez ce court texte, écrit par un élève de Première de Yaoundé. Soulignez mentalement tous les mots qui désignent des activités.

\begin{textebox}[My free time, by Amina, Form Five student]
In my opinion, free time is a gift, and we should use it wisely. In my school in Yaoundé, students take part in many recreational activities. I am fond of singing in the school choir, and my best friend is keen on drama: she acts in every school play. During the long break, many boys play football on the field, while some girls prefer volleyball. I believe that sports help us to keep fit and to make friends.

At the weekend, my family enjoys outdoor activities. Last Saturday, we went hiking on Mount Fébé, and we had a picnic at the top. My father loves fishing, but I find it a little boring! In the evening, I like listening to music or reading novels. My little brother spends hours playing video games, which my mother does not really like. As far as I am concerned, recreational activities reduce stress after a busy week. They also develop our skills: chess, for example, teaches us to think carefully. A balanced life needs both work and play.
\end{textebox}

\textbf{Glossaire}

\begin{center}
\begin{tabularx}{\linewidth}{|l|l|X|}
\hline
\rowcolor{popblueL} Mot anglais & Nature & Traduction \\
\hline
choir & nom & chœur, chorale \\
\hline
drama & nom & théâtre (activité) \\
\hline
to keep fit & locution verbale & rester en forme \\
\hline
hiking & nom & randonnée \\
\hline
boring & adjectif & ennuyeux \\
\hline
to reduce stress & locution verbale & réduire le stress \\
\hline
skills & nom (pluriel) & compétences, aptitudes \\
\hline
\end{tabularx}
\end{center}

\textbf{Questions de compréhension}
\begin{enumerate}
\item Which cultural activities does Amina mention?
\item What did Amina's family do last Saturday?
\item According to Amina, what are two benefits of recreational activities?
\end{enumerate}

\courssub{Les quatre grandes familles d'activités}

Classons le vocabulaire en quatre catégories. Cette organisation vous aidera à structurer le développement de votre composition : un paragraphe par catégorie, par exemple.

\begin{popfigure}
\begin{tikzpicture}[
  col/.style={rectangle, rounded corners=3pt, draw=popline, fill=popcream, text width=3.2cm, align=center, font=\small, inner sep=5pt},
  head/.style={font=\small\bfseries, align=center}
]
\node[head, text=popblue] (h1) at (0,0) {Sports and games};
\node[col, below=2pt of h1, align=center]{football\\basketball\\volleyball\\athletics\\chess, ludo, scrabble};
\node[head, text=poppurple] (h2) at (4.3,0) {Cultural activities};
\node[col, below=2pt of h2, align=center]{singing in a choir\\dancing\\drama\\reading novels\\storytelling};
\node[head, text=popgreen] (h3) at (8.6,0) {Outdoor activities};
\node[col, below=2pt of h3, align=center]{hiking\\swimming\\fishing\\picnicking\\cycling};
\node[head, text=poporange] (h4) at (12.9,0) {Modern hobbies};
\node[col, below=2pt of h4, align=center]{playing video games\\watching films\\browsing the internet\\listening to music};
\end{tikzpicture}
\legende{Carte mentale lexicale : les quatre familles d'activités récréatives.}
\end{popfigure}

\begin{definition}[Quelques précisions de vocabulaire]
\begin{itemize}
\item \eng{athletics} (prononcé « az-lé-tiks ») = l'athlétisme (course, saut, lancer). Attention : c'est un nom singulier en anglais : \eng{Athletics is my favourite sport.}
\item \eng{chess} = les échecs ; \eng{ludo} = le jeu de ludo (jeu de société avec dés) ; \eng{scrabble} = le scrabble. Ces noms de jeux ne prennent pas d'article : \eng{We play chess at break time.}
\item \eng{storytelling} = l'art de raconter des histoires ; \eng{browsing the internet} = naviguer sur Internet.
\end{itemize}
\end{definition}

\courssub{Exprimer ses goûts : les verbes de préférence}

Pour parler de vos activités favorites, l'anglais dispose de plusieurs expressions de goût. Observez :

\eng{I am fond of dancing.} « J'adore danser. »

\eng{She is keen on reading novels.} « Elle aime beaucoup lire des romans. »

\eng{We enjoy playing volleyball.} « Nous aimons jouer au volley. »

\begin{propriete}[Règle d'or : le verbe qui suit est en -ing]
Toutes les expressions de goût sont suivies d'un nom ou d'un verbe en \cle{-ing} (gérondif), jamais de l'infinitif avec \eng{to} (sauf \eng{love/like} qui acceptent les deux) :
\begin{itemize}
\item \eng{enjoy} + -ing : \eng{I enjoy swimming.} (jamais \eng{I enjoy to swim})
\item \eng{love / like} + -ing (ou + to + base verbale) : \eng{She loves dancing.} / \eng{She loves to dance.}
\item \eng{be fond of} + -ing ; \eng{be keen on} + -ing ; \eng{be interested in} + -ing
\item \eng{prefer} + -ing : \eng{I prefer cycling to fishing.} — notez la préposition \eng{to} dans la comparaison.
\end{itemize}
\end{propriete}

\begin{center}
\begin{tabularx}{\linewidth}{|X|X|X|}
\hline
\rowcolor{popblueL} Expression & Traduction & Exemple \\
\hline
enjoy + -ing & aimer (bien) & I enjoy fishing. (J'aime pêcher.) \\
\hline
love + -ing / to & adorer & He loves playing chess. / He loves to play chess. \\
\hline
like + -ing / to & aimer & We like hiking. / We like to hike. \\
\hline
be fond of + -ing & adorer, être passionné de & I am fond of dancing. (J'adore danser.) \\
\hline
be keen on + -ing & être passionné de & She is keen on drama. (Elle est passionnée de théâtre.) \\
\hline
be interested in + -ing & s'intéresser à & They are interested in athletics. (Ils s'intéressent à l'athlétisme.) \\
\hline
prefer + -ing & préférer & I prefer swimming to cycling. (Je préfère la natation au vélo.) \\
\hline
\end{tabularx}
\end{center}

\begin{savaistu}[Fond of, keen on : toujours le -ing !]
\eng{Be fond of} et \eng{be keen on} sont suivis du verbe en \cle{-ing} : \eng{I am keen on swimming} (jamais \eng{keen on to swim}). Pourquoi ? Parce que \eng{of} et \eng{on} sont des prépositions, et en anglais, tout verbe placé après une préposition prend la forme en -ing : \eng{interested in reading}, \eng{good at singing}, \eng{tired of waiting}. Retenez ce mécanisme, il est valable pour toute la langue !
\end{savaistu}

\courssub{Exprimer son opinion}

Une composition sur les loisirs n'est pas une simple liste : vous devez donner votre avis. Voici les expressions d'opinion à utiliser à l'introduction et dans le développement.

\begin{center}
\begin{tabularx}{\linewidth}{|X|X|}
\hline
\rowcolor{popblueL} English & Français \\
\hline
In my opinion, ... & À mon avis, ... \\
\hline
I think (that) ... & Je pense que ... \\
\hline
I believe (that) ... & Je crois que ... \\
\hline
As far as I am concerned, ... & En ce qui me concerne, ... / Pour ma part, ... \\
\hline
In my view, ... & À mon sens, ... \\
\hline
\end{tabularx}
\end{center}

\eng{As far as I am concerned, playing video games is a waste of time.} « Pour ma part, jouer aux jeux vidéo est une perte de temps. »

\begin{attention}[Deux erreurs classiques de francophones]
\begin{itemize}
\item « I am agree » est \textbf{FAUX} — dites \eng{I agree} (je suis d'accord). En anglais, \eng{agree} est un verbe, pas un adjectif : on ne met donc pas \eng{am} devant. De même : \eng{I disagree} (je ne suis pas d'accord).
\item « To make sport » est \textbf{FAUX} — dites \eng{to do sport} ou \eng{to play a sport}. Règle pratique : \eng{play} + sport avec ballon ou jeu (\eng{play football, play chess}) ; \eng{do} + activité individuelle ou discipline (\eng{do athletics, do sport}) ; \eng{go} + activité en -ing (\eng{go swimming, go hiking, go fishing}).
\end{itemize}
\end{attention}

\courssub{Les bénéfices des loisirs}

Pour conclure une composition, il est très utile d'expliquer \emph{pourquoi} les activités récréatives sont importantes. Voici les expressions-clés :

\begin{itemize}
\item \eng{to relax} = se détendre — \eng{Listening to music helps me to relax.} « Écouter de la musique m'aide à me détendre. »
\item \eng{to keep fit} = rester en forme — \eng{Hiking helps us to keep fit.} « La randonnée nous aide à rester en forme. »
\item \eng{to make friends} = se faire des amis — \eng{Team sports help students to make friends.} « Les sports d'équipe aident les élèves à se faire des amis. »
\item \eng{to reduce stress} = réduire le stress — \eng{Reading novels reduces stress before exams.} « Lire des romans réduit le stress avant les examens. »
\item \eng{to develop skills} = développer des compétences — \eng{Scrabble develops our vocabulary skills.} « Le scrabble développe notre vocabulaire. »
\end{itemize}

\begin{exemplebox}[Un modèle de phrase à bénéfice]
\eng{In my opinion, recreational activities are important because they help us to relax, to keep fit and to make new friends.} « À mon avis, les activités récréatives sont importantes parce qu'elles nous aident à nous détendre, à rester en forme et à nous faire de nouveaux amis. » Remarquez la structure : opinion + \eng{because} + liste de bénéfices à l'infinitif.
\end{exemplebox}

\courssub{Exercice résolu et pièges à éviter}

\begin{exoresolu}[Complétez avec le vocabulaire de la leçon]
Énoncé — Complétez les phrases avec : \eng{keen on, fond of, keep fit, make friends, browsing}.
\begin{enumerate}
\item My sister is \dots\ singing in the church choir.
\item I am \dots\ the internet every evening.
\item Jogging in the morning helps me to \dots.
\item Playing volleyball is a good way to \dots\ at school.
\item Are you \dots\ storytelling? Our club needs new members!
\end{enumerate}
\tcblower
Solution détaillée :
\begin{enumerate}
\item \eng{fond of} (ou \eng{keen on}) : les deux expressions sont possibles ; le verbe \eng{singing} est bien en -ing après la préposition.
\item \eng{browsing} : \eng{browsing the internet} = naviguer sur Internet.
\item \eng{keep fit} : rester en forme ; après \eng{helps me to}, on utilise la base verbale.
\item \eng{make friends} : se faire des amis ; même structure infinitive.
\item \eng{keen on} (ou \eng{fond of}) : dans une question, le sujet \eng{you} suit l'auxiliaire \eng{are}.
\end{enumerate}
\end{exoresolu}

\begin{attention}[Faux amis et pièges lexicaux]
\begin{itemize}
\item \eng{actually} ne veut pas dire « actuellement » mais « en fait, en réalité ». Pour « actuellement », dites \eng{currently} ou \eng{at the moment} : \eng{I currently play basketball.} « Je joue actuellement au basket. »
\item \eng{a novel} = un roman (littérature), pas « une nouvelle ». Une « nouvelle » se dit \eng{a short story}.
\item \eng{sensible} = raisonnable ; « sensible » se dit \eng{sensitive}.
\item On dit \eng{to watch films} (regarder des films), pas « to look at films » ; mais \eng{to listen to music} (avec la préposition \eng{to}, obligatoire).
\item \eng{free time} et \eng{spare time} = le temps libre ; évitez la traduction mot à mot « free hours ».
\end{itemize}
\end{attention}

\begin{exercice}[À vous de jouer]
\begin{enumerate}
\item Traduisez en anglais : a) J'adore nager. b) Mon frère est passionné de jeux vidéo. c) La randonnée nous aide à réduire le stress. d) À mon avis, le théâtre développe la confiance en soi.
\item Classez ces activités dans les quatre catégories de la carte mentale : \eng{fishing, drama, ludo, watching films, cycling, storytelling, volleyball, listening to music}.
\item Rédigez trois phrases sur vos propres loisirs avec trois expressions de goût différentes.
\end{enumerate}
\end{exercice}

\begin{corrige}[Exercice — Vocabulaire des activités récréatives]
\begin{enumerate}
\item a) \eng{I am fond of swimming.} / \eng{I love swimming.} b) \eng{My brother is keen on playing video games.} c) \eng{Hiking helps us to reduce stress.} d) \eng{In my opinion, drama develops self-confidence.}
\item Outdoor: \eng{fishing, cycling}. Cultural: \eng{drama, storytelling}. Sports and games: \eng{ludo, volleyball}. Modern hobbies: \eng{watching films, listening to music}.
\item Réponses personnelles ; vérifier que chaque expression de goût est bien suivie du verbe en -ing.
\end{enumerate}
\end{corrige}

\begin{aretenir}[L'essentiel de la section]
\begin{itemize}
\item Quatre familles d'activités : sports and games, cultural activities, outdoor activities, modern hobbies.
\item Goûts : \eng{enjoy, be fond of, be keen on, be interested in, prefer} + \cle{-ing} ; \eng{love, like} + \cle{-ing} \textbf{ou} \eng{to} + base verbale.
\item Opinions : \eng{in my opinion, I think that, I believe that, as far as I am concerned}.
\item Bénéfices : \eng{to relax, to keep fit, to make friends, to reduce stress, to develop skills}.
\item Jamais « I am agree » ni « make sport » : \eng{I agree}, \eng{do sport} ou \eng{play a sport}.
\end{itemize}
\end{aretenir}

\coursec{Planifier et rédiger sa composition}

Dans la section précédente, nous avons enrichi notre vocabulaire des activités récréatives : \eng{playing football}, \eng{listening to music}, \eng{reading novels}, \eng{dancing}, \eng{visiting friends}\ldots\ Nous savons maintenant \emph{de quoi parler}. Mais au baccalauréat comme à l'examen de fin d'année, avoir des idées ne suffit pas : il faut savoir les \cle{organiser}, les \cle{rédiger} et les \cle{présenter} dans le temps imparti. C'est précisément l'objet de cette section : transformer un sujet d'examen en une composition claire, bien structurée et sans fautes évitables.

\courssub{Pourquoi planifier avant de rédiger ?}

Observez ces deux débuts de copies, écrits par deux élèves de Première à Yaoundé sur le même sujet : \eng{Write a composition about the recreational activities you take part in.}

\begin{exemplebox}[Deux copies, deux méthodes]
\textbf{Copie A (sans plan) :} \eng{I like football. Football is good. My brother also likes music. Last year I went to the stadium. Reading is also nice. My friend John plays basketball\ldots} --- Les idées sautent d'un sujet à l'autre, sans ordre ni lien.

\textbf{Copie B (avec plan) :} \eng{Everyone needs free time to relax after work or school. In this composition, I will describe the recreational activities I enjoy most. First, I practise sports\ldots} --- Le lecteur sait immédiatement où va l'élève.

« Copie A (sans plan) : J'aime le football. Le football, c'est bien. Mon frère aime aussi la musique. L'année dernière, je suis allé au stade\ldots » / « Copie B (avec plan) : Chacun a besoin de temps libre pour se détendre après le travail ou l'école. Dans cette composition, je vais décrire les activités récréatives que je préfère. D'abord, je pratique un sport\ldots »
\end{exemplebox}

La différence ne tient pas au niveau d'anglais des deux élèves, mais à leur \cle{méthode}. L'élève B a passé cinq minutes à planifier ; l'élève A a commencé à écrire dès la lecture du sujet. Résultat : la copie B est plus courte, mais infiniment mieux notée.

\begin{propriete}[Longueur et qualité d'une composition de Première]
Une composition de classe de Première fait environ \cle{150 à 200 mots} (soit 3 à 4 paragraphes). Ce qui compte pour le correcteur, c'est d'abord la \cle{qualité} et l'\cle{organisation} : une introduction claire, des paragraphes bien délimités, des connecteurs logiques, des temps verbaux corrects. Une copie de 300 mots désorganisée obtiendra une moins bonne note qu'une copie de 160 mots bien construite.
\end{propriete}

\courssub{Les six étapes de la rédaction}

Rédiger une composition n'est pas un acte unique : c'est un \cle{processus} en six étapes. La frise suivante résume le parcours complet, du sujet à la copie finale.

\begin{popfigure}
\begin{center}
\begin{tikzpicture}[
  etape/.style={rectangle, rounded corners=3pt, draw=popblue, fill=popblueL, align=center, font=\small, minimum height=1.1cm, inner sep=4pt},
  fleche/.style={-{Stealth[length=2.5mm]}, thick, poporange}
]
\node[etape, align=center] (e1) at (0,0){1. Read the\\topic};
\node[etape, right=0.35cm of e1, align=center] (e2){2. Brain\\storm};
\node[etape, right=0.35cm of e2, align=center] (e3){3. Select\\ideas};
\node[etape, right=0.35cm of e3, align=center] (e4){4. Make\\a plan};
\node[etape, right=0.35cm of e4] (e5) {5. Draft};
\node[etape, right=0.35cm of e5] (e6) {6. Revise};
\node[etape, right=0.35cm of e6, draw=popgreen, fill=popgreenL, align=center] (e7){7. Copy\\neatly};
\draw[fleche] (e1) -- (e2);
\draw[fleche] (e2) -- (e3);
\draw[fleche] (e3) -- (e4);
\draw[fleche] (e4) -- (e5);
\draw[fleche] (e5) -- (e6);
\draw[fleche] (e6) -- (e7);
\end{tikzpicture}
\end{center}
\legende{Les étapes de la rédaction : lire le sujet, brainstormer, sélectionner, planifier, rédiger le brouillon, relire, recopier au propre.}
\end{popfigure}

Détaillons chaque étape avec le sujet d'examen suivant, que nous traiterons de bout en bout :

\begin{textebox}[Sujet d'examen]
\textbf{Topic:} Write a composition of about 150 to 200 words on the recreational activities you take part in during your free time. Say what they are, when and where you practise them, and why you enjoy them.

« Rédigez une composition d'environ 150 à 200 mots sur les activités récréatives auxquelles vous participez pendant votre temps libre. Dites quelles elles sont, quand et où vous les pratiquez, et pourquoi vous les aimez. »
\end{textebox}

\textbf{Étape 1 : lire attentivement le sujet et souligner les mots-clés.}\par
Lisez le sujet \emph{deux fois}. La première fois pour comprendre le sens général, la deuxième fois pour souligner les \cle{mots-clés} (\eng{key words}), c'est-à-dire les mots qui définissent la tâche. Dans notre sujet, les mots-clés sont :

\begin{itemize}
\item \eng{recreational activities} → le thème imposé (les loisirs, pas les études ni la famille) ;
\item \eng{you take part in} → il s'agit de \emph{vos} activités : il faudra écrire à la première personne (\eng{I}, \eng{my}, \eng{we}) ;
\item \eng{what / when / where / why} → quatre questions auxquelles la composition doit obligatoirement répondre ;
\item \eng{150 to 200 words} → la longueur attendue.
\end{itemize}

\begin{attention}[Rester dans le sujet !]
L'erreur la plus grave à l'examen est le \cle{hors-sujet} (\eng{going off the topic}). Si le sujet porte sur \emph{vos} activités récréatives, ne racontez pas celles de votre voisin, de votre frère ou de votre équipe préférée : utilisez \eng{I} et \eng{my}. De même, si le sujet demande \emph{deux} activités, n'en décrivez pas cinq superficiellement. Relisez le sujet après chaque paragraphe pour vérifier que vous répondez bien à la question posée.
\end{attention}

\textbf{Étape 2 : brainstormer --- noter toutes les idées.}\par
Le \cle{brainstorming} (remue-méninges) consiste à noter rapidement, au brouillon et sans trier, tout le vocabulaire et toutes les idées liés au sujet. Écrivez en anglais dès que possible ; si un mot anglais vous manque, notez-le en français et cherchez un équivalent simple que vous maîtrisez.

\begin{exemplebox}[Brainstorming sur notre sujet]
\eng{football --- Saturday mornings --- school field --- with classmates --- keep fit --- reading novels --- in the evening --- my room --- adventure stories --- relax --- listening to music --- Makossa and Afrobeats --- Sunday afternoons --- dancing --- church choir --- helping my mother --- tired after school --- enjoyable --- friends}

« football --- samedi matin --- terrain de l'école --- avec les camarades --- rester en forme --- lire des romans --- le soir --- ma chambre --- histoires d'aventure --- se détendre --- écouter de la musique --- Makossa et Afrobeats --- dimanche après-midi --- danser --- chœur de l'église --- aider ma mère --- fatigué après l'école --- agréable --- amis »
\end{exemplebox}

\textbf{Étape 3 : sélectionner 2 à 3 idées principales et les classer.}\par
Vous ne pouvez pas tout dire en 200 mots. Choisissez \cle{deux ou trois idées principales} --- celles pour lesquelles vous avez le plus de vocabulaire et de détails --- et barrez le reste. Classez ensuite ces idées dans un ordre logique : du plus général au plus précis, du plus important au moins important, ou par catégories (activités sportives / activités calmes).

Pour notre sujet, nous retenons : (1) le \eng{football} (activité sportive, dehors, le samedi) et (2) la \eng{lecture} et la \eng{musique} (activités calmes, à la maison, le soir). Nous écartons le chœur de l'église et l'aide à la mère, qui allongeraient la copie sans la renforcer.

\textbf{Étape 4 : rédiger un plan.}\par
Le plan est le squelette de votre composition. Il comporte toujours trois parties : une \cle{introduction} (une phrase générale + l'annonce du sujet), un \cle{développement} (un paragraphe par idée principale) et une \cle{conclusion} (une phrase de bilan ou d'opinion).

\begin{exemplebox}[Plan type de notre composition]
\textbf{Introduction :} \eng{Everyone needs free time.} + annonce : \eng{I will describe my favourite recreational activities.}

« Chacun a besoin de temps libre. » + « Je vais décrire mes activités récréatives préférées. »

\textbf{Paragraphe 1 --- sports :} \eng{football --- Saturday mornings --- school field --- with classmates --- keep fit}

« le football --- le samedi matin --- le terrain de l'école --- avec les camarades --- rester en forme »

\textbf{Paragraphe 2 --- quiet activities :} \eng{reading novels and listening to music --- in the evening --- my room --- relax}

« lire des romans et écouter de la musique --- le soir --- ma chambre --- me détendre »

\textbf{Conclusion :} \eng{Recreational activities make life enjoyable.}

« Les activités récréatives rendent la vie agréable. »
\end{exemplebox}

\textbf{Étape 5 : rédiger le brouillon en suivant le plan.}\par
Rédigez maintenant votre \cle{brouillon} (\eng{draft}) en suivant le plan à la lettre. Quelques règles d'or :

\begin{itemize}
\item une idée = un paragraphe ; commencez chaque paragraphe par un connecteur (\eng{First}, \eng{Moreover}, \eng{Finally}) ;
\item utilisez des temps cohérents : le \emph{present simple} pour les habitudes (\eng{I play football every Saturday.}), le \emph{present continuous} seulement pour une action en cours ;
\item préférez des phrases courtes et correctes à des phrases longues et risquées ;
\item si un mot vous manque, contournez la difficulté : au lieu de \eng{refreshing}, écrivez \eng{it makes me feel good}.
\end{itemize}

\textbf{Étape 6 : rédiger au propre.}\par
Après relecture (voir section suivante), recopiez votre texte \cle{au propre} : écriture lisible, alinéas visibles pour chaque paragraphe, pas de ratures. Une copie propre donne immédiatement une bonne impression au correcteur.

\begin{methode}[Réussir sa composition au BAC en cinq réflexes]
\begin{enumerate}
\item \textbf{Analyser le sujet} : lisez-le deux fois et soulignez les mots-clés (thème, personne, questions à traiter, nombre de mots).
\item \textbf{Brainstormer} : notez au brouillon toutes les idées et tout le vocabulaire qui vous viennent, sans trier.
\item \textbf{Planifier} : sélectionnez 2 ou 3 idées, classez-les et rédigez un plan en trois parties (introduction, développement, conclusion). \textbf{Ne sautez jamais cette étape !}
\item \textbf{Rédiger} : écrivez le brouillon en suivant le plan, avec des connecteurs et des phrases simples et correctes.
\item \textbf{Relire} : vérifiez les temps, les accords, l'orthographe et la ponctuation, puis recopiez au propre.
\end{enumerate}
\end{methode}

\courssub{Gérer son temps à l'examen}

À l'examen, vous disposez en général d'environ \cle{45 minutes} pour la composition. Voici comment les répartir :

\begin{center}
\begin{tabularx}{\linewidth}{|l|X|X|}
\hline
\rowcolor{popblueL} \textbf{Temps} & \textbf{Étape} & \textbf{Ce qu'il faut faire} \\
\hline
5 min & Plan & Lire le sujet, souligner les mots-clés, brainstormer, rédiger le plan. \\
\hline
30 min & Rédaction & Écrire le brouillon en suivant le plan, paragraphe par paragraphe. \\
\hline
10 min & Relecture & Corriger les fautes (temps, accords, orthographe), recopier au propre. \\
\hline
\end{tabularx}
\end{center}

Beaucoup d'élèves francophones consacrent 40 minutes à la rédaction et zéro à la relecture : c'est une erreur. Dix minutes de relecture permettent de corriger facilement cinq ou six fautes --- et de gagner des points précieux.

\begin{center}
\begin{tabularx}{\linewidth}{|X|X|X|}
\hline
\rowcolor{popblueL} \textbf{Français} & \textbf{English} & \textbf{Remarque} \\
\hline
le sujet (d'une rédaction) & the topic & Faux ami : \eng{a subject} désigne aussi une matière scolaire. \\
\hline
le brouillon & the draft & Prononciation : « dra-ft ». \\
\hline
le plan & the outline / the plan & \eng{Outline} est plus soutenu. \\
\hline
un paragraphe & a paragraph & Attention à l'orthographe : \eng{par-a-graph}. \\
\hline
recopier au propre & to copy neatly / to write a fair copy & \eng{Neatly} = proprement, soigneusement. \\
\hline
relire & to revise / to proofread & \eng{Revise} = relire et corriger ; \eng{proofread} = relire pour les fautes. \\
\hline
un mot-clé & a key word & Deux mots en anglais. \\
\hline
hors sujet & off the topic & \eng{Your composition is off the topic.} \\
\hline
\end{tabularx}
\end{center}

\begin{exoresolu}[Construire un plan à partir d'un brainstorming]
\textbf{Énoncé.} Voici le brainstorming d'un élève de Bafoussam sur le sujet : \eng{Write about the recreational activities you take part in.} --- \eng{watching football on TV --- playing volleyball --- Tuesday and Friday --- school playground --- my friends --- reading comics --- before sleeping --- relaxing --- cooking --- my uncle's farm}. Sélectionnez deux idées principales, classez-les et rédigez un plan complet (introduction, deux paragraphes, conclusion).

\tcblower
\textbf{Solution détaillée.}

\textbf{1. Sélection.} Les idées les plus riches sont \eng{playing volleyball} (avec jour, lieu, compagnons) et \eng{reading comics} (avec moment et effet). Nous écartons \eng{watching football on TV}, \eng{cooking} et \eng{my uncle's farm}, trop pauvres ou hors du cadre « activités régulières ».

\textbf{2. Classement.} Ordre logique : d'abord l'activité sportive (dehors, en groupe), puis l'activité calme (à la maison, seul).

\textbf{3. Plan rédigé :}

\textbf{Introduction :} \eng{Free time is precious for every student. I take part in two main recreational activities.}

« Le temps libre est précieux pour chaque élève. Je participe à deux activités récréatives principales. »

\textbf{Paragraphe 1 :} \eng{First, I play volleyball with my friends on the school playground every Tuesday and Friday. It helps me keep fit.}

« D'abord, je joue au volley avec mes amis sur le terrain de l'école chaque mardi et vendredi. Cela m'aide à rester en forme. »

\textbf{Paragraphe 2 :} \eng{Moreover, I enjoy reading comics in my room before sleeping. It helps me relax.}

« De plus, j'aime lire des bandes dessinées dans ma chambre avant de dormir. Cela m'aide à me détendre. »

\textbf{Conclusion :} \eng{These activities make my life balanced and enjoyable.}

« Ces activités rendent ma vie équilibrée et agréable. »
\end{exoresolu}

\begin{exercice}[À vous de planifier]
Sujet : \eng{Write a composition (150--200 words) about the recreational activities you take part in during the holidays.}
\begin{enumerate}
\item Soulignez les mots-clés du sujet.
\item Faites un brainstorming d'au moins huit idées ou mots de vocabulaire.
\item Sélectionnez deux idées principales et rédigez un plan en trois parties.
\item Rédigez le premier paragraphe du développement (4 à 5 phrases).
\end{enumerate}
\end{exercice}

\begin{corrige}[Exercice --- éléments de correction]
\textbf{1. Mots-clés à souligner :} \eng{recreational activities}, \eng{you take part in}, \eng{during the holidays}, \eng{150--200 words}.

\textbf{2. Brainstorming possible :} \eng{swimming --- river --- cousins --- village --- playing football --- every afternoon --- visiting grandparents --- helping on the farm --- reading --- listening to the radio --- dancing at parties --- resting}.

\textbf{3. Plan possible.} Introduction : \eng{Holidays give me time for my favourite activities.} --- ¶1 : \eng{playing football every afternoon with my cousins in the village} --- ¶2 : \eng{swimming in the river and listening to the radio in the evening} --- Conclusion : \eng{Holiday activities make me happy and healthy.}

\textbf{4. Paragraphe possible :} \eng{First, I play football with my cousins every afternoon. We meet on the field near my grandparents' house around four o'clock. We form two teams and play until sunset. This activity helps me keep fit and spend time with my family.}

« D'abord, je joue au football avec mes cousins chaque après-midi. Nous nous retrouvons sur le terrain près de la maison de mes grands-parents vers seize heures. Nous formons deux équipes et jouons jusqu'au coucher du soleil. Cette activité m'aide à rester en forme et à passer du temps avec ma famille. »
\end{corrige}

\begin{attention}[Les trois pièges de la rédaction pour les francophones]
\begin{itemize}
\item \textbf{Traduire mot à mot depuis le français.} \eng{I make sport} est un calque de « je fais du sport » ; dites \eng{I play sports} ou \eng{I do sport}. De même, « je m'amuse » se dit \eng{I have fun} ou \eng{I enjoy myself}, jamais \eng{I amuse myself}.
\item \textbf{Oublier le sujet ou la marque de la 3\textsuperscript{e} personne.} En anglais, le sujet est obligatoire : \eng{I play}, pas \eng{Play}. Et au présent simple : \eng{My brother plays}, avec un \eng{-s}.
\item \textbf{Négliger la présentation.} Pas d'alinéa, écriture illisible, ratures : le correcteur est découragé avant même de lire. Sautez une ligne entre les paragraphes ou faites un alinéa visible.
\end{itemize}
\end{attention}

\begin{aretenir}[L'essentiel de la section]
Rédiger une composition, c'est suivre un processus : \eng{Read the topic → Brainstorm → Select ideas → Make a plan → Draft → Revise → Copy neatly}. À l'examen : 5 minutes pour le plan, 30 minutes pour la rédaction, 10 minutes pour la relecture. Une composition de Première fait 150 à 200 mots, en trois parties (introduction, développement, conclusion), et l'organisation compte plus que la longueur. Ne sautez jamais l'étape du plan !
\end{aretenir}

\coursec{Relire, corriger et évaluer}

Dans la section précédente, vous avez appris à planifier puis à rédiger votre composition sur le thème \eng{taking part in recreational activities}. Vous avez maintenant un brouillon (\eng{a draft}) entre les mains. Mais attention : un bon écrivain ne rend jamais sa première version. En anglais comme en français, la différence entre une copie moyenne et une excellente copie se joue à la \cle{relecture} (\eng{proofreading}). Cette dernière étape vous apprend à vérifier votre texte, à traquer vos erreurs et à évaluer votre production avec la grille utilisée par les correcteurs.

\courssub{Pourquoi relire sa composition ?}

Observez ces deux phrases écrites par deux élèves de Première à Douala :

\eng{Last weekend, I play football with my friends. We wins the match and we was very happy.}

\eng{Last weekend, I played football with my friends. We won the match and we were very happy.}

Les deux phrases racontent exactement la même histoire. Pourtant, la première contient trois erreurs de temps et d'accord (\eng{play} au lieu de \eng{played}, \eng{wins} au lieu de \eng{won}, \eng{was} au lieu de \eng{were}), la seconde aucune. Au BAC, le correcteur lit des centaines de copies : un texte propre, correct et bien présenté gagne immédiatement sa confiance. La relecture n'est pas une perte de temps, c'est un \cle{gain de points}.

\begin{methode}[Relire efficacement : la méthode des trois lectures]
Ne relisez jamais votre composition « en général » : vous ne verrez rien. Relisez-la \textbf{trois fois}, avec un objectif différent à chaque lecture.
\begin{enumerate}
\item \textbf{Première lecture : le sens et l'organisation.} Lisez votre texte comme un lecteur qui le découvre. L'introduction présente-t-elle le sujet ? Chaque paragraphe développe-t-il une seule idée ? Y a-t-il des connecteurs (\eng{first, moreover, however, finally}) ? La conclusion répond-elle à la question posée ?
\item \textbf{Deuxième lecture : la grammaire.} Vérifiez chaque verbe : le temps est-il cohérent ? La troisième personne du singulier a-t-elle son \eng{-s} (\eng{he plays}) ? Les auxiliaires sont-ils corrects (\eng{I have played}, pas \eng{I have play}) ?
\item \textbf{Troisième lecture : l'orthographe, la ponctuation et les majuscules.} Traquez les fautes de frappe, les virgules oubliées, les majuscules manquantes (début de phrase, \eng{I}, jours, mois, nationalités : \eng{Monday, July, Cameroonian}).
\end{enumerate}
Astuce : lisez votre texte \textbf{à voix basse}, ou du dernier mot vers le premier pour l'orthographe : votre œil ne « devinera » plus les mots, il les verra vraiment.
\end{methode}

\courssub{Les erreurs à traquer en priorité}

\courssub{L'accord sujet-verbe et le -s de la 3e personne}

C'est l'erreur numéro un des élèves francophones, car en français la terminaison s'entend peu. En anglais, au présent simple, la troisième personne du singulier prend obligatoirement un \eng{-s} (ou \eng{-es}).

\begin{propriete}[Le -s de la 3e personne du singulier au present simple]
\textbf{Forme :} sujet (he / she / it / nom singulier) + verbe + \textbf{s}. Aux formes interrogative et négative, le \eng{-s} passe sur l'auxiliaire \eng{does} et le verbe revient à la base verbale.
\begin{itemize}
\item \eng{She plays basketball every Saturday.} « Elle joue au basket tous les samedis. »
\item \eng{My brother watches football matches on TV.} « Mon frère regarde les matchs de football à la télé. » (verbe en \eng{-ch} : on ajoute \eng{-es})
\item \eng{Does he enjoy swimming? — No, he does not (doesn't).} « Aime-t-il nager ? — Non. »
\end{itemize}
\textbf{Exception :} les auxiliaires modaux (\eng{can, must, should}) ne prennent jamais de \eng{-s} : \eng{she can dance}, jamais \eng{she cans dance}.
\end{propriete}

\begin{exemplebox}[Phrase fautive et phrase corrigée]
\begin{itemize}
\item Faute : \eng{He play football every day.}
\item Correction : \eng{He plays football every day.} « Il joue au football tous les jours. »
\item Faute : \eng{She don't like dancing.}
\item Correction : \eng{She doesn't like dancing.} « Elle n'aime pas danser. »
\end{itemize}
\end{exemplebox}

\courssub{La cohérence des temps verbaux}

Dans une composition sur un souvenir récréatif, le piège est de mélanger passé et présent sans raison. Règle simple : si vous racontez un événement terminé (\eng{last weekend, yesterday, two years ago}), tous les verbes du récit sont au passé ; si vous décrivez vos habitudes (\eng{every weekend, usually}), utilisez le présent simple ; si vous parlez d'une expérience ou d'une durée qui continue (\eng{for two years, since 2022}), utilisez le present perfect.

\begin{center}
\begin{tabularx}{\linewidth}{|X|X|X|}
\hline
\rowcolor{popblueL} Situation & Temps & Exemple \\
\hline
Habitude, goût général & Present simple & \eng{I play football every weekend.} « Je joue au foot chaque week-end. » \\
\hline
Événement passé daté & Past simple & \eng{We won the match last Saturday.} « Nous avons gagné le match samedi dernier. » \\
\hline
Durée commencée dans le passé, encore vraie & Present perfect + for/since & \eng{I have played football for two years.} « Je joue au foot depuis deux ans. » \\
\hline
\end{tabularx}
\end{center}

\begin{attention}[Erreurs fréquentes des francophones]
\begin{itemize}
\item \eng{I play football since two years.} FAUX. En français on dit « je joue depuis deux ans » au présent ; en anglais, une durée qui continue exige le present perfect : \eng{I have played football for two years.} (« depuis » + durée = \eng{for} ; « depuis » + point de départ = \eng{since} : \eng{since 2022}).
\item \eng{I like the football.} FAUX pour parler du sport en général. En anglais, les sports et activités en général se disent \textbf{sans article} : \eng{I like football, basketball and dancing.} « J'aime le football, le basket et la danse. »
\item \eng{I am 16 years.} FAUX. L'anglais exige \eng{old} : \eng{I am 16 years old.} « J'ai 16 ans. » (On ne dit jamais \eng{I have 16 years} : l'âge se dit avec \eng{to be}, pas \eng{to have}.)
\item \eng{I am agree.} FAUX. « Être d'accord » = \eng{I agree} (verbe simple, sans \eng{be}).
\end{itemize}
\end{attention}

\courssub{Vérifier les connecteurs et les paragraphes}

Lors de votre première lecture, posez-vous trois questions : mon texte a-t-il bien trois parties (introduction, développement en deux ou trois paragraphes, conclusion) ? Chaque paragraphe commence-t-il par une phrase d'idée (\eng{topic sentence}) ? Ai-je utilisé au moins quatre connecteurs variés ?

\begin{center}
\begin{tabularx}{\linewidth}{|X|X|}
\hline
\rowcolor{popblueL} Fonction & Connecteurs utiles \\
\hline
Ordonner & \eng{first, then, after that, finally} \\
\hline
Ajouter & \eng{moreover, in addition, also, besides} \\
\hline
Opposer & \eng{however, but, although, on the other hand} \\
\hline
Expliquer & \eng{because, since, that is why, so} \\
\hline
Conclure & \eng{to conclude, in conclusion, all in all} \\
\hline
\end{tabularx}
\end{center}

\courssub{Éviter les répétitions : synonymes et pronoms}

Un texte qui répète \eng{football} huit fois est pauvre. Deux solutions :
\begin{itemize}
\item \textbf{Les pronoms} : \eng{My friends and I play football. We love it because it is exciting.} (« Mes amis et moi jouons au foot. Nous adorons cela parce que c'est passionnant. ») — \eng{we} remplace \eng{my friends and I}, \eng{it} remplace \eng{football}.
\item \textbf{Les synonymes} : \eng{fun, enjoyable, exciting, entertaining} (amusant) ; \eng{activity, hobby, pastime, leisure activity} (activité de loisir) ; \eng{tiring, exhausting} (fatigant).
\end{itemize}

\begin{exoresolu}[Corriger un paragraphe d'élève]
Énoncé — Voici un extrait de la copie de Marlyse. Trouvez et corrigez les cinq erreurs : \eng{Last month, I go to a dance competition in Yaoundé with my school. We dances traditional dances and we was very proud. I am 17 years and I dance since three years. Dancing is my favourite hobby because the dancing makes me happy.}
\tcblower
Solution détaillée :
\begin{enumerate}
\item \eng{I go} $\rightarrow$ \eng{I went} : événement passé daté (\eng{last month}), donc past simple.
\item \eng{We dances} $\rightarrow$ \eng{we danced} : accord (pas de \eng{-s} au pluriel) et cohérence du passé.
\item \eng{we was} $\rightarrow$ \eng{we were} : \eng{was} est réservé à I/he/she/it.
\item \eng{I am 17 years} $\rightarrow$ \eng{I am 17 years old} : \eng{old} obligatoire.
\item \eng{I dance since three years} $\rightarrow$ \eng{I have danced for three years} : durée qui continue = present perfect + \eng{for}. Et \eng{the dancing makes me happy} $\rightarrow$ \eng{dancing makes me happy} : pas d'article devant une activité en général.
\end{enumerate}
Paragraphe corrigé : \eng{Last month, I went to a dance competition in Yaoundé with my school. We danced traditional dances and we were very proud. I am 17 years old and I have danced for three years. Dancing is my favourite hobby because it makes me happy.} « Le mois dernier, je suis allée à un concours de danse à Yaoundé avec mon école. Nous avons dansé des danses traditionnelles et nous étions très fiers. J'ai 17 ans et je danse depuis trois ans. La danse est mon passe-temps préféré parce qu'elle me rend heureuse. » (Notez aussi le pronom \eng{it} qui évite de répéter \eng{dancing}.)
\end{exoresolu}

\courssub{La grille d'évaluation}

Au BAC comme en classe, votre composition est notée selon quatre critères. Connaître cette grille, c'est savoir exactement ce que le correcteur attend.

\begin{popfigure}
\begin{tikzpicture}[
  cell/.style={draw=popline, fill=popcream, minimum width=3.4cm, minimum height=1.15cm, align=center, font=\small, text=popdark},
  head/.style={draw=popblue, fill=popblue, minimum width=3.4cm, minimum height=0.9cm, align=center, font=\small\bfseries, text=white},
  crit/.style={draw=poporange, fill=poporangeL, minimum width=3.4cm, minimum height=1.15cm, align=center, font=\small\bfseries, text=popdark}
]
\node[head] at (0,0) {Criterion};
\node[head] at (3.5,0) {Excellent};
\node[head] at (7.0,0) {Good};
\node[head] at (10.5,0) {Needs improvement};

\node[crit] at (0,-1.1) {Content};
\node[cell, align=center] at (3.5,-1.1){Relevant ideas,\\fully answers\\the topic};
\node[cell, align=center] at (7.0,-1.1){Mostly relevant\\ideas};
\node[cell, align=center] at (10.5,-1.1){Off topic or\\very few ideas};

\node[crit] at (0,-2.3) {Organisation};
\node[cell, align=center] at (3.5,-2.3){Clear intro, body,\\conclusion; varied\\linking words};
\node[cell, align=center] at (7.0,-2.3){Visible structure;\\some linking\\words};
\node[cell, align=center] at (10.5,-2.3){No paragraphs;\\no linking words};

\node[crit, align=center] at (0,-3.5){Grammar \&\\Spelling};
\node[cell, align=center] at (3.5,-3.5){Almost no\\mistakes; correct\\tenses};
\node[cell, align=center] at (7.0,-3.5){Some mistakes\\but meaning is\\clear};
\node[cell, align=center] at (10.5,-3.5){Many mistakes;\\meaning often\\unclear};

\node[crit] at (0,-4.7) {Vocabulary};
\node[cell, align=center] at (3.5,-4.7){Rich, precise,\\varied words};
\node[cell, align=center] at (7.0,-4.7){Correct but\\simple words};
\node[cell, align=center] at (10.5,-4.7){Very limited;\\many repetitions};
\end{tikzpicture}
\legende{Grille d'évaluation d'une composition : quatre critères, trois niveaux de performance.}
\end{popfigure}

\begin{definition}[Les quatre critères en détail]
\begin{itemize}
\item \textbf{Content} (contenu) : vos idées sont-elles pertinentes ? Répondez-vous vraiment au sujet posé ?
\item \textbf{Organisation} : structure en trois parties, paragraphes nets, connecteurs variés.
\item \textbf{Grammar \& Spelling} (grammaire et orthographe) : accords, temps, orthographe, ponctuation, majuscules.
\item \textbf{Vocabulary} (vocabulaire) : richesse, variété, exactitude des mots choisis.
\end{itemize}
\end{definition}

\courssub{Auto-évaluation et évaluation par les pairs}

Avant de rendre votre copie, évaluez-vous vous-même avec la grille : c'est l'\cle{auto-évaluation} (\eng{self-assessment}). Puis échangez votre brouillon avec un voisin : c'est l'\cle{évaluation par les pairs} (\eng{peer assessment}). Un regard neuf repère des erreurs que vous ne voyez plus. Pour chaque critère, entourez \eng{Excellent}, \eng{Good} ou \eng{Needs improvement}, puis écrivez un conseil précis : \eng{Add two linking words in paragraph 2} (« ajoute deux connecteurs au paragraphe 2 »), \eng{Check the -s in “he play”} (« vérifie le -s »).

\begin{experience}[Atelier de correction par les pairs]
Par groupes de deux : échangez vos brouillons. Chaque élève lit la copie de son voisin trois fois (sens, grammaire, orthographe), remplit la grille d'évaluation et écrit deux points forts (\eng{strengths}) et deux conseils (\eng{tips}). Ensuite, chacun réécrit la version finale de sa composition en tenant compte des remarques. L'enseignant circule et arbitre les désaccords.
\end{experience}

\courssub{Modèle commenté : une copie annotée}

\begin{textebox}[A student's composition — annotated model]
\textbf{My favourite recreational activity}\par
[1] \eng{Recreational activities are very important in the life of a teenager. My favourite one is football, and I would like to explain why I love it so much.}\par
[2] \eng{First, football keeps me fit and healthy. I have played football for five years with my friends at the Lycée de Nkolbisson. Every Saturday, we train for two hours, and after that we often play a friendly match. Moreover, football teaches me team spirit: when we play, we must help each other and share the ball.}\par
[3] \eng{However, football is not always easy. Sometimes we lose matches, and last month we were very disappointed when our team lost the final of the school tournament. But these moments also make us stronger.}\par
[4] \eng{To conclude, football is more than a game for me: it is a passion. It makes me happy, healthy and friendly. I think every teenager should practise a recreational activity.}
\end{textebox}

Glossaire : \emph{to keep fit} (v., se maintenir en forme) ; \emph{to train} (v., s'entraîner) ; \emph{team spirit} (n., esprit d'équipe) ; \emph{to share} (v., partager) ; \emph{disappointed} (adj., déçu) ; \emph{tournament} (n., tournoi).

\textbf{Points forts de cette copie :} [1] introduction claire qui annonce le sujet ; [2] connecteurs variés (\eng{first, moreover, after that}), present perfect correct (\eng{I have played for five years}), pas d'article devant \eng{football} ; [3] concession bien marquée par \eng{however} et temps du passé cohérents (\eng{were, lost}) ; [4] conclusion personnelle avec \eng{to conclude}. \textbf{Point à améliorer :} le mot \eng{football} revient souvent ; on pourrait écrire \eng{this sport} ou \eng{it} une ou deux fois.

\begin{savaistu}[Une copie organisée bat une copie longue]
Au BAC, une copie bien organisée avec quelques petites erreurs obtient souvent une meilleure note qu'une copie longue mais désorganisée. Les correcteurs lisent vite : une introduction visible, des paragraphes nets et des connecteurs clairs montrent immédiatement que vous maîtrisez l'écriture anglaise. Mieux vaut 180 mots propres et structurés que 350 mots confus. La qualité avant la quantité !
\end{savaistu}

\begin{exercice}[À vous de corriger]
Voici la conclusion écrite par Étienne. Trouvez les quatre erreurs, corrigez-les, puis notez sa copie avec la grille (critère \eng{Grammar \& Spelling} uniquement) : \eng{To conclude, I love the basketball very much. I play basketball since four years and it make me strong. I am 16 years and I think sport is very important for teenagers.}
\end{exercice}

\begin{corrige}[Exercice : À vous de corriger]
Erreurs et corrections : (1) \eng{the basketball} $\rightarrow$ \eng{basketball} (sport en général, sans article) ; (2) \eng{I play basketball since four years} $\rightarrow$ \eng{I have played basketball for four years} (durée qui continue : present perfect + \eng{for}) ; (3) \eng{it make me strong} $\rightarrow$ \eng{it makes me strong} (\eng{-s} à la 3e personne) ; (4) \eng{I am 16 years} $\rightarrow$ \eng{I am 16 years old}. Texte corrigé : \eng{To conclude, I love basketball very much. I have played basketball for four years and it makes me strong. I am 16 years old and I think sport is very important for teenagers.} « Pour conclure, j'aime beaucoup le basket. J'y joue depuis quatre ans et cela me rend fort. J'ai 16 ans et je pense que le sport est très important pour les adolescents. » Avec quatre erreurs en quatre phrases, le critère \eng{Grammar \& Spelling} se situe entre \eng{Good} et \eng{Needs improvement} : le sens reste clair, mais les erreurs sont fréquentes.
\end{corrige}

\begin{aretenir}[L'essentiel de la section]
\begin{itemize}
\item Relisez toujours votre composition \textbf{trois fois} : sens et organisation, puis grammaire, puis orthographe et ponctuation.
\item Traquez en priorité : le \eng{-s} de la 3e personne (\eng{he plays}), la cohérence des temps, le present perfect avec \eng{for/since}, l'absence d'article devant les sports, \eng{I am 16 years old}.
\item Remplacez les répétitions par des pronoms (\eng{it, we}) ou des synonymes (\eng{hobby, pastime, activity}).
\item Évaluez votre copie avec la grille à quatre critères : \eng{Content, Organisation, Grammar \& Spelling, Vocabulary} — seul ou avec un camarade.
\end{itemize}
\end{aretenir}

\newpage

\coursec{Activité d'intégration}

\coursec{Lesson 12 — Writing: Compositions about taking part in recreational activities}

\courssub{Objectifs de la leçon}

À la fin de cette leçon, tu seras capable de :

\begin{itemize}
\item comprendre les attendus d'une composition sur les activités récréatives (sujet, longueur, ton) ;
\item identifier et appliquer la structure en trois parties : \cle{introduction}, \cle{développement}, \cle{conclusion} ;
\item utiliser avec justesse les \cle{connecteurs logiques} (addition, opposition, cause, illustration, conclusion) ;
\item mobiliser les temps verbaux appropriés : \cle{present simple} (habitudes), \cle{past simple} (anecdotes), \cle{present continuous} (actions en cours) ;
\item employer un lexique varié et précis sur les loisirs (sports, arts, vie sociale, adjectifs de sentiment) ;
\item planifier, rédiger, relire et évaluer une composition de 150 à 200 mots selon les critères du Baccalauréat.
\end{itemize}

\courssub{Découverte : un modèle de composition}

\begin{textebox}[A message from your English Club]
Dear members,

Our school English Club is preparing the next issue of its wall magazine, \emph{Free Time Voices}. This month's topic is: \textbf{“My Favourite Recreational Activities”}.

Every student in Form Five is invited to write a composition of 150 to 200 words. Describe the recreational activities you enjoy (sports, music, games, reading, dancing, visiting friends...), say when and with whom you do them, and explain why they are important to you.

The best compositions will be read aloud during the club meeting and displayed on the classroom wall.

Deadline: Friday, 4 p.m. — Give your composition to the club secretary.

The English Club Committee,
Government High School, Buea
\end{textebox}

\textbf{Traduction du document.} Chers membres, notre club d'anglais prépare le prochain numéro de son magazine mural, \emph{Free Time Voices}. Le sujet du mois : « Mes activités récréatives préférées ». Chaque élève de Première est invité à rédiger une composition de 150 à 200 mots. Décris les activités que tu aimes (sport, musique, jeux, lecture, danse, visites aux amis...), dis quand et avec qui tu les pratiques, et explique pourquoi elles sont importantes pour toi. Les meilleures compositions seront lues à voix haute et affichées au mur de la classe.

\begin{exemplebox}[Modèle de composition — Analyse rapide]
\emph{My Favourite Recreational Activities}

\eng{Everybody needs free time to relax after long hours of study. In my free time, I enjoy several recreational activities that make me happy and keep me healthy.}

\eng{First of all, I love playing football with my friends on the community field every Saturday afternoon. Football is exciting and it teaches me teamwork. Moreover, I am fond of music. In the evenings, I often listen to Makossa and Afrobeats songs, and sometimes I dance with my sisters at home. It is always a wonderful moment.}

\eng{However, I also enjoy quiet activities. For example, I like reading adventure novels before going to sleep because they broaden my mind and improve my English. During the last holidays, I read three novels, and I was very proud of myself.}

\eng{In conclusion, recreational activities are very important in my life. They help me to relax, to stay fit and to share quality time with the people I love.}

\textbf{Observations guidées :}
\begin{enumerate}
\item Combien de paragraphes ? \textbf{5} (Introduction, 3 paragraphes de développement, Conclusion).
\item Quels connecteurs relient les paragraphes ? \eng{First of all, Moreover, However, For example, In conclusion}.
\item Quels temps verbaux dominent ? \textbf{Present simple} (\eng{I love, I often listen, they help}) pour les habitudes ; \textbf{Past simple} (\eng{I read, I was}) pour l'anecdote.
\item Repère les prépositions : \eng{listen \textbf{to}}, \eng{fond \textbf{of}}, \eng{\textbf{at} home}, \eng{\textbf{on} Saturday}, \eng{\textbf{in} the evenings}.
\end{enumerate}
\end{exemplebox}

\courssub{La structure d'une composition}

\begin{methode}[Comment structurer une composition en 3 parties]
\begin{enumerate}
\item \textbf{Introduction (1 paragraphe, 2--3 phrases)} :
      \begin{itemize}
      \item \emph{Accroche} : une phrase générale sur le thème (proverbe, constat, question).
      \item \emph{Annonce du sujet} : \eng{In my free time, I enjoy...} / \eng{My favourite recreational activities are...}
      \item \emph{Transition} : \eng{They help me relax and stay healthy.}
      \end{itemize}
\item \textbf{Développement (2 à 3 paragraphes, 1 idée principale par paragraphe)} :
      \begin{itemize}
      \item Paragraphe 1 : Activité 1 + détails (quand, avec qui, pourquoi). Connecteur d'ouverture : \eng{First of all,} / \eng{To begin with,}.
      \item Paragraphe 2 : Activité 2 + détails. Connecteur : \eng{Moreover,} / \eng{In addition,} / \eng{Besides,}.
      \item Paragraphe 3 (optionnel) : Activité 3 ou nuance (contraste). Connecteur : \eng{However,} / \eng{On the other hand,}.
      \item \textbf{Règle d'or} : Une idée = un paragraphe. Appuie-toi sur \eng{for example,} / \eng{such as} pour illustrer.
      \end{itemize}
\item \textbf{Conclusion (1 paragraphe, 2--3 phrases)} :
      \begin{itemize}
      \item Connecteur de conclusion : \eng{In conclusion,} / \eng{To sum up,} / \eng{Finally,}.
      \item \emph{Bilan} : résume l'importance des loisirs pour toi.
      \item \emph{Opinion personnelle / ouverture} : \eng{I think everyone should...} / \eng{Without them, life would be boring.}
      \end{itemize}
\end{enumerate}
\end{methode}

\begin{propriete}[Plan type « Mes activités récréatives »]
\begin{center}
\begin{tabularx}{\linewidth}{|X|X|}
\hline
\rowcolor{popblueL} \textbf{Partie} & \textbf{Contenu et phrases-types} \\
\hline
\textbf{Introduction} & \eng{Free time is essential for every student. Personally, I enjoy...} \\
\hline
\textbf{Développement 1} & \eng{First of all, I love [activity] with [people] at [time/place]. It is [adjective] because...} \\
\hline
\textbf{Développement 2} & \eng{Moreover, I am keen on [activity]... For instance, last week...} \\
\hline
\textbf{Développement 3} & \eng{However, I also like quiet moments... [activity] allows me to...} \\
\hline
\textbf{Conclusion} & \eng{In conclusion, these activities are vital. They help me [benefits]. I recommend them to everyone.} \\
\hline
\end{tabularx}
\end{center}
\end{propriete}

\courssub{Connecteurs logiques et temps verbaux}

\begin{propriete}[Les connecteurs logiques indispensables]
\begin{center}
\begin{tabularx}{\linewidth}{|l|X|X|}
\hline
\rowcolor{popblueL} \textbf{Fonction} & \textbf{Connecteurs (English)} & \textbf{Équivalent français} \\
\hline
\textbf{Ordre / Addition} & \eng{First of all, To begin with, First, Secondly, Moreover, Furthermore, In addition, Besides, Also} & Tout d'abord, ensuite, de plus, en outre, aussi \\
\hline
\textbf{Illustration} & \eng{For example, For instance, Such as, Like} & Par exemple, comme \\
\hline
\textbf{Cause / Conséquence} & \eng{Because, Since, As, So, Therefore, Consequently} & Parce que, donc, par conséquent \\
\hline
\textbf{Opposition / Contraste} & \eng{However, On the other hand, Although, Even though, But, Yet} & Cependant, bien que, mais \\
\hline
\textbf{Conclusion} & \eng{In conclusion, To sum up, Finally, All in all} & En conclusion, pour résumer, enfin \\
\hline
\end{tabularx}
\end{center}
\textbf{Attention :} \eng{However} et \eng{Therefore} sont souvent suivis d'une virgule en début de phrase ou séparés par un point-virgule ; \eng{Because} et \eng{Although} introduisent une subordonnée (sujet + verbe).
\end{propriete}

\begin{propriete}[Choix des temps verbaux dans une composition sur les loisirs]
\begin{center}
\begin{tabular}{|l|l|l|l|}
\hline
\rowcolor{popblueL} \textbf{Temps} & \textbf{Emploi} & \textbf{Marqueurs de temps} & \textbf{Exemple} \\
\hline
\cle{Present Simple} & Habitudes, goûts, vérités générales & \eng{every day, on Saturdays, usually, often, always} & \eng{I \textbf{play} football every weekend.} \\
\hline
\cle{Present Continuous} & Action en cours, projet proche & \eng{now, at the moment, currently, this week} & \eng{Right now, I \textbf{am reading} a novel.} \\
\hline
\cle{Past Simple} & Anecdote terminée, souvenir daté & \eng{yesterday, last week, two days ago, in 2023, during the holidays} & \eng{Last holiday, I \textbf{learnt} to swim.} \\
\hline
\end{tabular}
\end{center}
\end{propriete}

\begin{attention}[Erreurs fréquentes des francophones — Grammaire]
\begin{itemize}
\item \textbf{Oubli du -s à la 3\textsuperscript{e} personne du singulier (Present Simple)} : \eng{She \textbf{plays}} (pas \eng{play}). \eng{It \textbf{helps}} (pas \eng{help}).
\item \textbf{Prépositions verbes + activités} : \eng{listen \textbf{to} music}, \eng{keen \textbf{on} dancing}, \eng{fond \textbf{of} reading}, \eng{interested \textbf{in} sports}, \eng{good \textbf{at} playing}.
\item \textbf{Prépositions de temps} : \eng{\textbf{at} weekends} (UK), \eng{\textbf{on} Saturday}, \eng{\textbf{in} the evening}, \eng{\textbf{during} the holidays}.
\item \textbf{Faux ami \eng{actually}} : signifie « en fait / en réalité » (pas « actuellement » $\rightarrow$ \eng{at the moment} / \eng{currently}).
\item \textbf{\eng{I am agree} $\rightarrow$ INCORRECT}. Dire : \eng{I \textbf{agree}}.
\item \textbf{Concordance des temps} : Si tu commences au \cle{Past Simple} pour une anecdote, reste au Past Simple. \eng{Last year, I \textbf{went} to the stadium and \textbf{watched} a match.}
\item \textbf{Ordre des mots} : Sujet + Verbe + Complément. \eng{I \textbf{play football}} (pas \eng{I play the football} sauf cas précis).
\end{itemize}
\end{attention}

\courssub{Vocabulaire des activités récréatives}

\begin{aretenir}[Lexique essentiel : loisirs et temps libre]
\begin{center}
\begin{tabularx}{\linewidth}{|X|X|X|}
\hline
\rowcolor{popblueL} \textbf{Catégorie} & \textbf{Mots / Expressions (English)} & \textbf{Traduction / Note} \\
\hline
\textbf{Sports \& Action} & \eng{play football / basketball / handball}, \eng{go swimming / jogging / cycling}, \eng{do gymnastics / yoga}, \eng{work out}, \eng{keep fit}, \eng{teamwork}, \eng{coach}, \eng{playground}, \eng{stadium} & \eng{play} (ballon), \eng{go} + -ing, \eng{do} (sport individuel) \\
\hline
\textbf{Arts \& Culture} & \eng{listen to music}, \eng{play the guitar / piano}, \eng{sing}, \eng{dance}, \eng{draw}, \eng{paint}, \eng{read novels / comics / magazines}, \eng{watch films / series}, \eng{go to the cinema} & \eng{listen \textbf{to}}, \eng{play \textbf{the}} (instrument) \\
\hline
\textbf{Vie sociale \& Détente} & \eng{chat with friends}, \eng{hang out with friends}, \eng{visit relatives}, \eng{have a picnic}, \eng{play video games}, \eng{board games}, \eng{socialise}, \eng{community field}, \eng{youth centre} & \eng{hang out} (familier : traîner) \\
\hline
\textbf{Adjectifs de sentiment} & \eng{relaxing}, \eng{exciting}, \eng{thrilling}, \eng{boring}, \eng{tiring}, \eng{healthy}, \eng{beneficial}, \eng{creative}, \eng{sociable}, \eng{stressful} & \eng{-ing} = cause du sentiment \\
\hline
\textbf{Verbes d'opinion} & \eng{enjoy}, \eng{love}, \eng{like}, \eng{be fond of}, \eng{be keen on}, \eng{prefer}, \eng{hate}, \eng{can't stand} + \textbf{-ing} & Tous suivis du gérondif (\eng{-ing}) \\
\hline
\end{tabularx}
\end{center}
\end{aretenir}

\begin{savaistu}[Loisirs au Cameroun et dans le monde anglophone]
Au Cameroun, le \textbf{football} est roi (Lions Indomptables, championnats locaux). On joue sur le \eng{community field} ou le \eng{playground}. La musique (\textbf{Makossa}, \textbf{Bikutsi}, \textbf{Afrobeats}) accompagne les \eng{dance} improvisées. Les \eng{bendskins} (motos-taxis) mènent aux \eng{buvettes} ou \eng{call-box} pour \eng{chat with friends}. Au Nigeria ou au Ghana, on pratique le \eng{football} et l'\eng{athletics} ; au Royaume-Uni/USA, le \eng{cricket}, \eng{rugby}, \eng{baseball}, \eng{basketball} sont populaires. Le \eng{reading} de romans africains (\eng{Chinua Achebe, Mongo Beti, Chimamanda Ngozi Adichie}) est une activité scolaire valorisée.
\end{savaistu}

\courssub{Planifier et rédiger sa composition : tâche d'intégration}

\begin{experience}[Atelier d'écriture : « My Favourite Recreational Activities »]
\textbf{Situation.} Tu es élève en Première A au Government High School, Buea. Tu participes au concours du club d'anglais (voir message de la Découverte). Tu as 45 minutes en classe.

\textbf{Tâche 1 — Brainstorming (5 min).} Sur ton brouillon, remplis ce tableau en anglais (minimum 3 items par colonne) :
\begin{center}
\begin{tabularx}{\linewidth}{|X|X|X|X|}
\hline
\rowcolor{popblueL} \textbf{Activities} & \textbf{When / Frequency} & \textbf{With whom / Where} & \textbf{Why / Feelings} \\
\hline
\eng{playing football} & \eng{every Saturday afternoon} & \eng{with friends / at the community field} & \eng{exciting, teamwork, keep fit} \\
\hline
\eng{listening to music} & \eng{every evening} & \eng{at home / with sisters} & \eng{relaxing, happy, cultural} \\
\hline
\eng{reading novels} & \eng{before bed / during holidays} & \eng{alone / in my bedroom} & \eng{broaden mind, improve English} \\
\hline
\eng{...} & \eng{...} & \eng{...} & \eng{...} \\
\hline
\end{tabularx}
\end{center}

\textbf{Tâche 2 — Le plan structuré (5 min).} Rédige ton plan en phrases-clés (anglais).
\begin{enumerate}
\item \textbf{Introduction} : Accroche + Sujet. \textit{Ex: \eng{Free time is vital. I enjoy football, music and reading.}}
\item \textbf{Développement 1 (Sport)} : \eng{First of all, football... when/with whom... why (teamwork, health).}
\item \textbf{Développement 2 (Musique/Danse)} : \eng{Moreover, music... Makossa/Afrobeats... dance with sisters... relaxing.}
\item \textbf{Développement 3 (Lecture - contraste)} : \eng{However, I also like quiet activities... reading novels... example (last holidays)... benefits.}
\item \textbf{Conclusion} : \eng{In conclusion, essential for balance... relax, fit, socialise. Advice.}
\end{enumerate}

\textbf{Tâche 3 — Le brouillon (20 min).} Rédige 150--200 mots. Titre : \emph{My Favourite Recreational Activities}.
\begin{itemize}
\item Utilise \textbf{au moins 5 connecteurs différents} (cocher sur la liste de la Propriété).
\item Varie les temps : \cle{Present Simple} (base), \cle{Past Simple} (1 anecdote), \cle{Present Continuous} (si description scène).
\item Soigne les prépositions (\eng{listen to, fond of, at weekends, in the evening}).
\item Compte les mots (viser 170--180).
\end{itemize}

\textbf{Tâche 4 — Relecture guidée (10 min).} Utilise la \cle{Checklist de relecture} ci-dessous. Corrige au stylo vert.
\end{experience}

\begin{aretenir}[Checklist de relecture — \eng{Proofreading Checklist}]
\begin{enumerate}
\item \textbf{Structure} : Titre ? Introduction (accroche+sujet) ? 2--3 paragraphes développement ? Conclusion (opinion) ?
\item \textbf{Connecteurs} : Au moins 5 variés ? Bien placés (début phrase, après point-virgule) ?
\item \textbf{Grammaire} : \eng{-s} 3\textsuperscript{e} pers. sing. (Present Simple) ? \textbf{Past Simple} réguliers/irréguliers corrects (\eng{went, read, was}) ? Prépositions verbes (\eng{listen to, keen on}) ? Prépositions temps (\eng{on, at, in, during}) ?
\item \textbf{Ordre des mots} : Sujet--Verbe--Complément respecté ? Pas de traduction mot à mot ?
\item \textbf{Orthographe / Ponctuation} : Majuscules (I, jours, mois, noms propres, début phrase) ? Points, virgules ? Pas d'abréviations SMS (\eng{u, thx}) ?
\item \textbf{Vocabulaire} : Lexique loisirs varié ? Adjectifs \eng{-ing} ? Pas de répétitions excessives ?
\item \textbf{Nombre de mots} : Entre 150 et 200 ? (Compter : 1 ligne $\approx$ 10 mots).
\end{enumerate}
\end{aretenir}

\textbf{Tâche 5 — Évaluation par les pairs (5 min).} Échange ta copie avec ton binôme. Note sur 20 avec la grille. Écris \textbf{deux conseils précis et bienveillants} en anglais au dos de la copie.
\begin{center}
\begin{tabularx}{\linewidth}{|X|c|}
\hline
\rowcolor{popblueL} \textbf{Critère d'évaluation (Grille officielle type Bac)} & \textbf{Note} \\
\hline
Contenu : sujet traité, idées riches, personnelles, exemples concrets & /5 \\
\hline
Organisation : structure 3 parties, paragraphes, connecteurs variés & /5 \\
\hline
Grammaire : temps justes, accords, ordre mots, prépositions & /5 \\
\hline
Vocabulaire \& Orthographe : lexique loisirs précis, orthographe soignée & /5 \\
\hline
\textbf{Total} & \textbf{/20} \\
\hline
\end{tabularx}
\end{center}
\textit{Exemples de conseils :} \eng{Try to use more link words in the second paragraph.} / \eng{Check the past tense of 'go' (went).} / \eng{Good vocabulary! Next time, add a personal anecdote.}

\textbf{Tâche 6 — Publication.} Les 3 meilleures copies sont lues par leurs auteurs, affichées au mur (\eng{Free Time Voices}) et photographiées pour le groupe WhatsApp de la classe.


\courssub{Relire, corriger et évaluer : outils de référence}

\begin{methode}[La démarche en 5 étapes pour toute composition]
\begin{enumerate}
\item \cle{Brainstormer} : Jeter les idées (mots, expressions) sans phrases complètes. Classer par catégorie.
\item \cle{Planifier} : Choisir les 2--3 meilleures idées. Ordonner logiquement. Écrire le plan en mots-clés anglais.
\item \cle{Rédiger} : Suivre le plan. Écrire des phrases complètes. Ne pas chercher le mot parfait tout de suite (brouillon). Utiliser les connecteurs.
\item \cle{Relire} : Appliquer la Checklist (Grammaire $\rightarrow$ Vocabulaire $\rightarrow$ Structure $\rightarrow$ Mots). Corriger.
\item \cle{Évaluer} : Se noter soi-même ou faire noter par un pair avec la grille. Lire les conseils. Réécrire au propre si nécessaire.
\end{enumerate}
\end{methode}

\courssub{Production modèle commentée}

\begin{exoresolu}[My Favourite Recreational Activities — Modèle final commenté (182 mots)]
Rédige la composition pour le magazine \emph{Free Time Voices} (150--200 mots).

\tcblower

\textbf{My Favourite Recreational Activities}

\eng{Free time is not a luxury; it is a necessity for every student. In my free time, I enjoy several recreational activities that bring me joy and keep me balanced.}

\eng{First of all, I love playing football with my friends on the community field every Saturday afternoon. Football is thrilling and it teaches me teamwork and discipline. Moreover, I am keen on music. In the evenings, I often listen to Makossa and Afrobeats, and sometimes I dance with my sisters in the sitting room. It is always a relaxing and happy moment.}

\eng{However, I also appreciate quiet activities. For instance, I like reading adventure novels before going to sleep because they broaden my horizons and improve my vocabulary. During the last holidays, I read three novels, and I felt very proud of my progress.}

\eng{In conclusion, recreational activities are essential in my life. They help me to relax, to stay fit and to share quality time with my loved ones. I would advise every student to find a hobby they truly enjoy.}

\textbf{Commentaire détaillé des choix linguistiques :}
\begin{itemize}
\item \textbf{Structure} : Accroche proverbe (\eng{Free time is not a luxury...}), annonce claire (\eng{I enjoy several...}), 3 paragraphes développement distincts (sport, musique, lecture), conclusion avec conseil (\eng{I would advise...}).
\item \textbf{Connecteurs (6 variés)} : \eng{First of all} (ordre), \eng{Moreover} (addition), \eng{However} (contraste), \eng{For instance} (illustration), \eng{Because} (cause), \eng{In conclusion} (synthèse).
\item \textbf{Temps verbaux} : \cle{Present Simple} dominant pour habitudes (\eng{I love, I often listen, it teaches, they help}) ; \cle{Past Simple} pour anecdote datée (\eng{I read, I felt}) ; \cle{Present Continuous} absent ici mais possible (\eng{Currently, I am learning the guitar}).
\item \textbf{Grammaire soignée} : \eng{listen \textbf{to}}, \eng{keen \textbf{on}}, \eng{fond \textbf{of}} (non utilisé ici mais connu), \eng{good \textbf{at}}, \eng{\textbf{at} weekends} (implicite \eng{every Saturday}), \eng{\textbf{in} the evenings}, \eng{\textbf{during} the holidays}. Accords 3\textsuperscript{e} pers. : \eng{it teaches, it helps, she dances} (si elle était sujet).
\item \textbf{Vocabulaire riche} : \eng{thrilling, teamwork, discipline, keen on, sitting room, broaden my horizons, vocabulary, essential, quality time, loved ones, advise}.
\end{itemize}
\end{exoresolu}

\courssub{Bilan de la leçon}

Dans cette leçon, tu as maîtrisé le processus complet de l'écriture d'une composition argumentative et descriptive sur un thème familier.

\begin{itemize}
\item \textbf{Savoir} : Structure en 3 parties, connecteurs classés par fonction, temps verbaux (Present Simple, Past Simple), lexique loisirs + prépositions, grille d'évaluation.
\item \textbf{Savoir-faire} : Brainstorming efficace $\rightarrow$ Plan cohérent $\rightarrow$ Rédaction fluide $\rightarrow$ Relecture autonome (checklist) $\rightarrow$ Évaluation par les pairs (grille).
\item \textbf{Attitude} : Rigueur (respect consignes : nombre de mots, titre), autonomie (brouillon, correction), bienveillance (conseils constructifs aux pairs).
\end{itemize}

\cle{Retiens la routine « 5 étapes »} : \textbf{Brainstormer -- Planifier -- Rédiger -- Relire -- Évaluer}. Elle est transférable à tous les sujets de composition (Lettres, Articles, Récits, Essais) du programme de Première et du Baccalauréat.

\begin{attention}[Derniers rappels pour l'examen]
\begin{itemize}
\item \textbf{Lis le sujet 2 fois.} Souligne : thème, rôle, destinataire, nombre de mots, temps imposés.
\item \textbf{Ne commence JAMAIS au propre.} Le brouillon est obligatoire (noté parfois).
\item \textbf{Gère ton temps :} 10 min plan/brouillon, 25 min rédaction, 10 min relecture/copie propre.
\item \textbf{Soigne la présentation :} Titre centré, paragraphes sautés, écriture lisible, marge gauche.
\end{itemize}
\end{attention}

\newpage

\coursec{Méthodes et exercices résolus}

\courssub{Fiches méthodes et exercices résolus}

\begin{methode}[Planifier et rédiger une composition sur les activités récréatives]
Pour réussir une composition de type \eng{Write a composition of about 150 words on how you spend your leisure time}, suis toujours ces étapes :
\begin{enumerate}
\item \textbf{Lire le sujet deux fois} et souligner les mots-clés : le thème (\eng{recreational activities}), le nombre de mots, le destinataire éventuel.
\item \textbf{Brainstormer} : noter en vrac 5 ou 6 idées (activités, lieux, fréquence, raisons, sentiments). Exemple : \eng{playing football, reading novels, watching films, visiting friends, helping at home}.
\item \textbf{Classer les idées en trois parties} :
\begin{itemize}
\item \emph{Introduction} : présenter le thème en une ou deux phrases générales (\eng{Like most teenagers in Yaoundé, I enjoy several recreational activities after school.}).
\item \emph{Développement} : deux ou trois paragraphes, une idée principale par paragraphe, illustrée d'un exemple concret.
\item \emph{Conclusion} : résumer et donner son opinion personnelle (\eng{In short, leisure activities help me relax and make new friends.}).
\end{itemize}
\item \textbf{Rédiger au brouillon} en reliant les idées avec des connecteurs (\eng{first, moreover, however, finally}).
\item \textbf{Compter les mots} et vérifier que l'on reste dans la fourchette demandée.
\item \textbf{Recopier au propre} en écrivant lisiblement, sans ratures.
\end{enumerate}
Réflexe d'or : une phrase = une idée ; un paragraphe = une activité ou un argument.
\end{methode}

\begin{exoresolu}[Rédiger une composition guidée]
\textbf{Énoncé.} \eng{Write a composition of about 120 words on the topic: “My favourite recreational activity”. Say what the activity is, when and where you do it, and why you enjoy it.}
\tcblower
\textbf{Étape 1 — Brainstorming.} Activité choisie : \eng{playing football}. Idées : \eng{with classmates, on the school field, every Saturday, keeps me fit, teaches teamwork, fun}.
\textbf{Étape 2 — Plan.} Introduction : phrase générale. Développement : description (quand, où, avec qui) puis raisons. Conclusion : opinion.
\textbf{Étape 3 — Rédaction (modèle).}

\eng{Of all the recreational activities I know, playing football is my favourite. Every Saturday afternoon, I meet my classmates on the field behind our school in Buea, and we play for about two hours.}

\eng{First, football keeps me fit and healthy because we run a lot during the match. Moreover, it teaches me teamwork: to win, every player must pass the ball and help the others. Finally, I enjoy it because it is simply fun and it helps me forget my school worries.}

\eng{In short, football is not only a game for me; it is also a school of life. That is why I never miss a Saturday match.}

(Environ 110 mots.) \textbf{Justification} : le texte respecte le plan en trois parties, utilise les connecteurs \eng{first, moreover, finally, in short}, et les temps sont corrects : présent simple pour les habitudes (\eng{I meet, we play, it teaches}).
\end{exoresolu}

\begin{methode}[Organiser les idées avec des connecteurs logiques]
Les connecteurs donnent du relief à une composition et rapportent des points. Retiens cette boîte à outils :
\begin{center}
\begin{tabularx}{\linewidth}{|X|X|X|}\hline
\rowcolor{popblueL} Fonction & Connecteurs anglais & Équivalent français \\ \hline
Ajout & \eng{moreover, besides, also, in addition} & de plus, aussi \\ \hline
Ordre & \eng{first, then, next, finally} & d'abord, ensuite, enfin \\ \hline
Opposition & \eng{however, but, whereas, although} & cependant, mais, bien que \\ \hline
Cause & \eng{because, since, as} & parce que, puisque \\ \hline
Conséquence & \eng{so, therefore, that is why} & donc, c'est pourquoi \\ \hline
Exemple & \eng{for example, for instance, such as} & par exemple \\ \hline
Conclusion & \eng{in short, to conclude, in my opinion} & bref, pour conclure \\ \hline
\end{tabularx}
\end{center}
Démarche :
\begin{enumerate}
\item Identifier la relation logique entre deux idées (ajout ? opposition ? conséquence ?).
\item Choisir le connecteur correspondant et le placer correctement : en début de phrase suivi d'une virgule (\eng{However, I prefer reading.}) ou entre les deux propositions (\eng{I like dancing, but I prefer singing.}).
\item Ne jamais utiliser deux connecteurs pour la même relation (\eng{although\ldots but} est interdit).
\item Varier : ne pas répéter \eng{and} ou \eng{because} dans chaque phrase.
\end{enumerate}
\end{methode}

\begin{exoresolu}[Relier des phrases avec le bon connecteur]
\textbf{Énoncé.} \eng{Join each pair of sentences using the connector in brackets.}\\
1. \eng{I enjoy swimming. I do not know how to dive. (however)}\\
2. \eng{We stayed at home. It was raining heavily. (because)}\\
3. \eng{Mary sings well. She dances very well too. (moreover)}
\tcblower
\textbf{1.} La relation est une opposition entre aimer nager et ne pas savoir plonger. \eng{However} se place en tête de la deuxième phrase, suivi d'une virgule : \eng{I enjoy swimming; however, I do not know how to dive.} (« J'aime nager ; cependant, je ne sais pas plonger. »)

\textbf{2.} La pluie est la cause : \eng{because} introduit la subordonnée de cause. \eng{We stayed at home because it was raining heavily.} (« Nous sommes restés à la maison parce qu'il pleuvait fort. ») Remarque : pas de virgule avant \eng{because} quand la subordonnée suit la principale.

\textbf{3.} Il s'agit d'un ajout : \eng{Mary sings well; moreover, she dances very well too.} (« Marie chante bien ; de plus, elle danse très bien aussi. ») On peut aussi écrire : \eng{Moreover, Mary dances very well}, mais il faut éviter de répéter l'idée. \textbf{Conclusion} : toujours vérifier la ponctuation : virgule après le connecteur initial, point-virgule possible entre deux phrases indépendantes.
\end{exoresolu}

\begin{methode}[Choisir les bons temps verbaux dans une composition]
Le choix du temps dépend de ce que tu racontes :
\begin{center}
\begin{tabular}{|l|l|l|}\hline
\rowcolor{popblueL} Temps & Emploi & Exemple \\ \hline
Présent simple & habitudes, goûts, vérités générales & \eng{I play tennis every Friday.} \\ \hline
Présent continu (\eng{be} + \emph{-ing}) & action en cours, projet futur arrêté & \eng{We are training for the match.} / \eng{I am meeting him tomorrow.} \\ \hline
Prétérit simple & récit d'un événement passé daté & \eng{Last weekend, we visited the zoo.} \\ \hline
\eng{be going to} / Futur simple & intention, prédiction & \eng{I am going to join the choir.} / \eng{It will be fun.} \\ \hline
\end{tabular}
\end{center}
Démarche :
\begin{enumerate}
\item Décider si la composition décrit des habitudes (présent simple) ou raconte un souvenir (prétérit).
\item Rester cohérent : ne pas passer du prétérit au présent sans raison (changement de perspective).
\item Employer les marqueurs de temps adaptés : \eng{every day, usually, often} avec le présent ; \eng{last week, yesterday, two days ago, in 2022} avec le prétérit ; \eng{tomorrow, next week, soon} avec le futur.
\item Vérifier la troisième personne du singulier au présent : \eng{he plays, she watches, it relaxes}.
\item Conjuguer correctement les irréguliers au prétérit : \eng{go $\rightarrow$ went, have $\rightarrow$ had, take $\rightarrow$ took, swim $\rightarrow$ swam, do $\rightarrow$ did}.
\end{enumerate}
\end{methode}

\begin{exoresolu}[Mettre les verbes au bon temps]
\textbf{Énoncé.} \eng{Complete the paragraph with the correct form of the verbs in brackets.}\\
\eng{Last Saturday, my friends and I (1. go) to the municipal stadium in Bamenda. We (2. watch) a football match and we (3. have) a wonderful time. Usually, I (4. play) football myself, but that day I (5. be) just a supporter. Next weekend, we (6. organise) a match between our two schools.}
\tcblower
\textbf{1. went} : prétérit de \eng{go}, marqueur \eng{last Saturday}. \textbf{2. watched} : prétérit régulier en \emph{-ed}. \textbf{3. had} : prétérit irrégulier de \eng{have}. \textbf{4. play} : le marqueur \eng{usually} impose le présent simple (habitude), première personne : pas de \emph{-s}. \textbf{5. was} : on revient au récit passé ; prétérit de \eng{be} avec \eng{I}. \textbf{6. are going to organise} (ou \eng{will organise} / \eng{are organising}) : \eng{next weekend} annonce un projet futur.

\textbf{Paragraphe corrigé :} \eng{Last Saturday, my friends and I went to the municipal stadium in Bamenda. We watched a football match and we had a wonderful time. Usually, I play football myself, but that day I was just a supporter. Next weekend, we are going to organise a match between our two schools.} (« Samedi dernier, mes amis et moi sommes allés au stade municipal de Bamenda. Nous avons regardé un match et nous avons passé un moment merveilleux. D'habitude, je joue moi-même au football, mais ce jour-là j'étais simple supporter. Le week-end prochain, nous allons organiser un match entre nos deux écoles. »)
\end{exoresolu}

\begin{methode}[Relire et corriger sa production (la grille COCO)]
Avant de rendre ta copie, relis-la en suivant la grille \textbf{COCO} :
\begin{enumerate}
\item \textbf{C}ontenu : ai-je répondu à toutes les questions du sujet ? Ai-je respecté le nombre de mots (\eng{word count}) ?
\item \textbf{O}rganisation : introduction, développement en paragraphes, conclusion ? Connecteurs variés ?
\item \textbf{C}ohérence grammaticale : temps uniformes, sujet-verbe accordés (\eng{he plays}, pas \eng{he play}), articles (\eng{a, the}), pluriels (\eng{activities}).
\item \textbf{O}rthographe et ponctuation : majuscule en début de phrase et pour \eng{I}, point final, mots difficiles vérifiés (\eng{beautiful, because, favourite}).
\end{enumerate}
Réflexe : garder cinq minutes à la fin de l'épreuve uniquement pour la relecture ; corriger proprement au stylo, sans surligneur.
\end{methode}

\begin{exoresolu}[Corriger une composition fautive]
\textbf{Énoncé.} \eng{The following paragraph contains six mistakes. Find them and rewrite the paragraph correctly.}\\
\eng{Last sunday, i assist to a concert with my best friend. The musicians was very good and we dance all the evening. I am very happy because music make me relaxed.}
\tcblower
\textbf{Repérage des fautes :}
\begin{enumerate}
\item \eng{sunday} $\rightarrow$ \eng{Sunday} : majuscule obligatoire aux jours de la semaine.
\item \eng{i} $\rightarrow$ \eng{I} : le pronom personnel \eng{I} prend toujours une majuscule.
\item \eng{assist to} $\rightarrow$ \eng{attended} : faux ami ! \eng{to assist} signifie « aider » ; « assister à » se dit \eng{to attend} (sans préposition).
\item \eng{the musicians was} $\rightarrow$ \eng{the musicians were} : accord avec le sujet pluriel.
\item \eng{we dance} $\rightarrow$ \eng{we danced} : le récit est au prétérit (\eng{last Sunday}).
\item \eng{music make} $\rightarrow$ \eng{music makes} : troisième personne du singulier au présent simple. (\eng{all the evening} $\rightarrow$ \eng{all evening} est aussi plus naturel.)
\end{enumerate}
\textbf{Paragraphe corrigé :} \eng{Last Sunday, I attended a concert with my best friend. The musicians were very good and we danced all evening. I was very happy because music makes me relaxed.} (« Dimanche dernier, j'ai assisté à un concert avec mon meilleur ami. Les musiciens étaient très bons et nous avons dansé toute la soirée. J'étais très heureux car la musique me détend. ») \textbf{Conclusion} : la méthode COCO permet de repérer méthodiquement chaque type de faute.
\end{exoresolu}

\begin{methode}[Utiliser le vocabulaire des activités récréatives avec les bons verbes]
En anglais, chaque activité réclame son verbe : \eng{play} (sports avec ballon, jeux, instruments), \eng{go} + verbe en \emph{-ing} (activités de déplacement/loisir), \eng{do} (activités individuelles, exercices).
\begin{center}
\begin{tabularx}{\linewidth}{|X|X|X|}\hline
\rowcolor{popblueL} Verbe & Activités & Exemple \\ \hline
\eng{play} & \eng{football, basketball, tennis, the guitar, the piano, cards, ludo} & \eng{We play football on Saturdays.} \\ \hline
\eng{go} + \emph{-ing} & \eng{swimming, dancing, shopping, fishing, hiking, jogging, cycling} & \eng{They go swimming at the pool.} \\ \hline
\eng{do} & \eng{gymnastics, karate, judo, exercise, yoga, puzzles, homework} & \eng{She does karate after school.} \\ \hline
\eng{watch / listen to / read} & \eng{films, TV, music, the radio, novels, comics} & \eng{I listen to music every evening.} \\ \hline
\end{tabularx}
\end{center}
Réflexes : pas d'article après \eng{play} + sport (\eng{I play football}, pas \eng{the football}), mais article défini avec un instrument (\eng{I play the guitar}). Pour exprimer le goût : \eng{I enjoy / like / love / prefer + \emph{-ing}} : \eng{I enjoy reading novels.} (« J'aime lire des romans. »)
\end{methode}

\begin{exoresolu}[Compléter avec le bon verbe d'activité]
\textbf{Énoncé.} \eng{Complete each sentence with play, go or do in the correct form.}\\
1. \eng{Every morning, my father \ldots\ exercise before going to work.}\\
2. \eng{At weekends, we often \ldots\ fishing on the River Wouri.}\\
3. \eng{Nadia \ldots\ the piano very well; she practises every day.}\\
4. \eng{My friends \ldots\ volleyball in the school yard during the break.}
\tcblower
\textbf{1. does} : \eng{exercise} est une activité individuelle $\rightarrow$ \eng{do} ; présent simple, 3\up{e} personne : \eng{does}. (« Chaque matin, mon père fait de l'exercice avant d'aller au travail. »)

\textbf{2. go} : \eng{fishing} est un verbe en \emph{-ing} $\rightarrow$ \eng{go fishing} ; sujet pluriel \eng{we} : \eng{go}. (« Le week-end, nous allons souvent pêcher sur le Wouri. »)

\textbf{3. plays} : instrument de musique $\rightarrow$ \eng{play the piano} ; 3\up{e} personne : \eng{plays}. (« Nadia joue très bien du piano ; elle s'entraîne chaque jour. »)

\textbf{4. play} : sport avec ballon $\rightarrow$ \eng{play volleyball} ; sujet pluriel : \eng{play}. (« Mes amis jouent au volley dans la cour pendant la récréation. »)

\textbf{Conclusion} : mémoriser les triplets \eng{play / go / do} évite les calques du français « faire du sport » traduit mot à mot.
\end{exoresolu}

\begin{attention}[Les 5 erreurs qui coûtent des points au Bac]
\begin{enumerate}
\item \textbf{Faux ami \eng{assist}.} « Assister à un match » se dit \eng{to attend a match} (sans \eng{to}). \eng{To assist} signifie « aider ». On écrit : \eng{I attended a football match.} et non \eng{I assisted to a match}.
\item \textbf{Calque « faire du sport ».} Ne jamais traduire par \eng{make sport}. On dit \eng{to do sport}, \eng{to play football}, \eng{to do exercise}. De même, « passer du bon temps » = \eng{to have a good time}, pas \eng{to pass a good time}.
\item \textbf{Oubli du \emph{-s} à la 3\up{e} personne.} \eng{She play} est fautif : \eng{She plays tennis every weekend.} Vérifier chaque verbe au présent simple.
\item \textbf{Mélange des temps.} Un récit commencé au prétérit doit y rester : \eng{We went to the beach and we swam}, pas \eng{we go to the beach and we swim}. Surveiller les marqueurs : \eng{last} $\rightarrow$ prétérit, \eng{every} $\rightarrow$ présent.
\item \textbf{Majuscules et orthographe.} \eng{I} et les jours/mois prennent toujours une majuscule (\eng{Saturday, January}). Attention aux mots pièges : \eng{because} (pas \eng{becauce}), \eng{favourite} (orthographe britannique), \eng{beautiful}, \eng{activities} (pluriel en \emph{-ies}).
\end{enumerate}
\end{attention}

\newpage

\coursec{Exercices}

\courssub{Je vérifie mes connaissances}

\begin{exercice}[QCM : la structure d'une composition]
Entoure la bonne réponse.
\begin{enumerate}
\item A composition about recreational activities should normally contain :
\begin{enumerate}
\item only one long paragraph
\item an introduction, a body and a conclusion
\item a title, a signature and a date
\end{enumerate}
\item The sentence that presents the topic at the beginning of a composition is called :
\begin{enumerate}
\item a topic sentence
\item a linking word
\item a conclusion
\end{enumerate}
\item To add an idea, the best connector is :
\begin{enumerate}
\item however
\item moreover
\item in conclusion
\end{enumerate}
\item To introduce the final paragraph, we usually write :
\begin{enumerate}
\item first of all
\item for example
\item to sum up
\end{enumerate}
\end{enumerate}
\end{exercice}

\begin{exercice}[Vrai ou faux ? Justifie]
Réponds par \eng{True} ou \eng{False} et justifie ta réponse en une phrase.
\begin{enumerate}
\item In a composition, you can write your ideas in any order.
\item The connector \eng{however} is used to express a contrast.
\item When you describe a hobby you practise regularly, you generally use the present simple.
\item It is a good idea to repeat the same word, like \eng{like}, in every sentence.
\item You should always reread your composition to correct mistakes before handing it in.
\end{enumerate}
\end{exercice}

\begin{exercice}[Vocabulaire des loisirs]
Complète chaque phrase avec le mot qui convient, choisi dans la liste : \eng{leisure} — \eng{team} — \eng{indoor} — \eng{outdoor} — \eng{relaxing} — \eng{tournament}.
\begin{enumerate}
\item Football and athletics are \ldots\ldots\ldots activities because they are practised outside.
\item Reading and playing chess are \ldots\ldots\ldots activities: you can do them at home.
\item Our school organised a football \ldots\ldots\ldots last term; eight schools took part in it.
\item After a long week at school, listening to music is very \ldots\ldots\ldots .
\item In my free time, or \ldots\ldots\ldots time, I enjoy swimming with my friends.
\item A volleyball \ldots\ldots\ldots is made up of six players.
\end{enumerate}
\end{exercice}

\begin{exercice}[Conjugaison : les bons temps]
Mets le verbe entre parenthèses au temps qui convient (présent simple, présent en \eng{-ing} ou prétérit).
\begin{enumerate}
\item Every Saturday, my friends and I \ldots\ldots\ldots (play) football in the neighbourhood.
\item Look! The children \ldots\ldots\ldots (dance) traditional dances at the school party.
\item Last holidays, we \ldots\ldots\ldots (visit) the Limbe Wildlife Centre.
\item She \ldots\ldots\ldots (enjoy) singing in the church choir.
\item While I \ldots\ldots\ldots (watch) a film yesterday evening, the lights went off.
\item My brother usually \ldots\ldots\ldots (spend) his free time drawing cartoons.
\end{enumerate}
\end{exercice}

\courssub{Je m'entraîne}

\begin{exercice}[Traduction]
Traduis les phrases suivantes en anglais.
\begin{enumerate}
\item Pendant mon temps libre, j'aime jouer au football avec mes camarades.
\item D'abord, les activités récréatives nous aident à nous détendre après les cours.
\item De plus, le sport est bon pour la santé.
\item Cependant, certains élèves préfèrent regarder la télévision toute la journée.
\item En conclusion, je pense que les loisirs sont très importants dans la vie des jeunes.
\end{enumerate}
\end{exercice}

\begin{exercice}[Composition à trous : les connecteurs]
Complète ce texte avec les connecteurs de la liste : \eng{first of all} — \eng{moreover} — \eng{however} — \eng{for example} — \eng{in conclusion}.
\begin{textebox}[My favourite recreational activity]
\ldots\ldots\ldots , I must say that my favourite recreational activity is playing basketball. I started playing it when I was in Form Three. \ldots\ldots\ldots , basketball keeps me fit and teaches me how to work in a team. \ldots\ldots\ldots , I have learnt to pass the ball instead of playing alone. \ldots\ldots\ldots , some of my classmates find basketball too tiring; they prefer quieter games like ludo. \ldots\ldots\ldots , I would say that basketball is more than a game for me: it is a school of life.
\end{textebox}
\end{exercice}

\begin{exercice}[Corrige les erreurs]
Chaque phrase contient une erreur typique des francophones. Réécris-la correctement.
\begin{enumerate}
\item I am playing football every weekend.
\item She enjoy reading novels very much.
\item They are fond to swim in the river.
\item Last year, I have joined the school choir.
\item My sister she likes dancing makossa.
\item We played football during two hours yesterday.
\item He is very good in mathematics and in tennis.
\item I like very much listening to music.
\end{enumerate}
\end{exercice}

\begin{exercice}[Appariements : varier les expressions]
Réécris le paragraphe suivant en remplaçant les quatre occurrences de \eng{like} par des expressions variées : \eng{enjoy} — \eng{be fond of} — \eng{be keen on} — \eng{prefer}. Attention à la construction de chaque expression.
\begin{textebox}[Free time in my family]
In my family, everybody has a different taste. I like playing football with my friends on Saturdays. My elder brother likes video games and can spend hours in front of his screen. My mother likes gardening: she says it relaxes her after work. As for my little sister, she likes drawing more than any other activity.
\end{textebox}
\end{exercice}

\courssub{Je me prépare au Bac}

\begin{exercice}[Compréhension et langue]
Lis le texte suivant, puis réponds aux questions en anglais, avec tes propres mots.
\begin{textebox}[Recreation in a Cameroonian school]
At Government High School Mbankomo, Fridays are special. After classes end at two o'clock, the school field comes alive. Some students rush to the football pitch, where the senior team is training for the inter-school competition. Others gather under the mango trees to play draughts or sing with the school choir. A group of girls has recently started a dance club, and their performances now attract the whole school.

The principal, Mr Etoundi, believes that recreational activities are not a waste of time. “A student who only reads books becomes dull,” he says. “Sports and games teach discipline, cooperation and leadership.” Indeed, teachers have noticed that students who take part in clubs are often more confident in class. However, a few parents complain that their children come home late on Fridays. The school is now discussing the problem with the Parents' Association so that everyone can benefit from recreation without neglecting studies.
\end{textebox}
\begin{enumerate}
\item Why are Fridays special at Government High School Mbankomo? (2 marks)
\item Mention three recreational activities practised in the school. (3 marks)
\item According to the principal, what do sports and games teach? (2 marks)
\item What problem do some parents complain about? (2 marks)
\item Find in the text a word or expression that means: (a) \eng{full of life and activity}; (b) \eng{pay no attention to}. (2 marks)
\item Grammar: rewrite in the passive voice — \eng{A group of girls has recently started a dance club.} (2 marks)
\item Grammar: put the verb in the correct tense — \eng{Teachers (notice) that students who take part in clubs are more confident.} (2 marks)
\end{enumerate}
\end{exercice}

\begin{exercice}[Traduction]
Traduis le passage suivant en anglais (environ 60 mots). Sois attentif aux temps et aux prépositions.
\begin{quote}
Dans mon quartier à Yaoundé, les jeunes adorent les activités récréatives. Le week-end, certains jouent au football sur le terrain du quartier, tandis que d'autres dansent ou écoutent de la musique. Moi, je préfère jouer aux dames avec mon grand-père. Ces moments nous permettent de nous détendre et d'oublier les soucis de l'école.
\end{quote}
\end{exercice}

\begin{exercice}[Sujet de composition type Bac]
\begin{textebox}[Composition topic]
Write a composition of about 150–200 words on the following topic:

“Describe your favourite recreational activity and explain why you enjoy it.”

Your composition should have an introduction, two body paragraphs and a conclusion. Use appropriate linking words (first of all, moreover, however, in conclusion\ldots) and varied vocabulary.
\end{textebox}
Avant de rédiger, construis ton plan :
\begin{enumerate}
\item Introduction : présente l'activité choisie en une ou deux phrases.
\item Paragraphe 1 : décris l'activité (où, quand, avec qui, comment).
\item Paragraphe 2 : explique pourquoi tu l'aimes (au moins deux raisons avec des exemples).
\item Conclusion : résume ton idée principale et ouvre sur un conseil aux autres élèves.
\end{enumerate}
Barème indicatif : respect du sujet et du plan (5 pts), richesse du vocabulaire (5 pts), correction grammaticale (5 pts), connecteurs et organisation (5 pts).
\end{exercice}

\newpage

\coursec{Corrigés des exercices}

\begin{corrige}[Exercice 1 — QCM : la structure d'une composition]
\begin{enumerate}
\item \textbf{(b) an introduction, a body and a conclusion.} Une composition se compose toujours de trois parties : une introduction qui présente le sujet, un développement (\eng{body}) en un ou plusieurs paragraphes, et une conclusion.
\item \textbf{(a) a topic sentence.} La phrase qui présente le sujet ou l'idée principale s'appelle une \eng{topic sentence}. Un \eng{linking word} est un connecteur ; la \eng{conclusion} clôt le texte.
\item \textbf{(b) moreover.} \eng{Moreover} (« de plus ») ajoute une idée. \eng{However} exprime un contraste ; \eng{in conclusion} introduit la conclusion.
\item \textbf{(c) to sum up.} \eng{To sum up} (« pour résumer ») introduit le paragraphe final. \eng{First of all} sert à commencer le développement ; \eng{for example} introduit un exemple.
\end{enumerate}
\end{corrige}

\begin{corrige}[Exercice 2 — Vrai ou faux ? Justifie]
\begin{enumerate}
\item \textbf{False.} \eng{Ideas must be organised logically with linking words: introduction, body and conclusion.} (Les idées doivent suivre un plan logique, sinon le texte est incompréhensible.)
\item \textbf{True.} \eng{However is a connector of contrast, like but or whereas.} Exemple : \eng{Some students prefer indoor games; however, I enjoy outdoor sports.}
\item \textbf{True.} \eng{The present simple is used for habits and regular activities.} Exemple : \eng{I play basketball every Friday.}
\item \textbf{False.} \eng{Repeating the same word makes the composition poor; you should vary your vocabulary (enjoy, be fond of, be keen on, prefer).}
\item \textbf{True.} \eng{Rereading helps you correct grammar, spelling and punctuation mistakes before handing in your work.}
\end{enumerate}
\end{corrige}

\begin{corrige}[Exercice 3 — Vocabulaire des loisirs]
\begin{enumerate}
\item \eng{Football and athletics are \textbf{outdoor} activities because they are practised outside.} (\eng{outdoor} = en plein air, à l'extérieur.)
\item \eng{Reading and playing chess are \textbf{indoor} activities: you can do them at home.} (\eng{indoor} = en salle, à l'intérieur.)
\item \eng{Our school organised a football \textbf{tournament} last term; eight schools took part in it.} (\eng{tournament} = tournoi.)
\item \eng{After a long week at school, listening to music is very \textbf{relaxing}.} (\eng{relaxing} = relaxant ; attention, \eng{relaxed} décrit la personne détendue.)
\item \eng{In my free time, or \textbf{leisure} time, I enjoy swimming with my friends.} (\eng{leisure} = loisirs ; prononciation « lè-jeur ».)
\item \eng{A volleyball \textbf{team} is made up of six players.} (\eng{team} = équipe.)
\end{enumerate}
\end{corrige}

\begin{corrige}[Exercice 4 — Conjugaison : les bons temps]
\begin{enumerate}
\item \eng{Every Saturday, my friends and I \textbf{play} football in the neighbourhood.} Présent simple : habitude marquée par \eng{every Saturday}.
\item \eng{Look! The children \textbf{are dancing} traditional dances at the school party.} Présent en \eng{-ing} : action en cours, signalée par \eng{Look!}
\item \eng{Last holidays, we \textbf{visited} the Limbe Wildlife Centre.} Prétérit simple : action passée et terminée (\eng{last holidays}).
\item \eng{She \textbf{enjoys} singing in the church choir.} Présent simple, 3\textsuperscript{e} personne : \eng{enjoy} + \textbf{s}. Rappel : \eng{enjoy} est suivi d'un verbe en \eng{-ing}.
\item \eng{While I \textbf{was watching} a film yesterday evening, the lights went off.} Prétérit en \eng{-ing} après \eng{while} : action longue interrompue par une action courte (\eng{went off}).
\item \eng{My brother usually \textbf{spends} his free time drawing cartoons.} Présent simple avec \eng{usually} ; 3\textsuperscript{e} personne : \eng{spend} + \textbf{s}.
\end{enumerate}
\end{corrige}

\begin{corrige}[Exercice 5 — Traduction]
\begin{enumerate}
\item \eng{In my free time, I like playing football with my classmates.} (Après \eng{like}, verbe en \eng{-ing} ; \eng{classmates} = camarades de classe.)
\item \eng{First of all, recreational activities help us (to) relax after classes.} (\eng{help} + quelqu'un + base verbale, avec ou sans \eng{to}.)
\item \eng{Moreover, sport is good for our health.} (Attention à la préposition : \eng{good \textbf{for}} ; \eng{health} se prononce « hèlth ».)
\item \eng{However, some students prefer watching television all day long.} (\eng{prefer} + \eng{-ing} ; \eng{all day long} = toute la journée.)
\item \eng{In conclusion, I think (that) leisure activities are very important in young people's lives.} (Ne pas oublier la majuscule à \eng{In conclusion} en début de phrase.)
\end{enumerate}
\end{corrige}

\begin{corrige}[Exercice 6 — Composition à trous : les connecteurs]
\begin{enumerate}
\item \eng{\textbf{First of all}, I must say that my favourite recreational activity is playing basketball.} (Ouverture du texte : on présente le sujet.)
\item \eng{\textbf{Moreover}, basketball keeps me fit and teaches me how to work in a team.} (Ajout d'un premier argument.)
\item \eng{\textbf{For example}, I have learnt to pass the ball instead of playing alone.} (Illustration de l'argument précédent par un exemple.)
\item \eng{\textbf{However}, some of my classmates find basketball too tiring; they prefer quieter games like ludo.} (Contraste avec un avis différent.)
\item \eng{\textbf{In conclusion}, I would say that basketball is more than a game for me: it is a school of life.} (Formule de clôture.)
\end{enumerate}
\textbf{Point de méthode :} chaque connecteur est suivi d'une virgule en tête de phrase ; la logique du texte (présentation $\rightarrow$ argument $\rightarrow$ exemple $\rightarrow$ contraste $\rightarrow$ conclusion) doit guider le choix.
\end{corrige}

\begin{corrige}[Exercice 7 — Corrige les erreurs]
\begin{enumerate}
\item \eng{I \textbf{play} football every weekend.} Habitude (\eng{every weekend}) $\rightarrow$ présent simple, pas de présent en \eng{-ing}.
\item \eng{She \textbf{enjoys} reading novels very much.} À la 3\textsuperscript{e} personne du singulier, le verbe prend un \textbf{s}.
\item \eng{They are fond \textbf{of swimming} in the river.} Construction : \eng{be fond \textbf{of} + -ing} (jamais \eng{fond to}).
\item \eng{Last year, I \textbf{joined} the school choir.} Avec un repère passé précis (\eng{last year}), on utilise le prétérit simple, jamais le present perfect.
\item \eng{My sister likes dancing makossa.} Pas de double sujet : on supprime \eng{she} (calque du français « ma sœur, elle… »).
\item \eng{We played football \textbf{for} two hours yesterday.} Durée $\rightarrow$ \eng{for} ; \eng{during} s'emploie devant un nom (\eng{during the match}).
\item \eng{He is very good \textbf{at} mathematics and \textbf{at} tennis.} Construction : \eng{be good \textbf{at}} (jamais \eng{good in}).
\item \eng{I like listening to music very much.} En anglais, \eng{very much} se place en fin de phrase, jamais entre le verbe et son complément.
\end{enumerate}
\end{corrige}

\begin{corrige}[Exercice 8 — Appariements : varier les expressions]
\textbf{Paragraphe réécrit :}

\eng{In my family, everybody has a different taste. I \textbf{enjoy playing} football with my friends on Saturdays. My elder brother \textbf{is fond of} video games and can spend hours in front of his screen. My mother \textbf{is keen on gardening}: she says it relaxes her after work. As for my little sister, she \textbf{prefers} drawing \textbf{to} any other activity.}

\textbf{Constructions à respecter :}
\begin{center}
\begin{tabularx}{\linewidth}{|X|X|}\hline
\rowcolor{popblueL} Expression & Construction \\ \hline
\eng{enjoy} & \eng{enjoy + -ing} (\eng{I enjoy playing}) \\ \hline
\eng{be fond of} & \eng{be fond of + nom ou -ing} \\ \hline
\eng{be keen on} & \eng{be keen on + nom ou -ing} \\ \hline
\eng{prefer} & \eng{prefer + nom/-ing + to + nom/-ing} \\ \hline
\end{tabularx}
\end{center}
\textbf{Attention :} après \eng{prefer}, la comparaison se fait avec \eng{to} : \eng{She prefers drawing to dancing}, et non \eng{than}.
\end{corrige}

\begin{corrige}[Exercice 9 — Compréhension et langue]
\begin{enumerate}
\item \eng{Fridays are special because classes end at two o'clock and students take part in recreational activities on the school field.} (2 marks)
\item \eng{Three recreational activities are: playing football, playing draughts and singing with the school choir.} (On accepte aussi \eng{dancing in the dance club}.) (3 marks)
\item \eng{According to the principal, sports and games teach discipline, cooperation and leadership.} (2 marks)
\item \eng{Some parents complain that their children come home late on Fridays.} (2 marks)
\item (a) \eng{comes alive} (= full of life and activity) ; (b) \eng{neglecting} (= pay no attention to). (2 marks)
\item Voix passive : \eng{A dance club has recently been started by a group of girls.} Auxiliaire \eng{be} au présent perfect (\eng{has been}) + participe passé \eng{started}. (2 marks)
\item \eng{Teachers \textbf{have noticed} that students who take part in clubs are more confident.} Present perfect : constat dont le résultat est visible aujourd'hui (le texte dit \eng{have noticed}). (2 marks)
\end{enumerate}
\textbf{Barème indicatif :} 15 points au total ; 1 point par réponse correcte en compréhension, 1 point pour la correction grammaticale de la formulation.
\end{corrige}

\begin{corrige}[Exercice 10 — Traduction]
\textbf{Proposition de traduction :}

\eng{In my neighbourhood in Yaoundé, young people love recreational activities. At the weekend, some play football on the neighbourhood field, while others dance or listen to music. As for me, I prefer playing draughts with my grandfather. These moments allow us to relax and to forget the worries of school.}

\textbf{Points de vigilance :}
\begin{itemize}
\item \eng{In my neighbourhood} : préposition \eng{in} (pas \eng{at}).
\item \eng{At the weekend} : préposition \eng{at} devant \eng{weekend}.
\item \eng{some\ldots while others\ldots} : bonne traduction de « certains\ldots tandis que d'autres\ldots ».
\item \eng{playing draughts} : \eng{prefer} + \eng{-ing} ; « les dames » (jeu) = \eng{draughts}, invariable.
\item \eng{allow us to relax} : \eng{allow} + quelqu'un + \eng{to} + base verbale (« permettre à quelqu'un de »).
\item Tous les verbes au présent simple : il s'agit d'habitudes.
\end{itemize}
\end{corrige}

\begin{corrige}[Exercice 11 — Sujet de composition type Bac]
\textbf{Plan attendu :}
\begin{enumerate}
\item Introduction : présentation de l'activité favorite en une ou deux phrases (\eng{topic sentence}).
\item Paragraphe 1 : description (où, quand, avec qui, comment).
\item Paragraphe 2 : au moins deux raisons d'aimer cette activité, avec des exemples.
\item Conclusion : résumé + conseil aux autres élèves.
\end{enumerate}

\textbf{Proposition de rédaction modèle (environ 180 mots) :}

\eng{Of all recreational activities, the one I enjoy most is playing football. I have practised it since I was in primary school, and it has become an important part of my life.}

\eng{Every Saturday morning, I meet my friends on the field behind our school in Buea. We divide ourselves into two teams and play for about two hours. Sometimes, our coach organises small tournaments with neighbouring schools, which makes the game even more exciting.}

\eng{I love football for several reasons. First of all, it keeps me fit and healthy after a long week of studies. Moreover, it teaches me cooperation and discipline: for example, I have learnt to pass the ball instead of playing alone. However, football is not only about winning; it also helps me make new friends and forget my worries.}

\eng{In conclusion, football is much more than a game to me: it is a source of health, friendship and joy. I advise every student to take part in a recreational activity, because a healthy mind lives in a healthy body.}

\textbf{Barème indicatif (20 pts) :} respect du sujet et du plan en quatre parties (5 pts) ; richesse et variété du vocabulaire des loisirs (5 pts) ; correction grammaticale — temps, 3\textsuperscript{e} personne, prépositions (5 pts) ; connecteurs et organisation logique (5 pts).

\textbf{Points de vigilance :}
\begin{itemize}
\item Bien séparer les paragraphes ; une idée principale par paragraphe.
\item Varier les expressions de goût : \eng{enjoy}, \eng{be fond of}, \eng{love} + \eng{-ing}.
\item Employer le présent simple pour les habitudes, le présent perfect pour l'expérience (\eng{I have practised\ldots}, \eng{I have learnt\ldots}).
\item Relire pour vérifier le \textbf{s} de la 3\textsuperscript{e} personne, les prépositions (\eng{good at}, \eng{take part in}, \eng{listen to}) et l'orthographe.
\end{itemize}
\end{corrige}

\newpage

\coursec{Fiche bilan}

\begin{popfigure}
\begin{tikzpicture}[mindmap, grow cyclic, every node/.style={concept, font=\small}, root concept/.append style={concept color=popblue, font=\bfseries\small}, level 1/.append style={level distance=3.2cm, sibling angle=60}, text=white]
\node[root concept]{Writing: recreational activities}
  child[concept color=poporange]{ node[align=center]{Structure\\ intro -- body -- conclusion} }
  child[concept color=popgreen]{ node[align=center]{Connectors\\ first, then, moreover, finally} }
  child[concept color=poppurple]{ node[align=center]{Tenses\\ past simple, present perfect} }
  child[concept color=poppink]{ node[align=center]{Vocabulary\\ games, sports, music, outings} }
  child[concept color=popgold]{ node[align=center]{Planning\\ brainstorm, outline, draft} }
  child[concept color=popblue!70]{ node[align=center]{Checking\\ grammar, spelling, linking} };
\end{tikzpicture}
\legende{Carte mentale de la leçon : rédiger une composition sur les activités récréatives.}
\end{popfigure}

\begin{bilan}
\textbf{1. Lexique clé (English -- Français)}
\begin{itemize}
  \item \eng{recreational activities} : activités récréatives ; \eng{leisure / free time} : loisirs / temps libre
  \item \eng{to take part in} : participer à ; \eng{to join a club} : adhérer à un club
  \item \eng{to relax / to unwind} : se détendre ; \eng{to have fun / to enjoy oneself} : s'amuser
  \item \eng{a football match} : un match de football ; \eng{a choir} : une chorale ; \eng{a board game} : un jeu de société
  \item \eng{an outing / a trip} : une sortie / une excursion ; \eng{team spirit} : l'esprit d'équipe
  \item \eng{to keep fit} : se maintenir en forme ; \eng{to make new friends} : se faire de nouveaux amis
\end{itemize}

\textbf{2. Règles à connaître par cœur}
\begin{center}
\begin{tabularx}{\linewidth}{|X|X|}
\hline
\rowcolor{popblueL} Règle & Exemple \\
\hline
Récit d'une activité passée : \cle{past simple} (V-ed ou 2\textsuperscript{e} colonne) & \eng{Last Saturday, we played football in Buea.} \\
\hline
Expérience sans date précise : \cle{present perfect} (have/has + participe passé) & \eng{I have joined the school choir.} \\
\hline
Habitudes et goûts : \cle{present simple} (+ -s à la 3\textsuperscript{e} personne) & \eng{She enjoys singing every weekend.} \\
\hline
\eng{enjoy / practise / suggest} + \cle{verbe en -ing} (jamais + to) & \eng{I enjoy playing draughts.} \\
\hline
\eng{to be good at} + -ing ; \eng{to be interested in} + -ing & \eng{He is good at dancing.} \\
\hline
\end{tabularx}
\end{center}

\textbf{3. Connecteurs logiques indispensables}
\begin{itemize}
  \item Ordre : \eng{first, then, next, after that, finally}
  \item Addition : \eng{moreover, besides, also, in addition}
  \item Contraste : \eng{however, but, although, whereas}
  \item Cause et conséquence : \eng{because, so, therefore, as a result}
  \item Conclusion : \eng{to conclude, in short, all in all}
\end{itemize}

\textbf{4. Méthode express en 5 étapes}
\begin{enumerate}
  \item \cle{Comprendre} le sujet et noter les mots-clés.
  \item \cle{Brainstormer} : lister idées et vocabulaire en anglais.
  \item \cle{Planifier} : introduction (sujet + opinion), 2--3 paragraphes de développement (une idée + un exemple chacun), conclusion (bilan + ouverture).
  \item \cle{Rédiger} avec connecteurs et temps corrects.
  \item \cle{Relire} : accords, temps, orthographe, ponctuation, majuscules.
\end{enumerate}
\end{bilan}

\begin{aretenir}[Checklist avant le Bac]
Avant de rendre ma composition, je vérifie :
\begin{itemize}
  \item $\square$ J'ai une introduction, un développement (2--3 paragraphes) et une conclusion.
  \item $\square$ Chaque paragraphe contient une idée principale et un exemple.
  \item $\square$ J'ai utilisé au moins quatre connecteurs logiques variés.
  \item $\square$ Le past simple est correct pour les actions passées datées.
  \item $\square$ Les verbes après \eng{enjoy}, \eng{practise}, \eng{be good at} sont en -ing.
  \item $\square$ Le présent prend bien un -s à la troisième personne du singulier.
  \item $\square$ Le pronom \eng{I} est toujours en majuscule.
  \item $\square$ J'ai évité les faux amis : \eng{actually} (en fait), \eng{to assist} n'est pas « assister à » (\eng{to attend}).
  \item $\square$ J'ai utilisé le vocabulaire des loisirs vu en classe.
  \item $\square$ J'ai vérifié l'orthographe et la ponctuation.
  \item $\square$ Mon écriture est lisible et le nombre de mots est respecté.
\end{itemize}
\end{aretenir}

\findecours

\end{document}
